\documentclass[10pt]{article}
\usepackage[a4paper,margin=18mm]{geometry}
\usepackage[no-math]{fontspec}\setmainfont{Latin Modern Roman}
\usepackage{amsmath,amssymb,amsthm,mathrsfs,enumitem,tikz,graphicx,multirow,float}
\usepackage[bookmarks=false]{hyperref}\hypersetup{hidelinks}
\newtheorem{theorem}{Theorem}[section]
\newtheorem{proposition}[theorem]{Proposition}
\newtheorem{lemma}[theorem]{Lemma}
\newtheorem{fact}[theorem]{Fact}
\newtheorem{definition}[theorem]{Definition}
\newtheorem{notation}[theorem]{Notation}
\theoremstyle{definition}
\newtheorem{remark}[theorem]{Remark}
\newtheorem{example}[theorem]{Example}
\newtheorem*{propositionC}{Proposition C}
\numberwithin{equation}{section}
\emergencystretch=3em
\makeatletter
\renewcommand\subsection{\@startsection{subsection}{2}{\z@}{-3.25ex\@plus -1ex\@minus -.2ex}{1.5ex\@plus .2ex}{\normalfont\normalsize\bfseries}}
\makeatother
\newcommand{\jepPubCoroot}{{\scriptscriptstyle\vee}}
\newcommand{\jepPubBullet}{{\scriptscriptstyle\bullet}}
\newcommand{\jepPubRestrict}{{\scriptstyle|}}
\newcommand{\jepPubOplus}{\mathop{\textstyle\bigoplus}}
\renewcommand{\qedsymbol}{$\square$}
\DeclareMathAlphabet{\jepEuler}{U}{eus}{m}{n}
\DeclareRobustCommand{\jepPubScript}[1]{\begingroup\def\jepArg{#1}\def\jepTestC{C}\def\jepTestE{E}\def\jepTestF{F}\def\jepTestG{G}\def\jepTestL{L}\def\jepTestM{M}\def\jepTestO{O}\def\jepTestR{R}\def\jepTestS{S}\def\jepTestT{T}\def\jepTestU{U}\def\jepTestV{V}\def\jepTestW{W}\def\jepTestX{X}\ifx\jepArg\jepTestC\mathscr{C}\else\ifx\jepArg\jepTestE\mathscr{E}\else\ifx\jepArg\jepTestF\mathscr{F}\else\ifx\jepArg\jepTestG\mathscr{G}\else\ifx\jepArg\jepTestL\mathscr{L}\else\ifx\jepArg\jepTestM\mathscr{M}\else\ifx\jepArg\jepTestO\mathscr{O}\else\ifx\jepArg\jepTestR\mathscr{R}\else\ifx\jepArg\jepTestS\mathscr{S}\else\ifx\jepArg\jepTestT\mathscr{T}\else\ifx\jepArg\jepTestU\mathscr{U}\else\ifx\jepArg\jepTestV\mathscr{V}\else\ifx\jepArg\jepTestW\mathscr{W}\else\ifx\jepArg\jepTestX\mathscr{X}\else\PackageError{scoped-published-font}{Unsupported literal in jepPubScript}{Only inventoried published arguments are allowed.}\fi\fi\fi\fi\fi\fi\fi\fi\fi\fi\fi\fi\fi\fi\endgroup}
\DeclareRobustCommand{\jepPubBbb}[1]{\begingroup\def\jepArg{#1}\def\jepTestE{E}\def\jepTestF{F}\def\jepTestS{S}\def\jepTestU{U}\def\jepTestV{V}\def\jepTestW{W}\ifx\jepArg\jepTestE\mathbb{E}\else\ifx\jepArg\jepTestF\mathbb{F}\else\ifx\jepArg\jepTestS\mathbb{S}\else\ifx\jepArg\jepTestU\mathbb{U}\else\ifx\jepArg\jepTestV\mathbb{V}\else\ifx\jepArg\jepTestW\mathbb{W}\else\PackageError{scoped-published-font}{Unsupported literal in jepPubBbb}{Only inventoried published arguments are allowed.}\fi\fi\fi\fi\fi\fi\endgroup}
\DeclareRobustCommand{\jepPubBmi}[1]{\begingroup\def\jepArg{#1}\def\jepTestG{G}\def\jepTesty{y}\ifx\jepArg\jepTestG\boldsymbol{G}\else\ifx\jepArg\jepTesty\boldsymbol{y}\else\PackageError{scoped-published-font}{Unsupported literal in jepPubBmi}{Only inventoried published arguments are allowed.}\fi\fi\endgroup}
\DeclareRobustCommand{\jepPubLeaf}[1]{\begingroup\def\jepArg{#1}\def\jepTestL{L}\ifx\jepArg\jepTestL\jepEuler{L}\else\PackageError{scoped-published-font}{Unsupported literal in jepPubLeaf}{Only inventoried published arguments are allowed.}\fi\endgroup}
\newcommand{\jepGraph}[5]{%
\scalebox{0.7}{\begin{tikzpicture}[x=14.25pt,y=14.25pt,baseline=-3pt,line width=.24pt]
% #1 node records index/x/y/coefficient/above label/filled flag; #2 ordinary edges; #3 dotted edges; #4 double edges with direction; #5 minimum extent
\foreach \i/\x/\y/\coef/\above/\filled in {#1}{\coordinate (g\i) at (\x,\y);}
\foreach \a/\b in {#2}{\draw (g\a)--(g\b);}
\foreach \a/\b in {#3}{\path (g\a)--coordinate[midway] (mid\a) (g\b);\draw (g\a)--([xshift=-4.5pt]mid\a);\draw ([xshift=4.5pt]mid\a)--(g\b);\foreach \shift in {-2.5pt,0pt,2.5pt}{\fill ([xshift=\shift]mid\a) circle[radius=.25pt];}}
\foreach \a/\b/\dir in {#4}{\draw ([yshift=.55pt]g\a)--([yshift=.55pt]g\b);\draw ([yshift=-.55pt]g\a)--([yshift=-.55pt]g\b);\path (g\a)--coordinate[midway] (dm\a) (g\b);\ifnum\dir=1\draw ([xshift=-1.4pt,yshift=2.6pt]dm\a)--([xshift=1.4pt]dm\a)--([xshift=-1.4pt,yshift=-2.6pt]dm\a);\else\draw ([xshift=1.4pt,yshift=2.6pt]dm\a)--([xshift=-1.4pt]dm\a)--([xshift=1.4pt,yshift=-2.6pt]dm\a);\fi}
\foreach \i/\x/\y/\coef/\above/\filled in {#1}{\ifnum\filled=1\node[circle,draw,fill=black,minimum size=3.8pt,inner sep=0pt] at (g\i) {};\else\node[circle,draw,fill=white,minimum size=3.8pt,inner sep=0pt] at (g\i) {};\fi\ifx\coef\empty\else\pgfmathtruncatemacro{\branchlabel}{abs(\y)>0 ? 1:0}\ifnum\branchlabel=1\node[inner sep=0pt,anchor=west,xshift=3.5pt] at (g\i) {$\coef$};\else\node[inner sep=0pt,anchor=north,yshift=-3.5pt] at (g\i) {$\coef$};\fi\fi\ifx\above\empty\else\pgfmathtruncatemacro{\branchlabel}{abs(\y)>0 ? 1:0}\ifnum\branchlabel=1\node[inner sep=0pt,anchor=west,xshift=3.5pt] at (g\i) {$\above$};\else\node[inner sep=0pt,anchor=south,yshift=3.5pt] at (g\i) {$\above$};\fi\fi}
\end{tikzpicture}}}
\begin{document}
\begin{center}{\Large Stable item 1905: checked published transcription}\end{center}
\paragraph{Source and attribution.}
Pierre-Emmanuel Chaput and Julien Maubon, \emph{Maximal representations of cocompact complex hyperbolic lattices, a uniform approach}, Journal de l’École polytechnique -- Mathématiques 6 (2019), 231--281, DOI 10.5802/jep.93. \href{https://jep.centre-mersenne.org/articles/10.5802/jep.93/}{Publisher article}; \href{https://creativecommons.org/licenses/by/4.0/}{CC BY 4.0}.
\paragraph{Selected independent problem.}
Original Proposition C restricted to the simple tube-type real algebraic Hermitian target with simply connected complexification. The lattice is uniform in $\operatorname{SU}(1,n)$, with $n\geqslant2$, and the Toledo invariant uses the source curvature and volume normalizations. Complete original Sections 2--4, the Section 5 opening, complete Section 5.1, and the full 52-entry bibliography are retained with the introductory definitions and normalizations. The complete 52-page original publisher PDF is retained separately. Sections 5.2--5.3 and the extension to nonsimply connected targets are outside the selected unit; their original references remain printed. All supporting results are uncounted context.
\paragraph{Difficulty (editorial): H3.}
Cominuscule root calculations construct a graded nilpotent Higgs submodule, its parabolic stabilizer and arithmetic-progression character slopes. Its saturated leafwise Higgs sheaf gives the rank-sensitive foliated Milnor-Wood inequality. Full-rank Higgs iteration and semistability of symmetric cotangent powers then yield the stronger tube-type bound.
\paragraph{Representation disclosure.}
This is an editable publisher-authority transcription. The nonexclusive arXiv v1 archive was read only as an unexecuted private aid; no private archive, preamble, class, comment, or unpublished difference is exported or loaded. Every published Dynkin node, edge, double-edge direction, branch, coefficient and index in Tables 1 and 3 is native editable TeX. Literal source quirks, including $\mathfrak g_a$ and unindexed $Q$ in Proposition 2.33, the printed coroot pairing in Proposition 2.40, and the Proposition 2.35 citation in Theorem 4.6, are retained without mathematical repair. The source review checks complete statement/proof correspondence and published content, not independent mathematical refereeing.
\clearpage
\begin{center}{\Large Maximal representations of cocompact complex hyperbolic lattices, a uniform approach}\\[1ex]Pierre-Emmanuel Chaput \& Julien Maubon\end{center}
\section*{Original introductory normalizations}
This paper deals with maximal representations of complex hyperbolic lattices in semisimple Hermitian Lie groups with no compact factors.

A complex hyperbolic lattice $\Gamma$ is a lattice in the Lie group $\mathrm{SU}(1,n)$, a finite cover of the group of biholomorphisms of the $n$-dimensional complex hyperbolic space $\mathbb H^n_{\mathbb C}=\mathrm{SU}(1,n)/\mathrm U(n)$. We shall always assume that our lattice $\Gamma$ is cocompact, or uniform, meaning that the quotient $X:=\Gamma\backslash\mathbb H^n_{\mathbb C}$ is compact. To simplify matters, we also assume in this introduction (except in the statements of our main results) that $\Gamma$ is torsion free, so that $X$ is also a manifold. The $\mathrm{SU}(1,n)$-invariant Kähler form on $\mathbb H^n_{\mathbb C}$ with constant holomorphic sectional curvature $-1$ and the Kähler form it induces on $X$ will both be denoted by $\omega$.

A Lie group $G_{\mathbb R}$ is said to be a \emph{real algebraic Hermitian Lie group} (or a Hermitian group for short) if it is the connected component $\jepPubBmi{G}(\mathbb R)^{\circ}$ of the group of real point of an algebraic group $\jepPubBmi{G}$ defined over $\mathbb R$, if it is semisimple with no compact factors, and if its associated symmetric space $M=G_{\mathbb R}/K_{\mathbb R}$ is a Hermitian symmetric space (of the noncompact type). This means that the noncompact symmetric space $M$ admits a $G_{\mathbb R}$-invariant complex structure, with respect to which the $G_{\mathbb R}$-invariant Riemannian metrics are (necessarily) Kähler. The real rank $\operatorname{rk}_{\mathbb R}G_{\mathbb R}$ of $G_{\mathbb R}$ coincides with the rank $r_M$ of $M$ as a symmetric space, namely the maximal dimension of a flat subspace in $M$. Simple Hermitian groups are classified: there are four infinite families of classical groups, which up to isogeny are $\mathrm{SU}(p,q)$ with $1\leqslant p\leqslant q$, $\mathrm{SO}_0(2,p)$ with $p\geqslant3$, $\mathrm{Sp}(2m,\mathbb R)$ with $m\geqslant2$ and $\mathrm{SO}^{\star}(2m)$ with $m\geqslant4$; and two exceptional groups $\mathrm E_{6(-14)}$ and $\mathrm E_{7(-25)}$. The real ranks of these groups are respectively $p$, $2$, $m$, $\lfloor m/2\rfloor$, $2$ and $3$.

Let $\omega_M$ denote the $G_{\mathbb R}$-invariant Kähler form of $M$, uniquely normalized so that its holomorphic sectional curvatures lie between $-1$ and $-1/r_M$.

If $\rho$ is a representation (a group homomorphism) from $\Gamma$ to $G_{\mathbb R}$, we define its \emph{Toledo invariant} $\tau(\rho)$ as follows:
\[
\tau(\rho)=\frac1{n!}\int_X f^{\star}\omega_M\wedge\omega^{n-1},
\]
where $f:\mathbb H^n_{\mathbb C}\to M$ is a $\jepPubScript{C}^{\infty}$ and $\rho$-equivariant map and $f^{\star}\omega_M$ is seen as a 2-form on $X$ by $\Gamma$-invariance. The Toledo invariant does not depend on the choice of the map $f$, it depends only on $\rho$, and in fact, only on the connected component of $\rho$ in $\operatorname{Hom}(\Gamma,G_{\mathbb R})$. Moreover, it satisfies the following Milnor-Wood type inequality:
\[
|\tau(\rho)|\leqslant r_M\operatorname{vol}(X),
\]
a fundamental property established in full generality in \cite{jep93BI07}.
The \emph{maximal representations} $\rho:\Gamma\to G_{\mathbb R}$ are those representations for which the Milnor-Wood inequality is an equality.

One interesting by-product of this unified approach is that it allows to exclude a priori the possibility of maximal representations in any tube type real algebraic Hermitian Lie group, and in particular in $\operatorname E_{7(-25)}$. Recall that up to isogeny the simple tube type Hermitian groups are $\operatorname{SU}(p,p)$, $\operatorname{SO}_0(2,p)$, $\operatorname{Sp}(2m,\mathbb R)$, $\operatorname{SO}^{\star}(2m)$ with $m$ even, and $\operatorname E_{7(-25)}$. See also Section~2.4.4. Indeed, we prove in Section~5.1 that for tube type targets the representation $\rho$ satisfies a stronger inequality than the Milnor-Wood inequality.
\begin{propositionC}
Let $\Gamma$ be a uniform lattice in $\mathrm{SU}(1,n)$, with $n\geqslant2$, and let $X=\Gamma\backslash\mathbb H^n_{\mathbb C}$. Assume that the real algebraic Hermitian Lie group $G_{\mathbb R}$ has tube type and let $r_M$ be the real rank of $G_{\mathbb R}$. Let $\rho$ be a representation $\Gamma\to G_{\mathbb R}$. Then
\[
|\tau(\rho)|\leqslant\max\Big\{r_M-1,\frac{r_M}{2}\cdot\frac{n+1}{n}\Big\}\operatorname{vol}(X)<r_M\operatorname{vol}(X).
\]
\end{propositionC}

\setcounter{section}{1}
\section{Submodule of a cominuscule representation associated with a nilpotent element}
Here we develop the algebraic material that we will need in Section~4 to give a new and unified proof of the Milnor-Wood inequality.

\subsection{Hermitian symmetric spaces}
We begin by recalling useful facts and definitions concerning Hermitian symmetric spaces and setting up our notation. Good references for what follows are \cite{jep93Wol72,jep93AMRT10,jep93Hel01,jep93Kna02}.

Let $G_{\mathbb R}$ be a simple real algebraic Hermitian Lie group and $K_{\mathbb R}$ a maximal compact subgroup of $G_{\mathbb R}$. By our convention, this means that $G_{\mathbb R}$ is a simple noncompact Lie group whose associated symmetric space $M=G_{\mathbb R}/K_{\mathbb R}$ is a Hermitian symmetric space of the noncompact type ($M$ is irreducible since $G_{\mathbb R}$ is simple), and that moreover $G_{\mathbb R}$ is the connected component of the group of real points $\jepPubBmi{G}(\mathbb R)$ of an algebraic group $\jepPubBmi{G}$ defined over $\mathbb R$. Observe that for example the connected component $\operatorname{Isom}^{\circ}(M)$ of the group of isometries of a Hermitian symmetric space of the noncompact type $M$ is a real algebraic Hermitian Lie group, see e.g. \cite[p.~106]{jep93AMRT10}. The group of complex points $\jepPubBmi{G}(\mathbb C)$ of $\jepPubBmi{G}$ will be denoted by $G$.

We shall first assume that the complex Lie group $G$ is simply connected. This assumption simplifies the exposition of the paper and the arguments of our proofs. We shall see in Remarks~5.14 and~5.15 how to deal with the general case.

The rank of the symmetric space $M$, or equivalently the real rank of $G_{\mathbb R}$, will be denoted by $r_M$. Recall that a $k$-dimensional \emph{flat} in $M$ is the image of a totally geodesic isometric embedding of $\mathbb R^k$ in $M$. By definition the rank $r_M$ of $M$ is the maximal dimension of a flat in $M$. Maximal flats are conjugated under $G_{\mathbb R}$ and maximal flats through the origin $eK_{\mathbb R}$ are conjugated under $K_{\mathbb R}$. Equivalently, in the Hermitian setting, one can consider \emph{polydiscs} in $M$, that is, images of totally geodesic homomorphic embeddings of a polydisc $\Delta^k$ in $M$. Then $r_M$ is also the maximal (complex) dimension of a polydisc in $M$ (maximal polydiscs are complexifications of maximal flats).

The fact that $M$ is an irreducible Hermitian symmetric space is equivalent to the fact that the Lie algebra $\mathfrak k_{\mathbb R}$ of $K_{\mathbb R}$ has a 1-dimensional center $\mathfrak z_{\mathbb R}$ and that the centralizer of $\mathfrak z_{\mathbb R}$ in the Lie algebra $\mathfrak g_{\mathbb R}$ of $G_{\mathbb R}$ is $\mathfrak k_{\mathbb R}$, see e.g. \cite{jep93Kna02}. If $\mathfrak g_{\mathbb R}=\mathfrak k_{\mathbb R}\oplus\mathfrak m_{\mathbb R}$ is a Cartan decomposition of $\mathfrak g_{\mathbb R}$, the complex structure at the origin is given by the adjoint action on $\mathfrak m_{\mathbb R}$ of a suitably normalized element of $\mathfrak z_{\mathbb R}$. It follows that any maximal Abelian subspace $\mathfrak t_{\mathbb R}$ of $\mathfrak k_{\mathbb R}$ contains $\mathfrak z_{\mathbb R}$ and is a Cartan subalgebra of $\mathfrak g_{\mathbb R}$. The corresponding subgroup $T_{\mathbb R}\subset K_{\mathbb R}\subset G_{\mathbb R}$ is a torus.

The Lie algebra $\mathfrak g$ of $G$ is the complexification of the Lie algebra $\mathfrak g_{\mathbb R}$ of $G_{\mathbb R}$. We denote by $\mathfrak z\subset\mathfrak t\subset\mathfrak k\subset\mathfrak g$ the complexifications of the subalgebras $\mathfrak z_{\mathbb R}\subset\mathfrak t_{\mathbb R}\subset\mathfrak k_{\mathbb R}$ and by $Z\subset T\subset K$ the corresponding complex algebraic subgroups of $G$. We also let $\mathfrak m$ be the complexification of $\mathfrak m_{\mathbb R}$, so that $\mathfrak g=\mathfrak k\oplus\mathfrak m$ is a Cartan decomposition of $\mathfrak g$. In particular, $[\mathfrak k,\mathfrak k]\subset\mathfrak k$, $[\mathfrak k,\mathfrak m]\subset\mathfrak m$ and $[\mathfrak m,\mathfrak m]\subset\mathfrak k$.

Let $z\neq0$ be an element in the center $\mathfrak z$ of $\mathfrak k$. Since the adjoint action of $\mathfrak z_{\mathbb R}$ gives the complex structure of $\mathfrak m_{\mathbb R}$, $\operatorname{ad}(z)\jepPubRestrict_{\mathfrak m}$ has exactly two opposite eigenvalues. We choose $z$ so that these eigenvalues are $2$ and $-2$ and we let
\[
\mathfrak m^+=\{x\in\mathfrak m\mid\operatorname{ad}(z)x=2x\},\qquad
\mathfrak m^-=\{x\in\mathfrak m\mid\operatorname{ad}(z)x=-2x\}
\]
be the corresponding eigenspaces. These are abelian subspaces of $\mathfrak m$.

This in turn implies that $\mathfrak p=\mathfrak k\oplus\mathfrak m^+$ is a maximal parabolic subalgebra of $\mathfrak g$. Let $P$ be the corresponding parabolic subgroup of $G$ and $\check M$ the projective variety $G/P$. Then $\check M$ is a symmetric space of the compact type called the compact dual of $M$. Indeed if we let $\check{\mathfrak g}=\mathfrak k_{\mathbb R}\oplus i\mathfrak m_{\mathbb R}$, $\check{\mathfrak g}$ is the Lie algebra of a real compact form $\check G_{\mathbb R}$ of $G$, and $\check M=\check G_{\mathbb R}/K_{\mathbb R}$. The symmetric space $M$ is then an open $G_{\mathbb R}$-orbit in $\check M=G/P$.

\begin{example}
For $p\leqslant q$ let $G_{\mathbb R}=\mathrm{SU}(p,q)$ be the special unitary group of the Hermitian form of signature $(p,q)$ on $\mathbb C^{p+q}$ whose matrix in the canonical basis is $I_{p,q}=\operatorname{diag}(-1,\ldots,-1,1,\ldots,1)$. Then one can choose $K_{\mathbb R}=\mathrm S(\mathrm U(p)\times\mathrm U(q))$ and the symmetric space $M$ identifies with the set of $p$-planes in $\mathbb C^{p+q}$ on which the Hermitian form is negative definite. It has rank $p$. As a bounded symmetric domain, it identifies with $\{Y\in \jepPubScript{M}_{p,q}(\mathbb C)\mid I_p-Y^{\star}Y\text{ is positive definite}\}$. We have $G=\mathrm{SL}(p+q,\mathbb C)$, $K=\mathrm S(\mathrm{GL}(p)\times\mathrm{GL}(q))$, $T$ is the torus of diagonal matrices in $G$, $Z$ is the 1-dimensional subgroup of diagonal matrices of the form $\operatorname{diag}(\lambda^q,\ldots,\lambda^q,\lambda^{-p},\ldots,\lambda^{-p})$ for $\lambda\in\mathbb C^{\star}$, $\mathfrak m\simeq\mathbb C^{pq}\times\mathbb C^{pq}$ is the subspace of block anti-diagonal matrices in $\mathfrak{sl}(p+q,\mathbb C)$. The compact real form of $G$ is $\mathrm{SU}(p+q)$ and the compact dual $\check M$ identifies with the Grassmannian of $p$-planes in $\mathbb C^{p+q}$, of which $M$ is an open subset.
\end{example}

Let $R$ be the set of roots of $G$. For $\alpha\in R$, $\mathfrak g_\alpha\subset\mathfrak g$ is the root space of $\alpha$, $\mathfrak{sl}(\alpha)$ is the Lie subalgebra of $\mathfrak g$ generated by $\mathfrak g_\alpha$ and $\mathfrak g_{-\alpha}$ and $\mathrm{SL}(\alpha)$ the corresponding subgroup of $G$.

A root space $\mathfrak g_\alpha$, $\alpha\in R$, is either a subspace of $\mathfrak k$ or a subspace of $\mathfrak m$, and the root $\alpha$ is said to be \emph{compact} or \emph{noncompact} accordingly. Equivalently, a root $\alpha$ is compact if and only if $\langle\alpha,z\rangle=0$. We have:
\[
\mathfrak k=\mathfrak t\oplus\jepPubOplus_{\alpha:\langle\alpha,z\rangle=0}\mathfrak g_\alpha,\quad
\mathfrak m=\jepPubOplus_{\alpha:\langle\alpha,z\rangle\neq0}\mathfrak g_\alpha,\quad
\mathfrak m^+=\jepPubOplus_{\alpha:\langle\alpha,z\rangle=2}\mathfrak g_\alpha,\quad
\mathfrak m^-=\jepPubOplus_{\alpha:\langle\alpha,z\rangle=-2}\mathfrak g_\alpha.
\]
We choose a basis $\Pi$ of $R$ so that the roots in $R(\mathfrak m^+)$ are positive roots. If $\alpha\in R$, we write $\alpha=\sum_{\beta\in\Pi}n_\beta(\alpha)\beta$ the expression of $\alpha$ in terms of the simple roots. The support $\operatorname{supp}(\alpha)$ of $\alpha\in R$ is the set $\{\beta\in\Pi\mid n_\beta(\alpha)\neq0\}$.

The Weyl group of $R$ is denoted by $W$. If $\alpha\in R$, $\alpha^\jepPubCoroot\in\mathfrak t$ denotes the corresponding coroot: it is defined by the relation $s_\alpha(\beta)=\beta-\langle\beta,\alpha^\jepPubCoroot\rangle\alpha$ for all $\beta\in\mathfrak t^*$, where $s_\alpha\in W$ is the reflexion corresponding to $\alpha$ (in particular, $s_\alpha(\alpha)=-\alpha$ and hence $\langle\alpha,\alpha^\jepPubCoroot\rangle=2$). If $\alpha=\sum_{\beta\in\Pi}n_\beta(\alpha)\beta$ is the expression of a root $\alpha$ in the simple roots, then $\alpha^\jepPubCoroot=\sum_{\beta\in\Pi}n_\beta(\alpha)\frac{\langle\beta,\alpha^\jepPubCoroot\rangle}{\langle\alpha,\beta^\jepPubCoroot\rangle}\beta^\jepPubCoroot$. If $\beta\in\Pi$ is a simple root, $\varpi_\beta$ denotes the fundamental weight corresponding to $\beta$. The set of simple roots $\Pi$ is a basis of $\mathfrak t^{\star}$, $\Pi^\jepPubCoroot=\{\beta^\jepPubCoroot\mid\beta\in\Pi\}$ is a basis of $\mathfrak t$ and $\{\varpi_\beta\mid\beta\in\Pi\}$ is by definition the basis of $\mathfrak t^{\star}$ dual to $\Pi^\jepPubCoroot$.

In this paper, we use the convention that if the root system $R$ of $\mathfrak g$ is simply laced (equivalently, of type A, D or E), then all the roots are long. Therefore short roots exist only if $R$ is not simply laced. There can be at most two root lengths in $R$. Moreover if there are two root lengths the ratio $\frac{\text{long root length}}{\text{short root length}}$ equals $\sqrt2$ because Hermitian symmetric spaces exist only when the type of $G$ is $A$, $B$, $C$, $D$, $E_6$ or $E_7$ (in particular, $G_2$ does not appear). We recall that for $\alpha,\beta\in R$, the ratio $\langle\alpha,\beta^\jepPubCoroot\rangle/\langle\beta,\alpha^\jepPubCoroot\rangle$ equals $(\text{root length of }\alpha)^2/(\text{root length of }\beta)^2$.

Linearly independent positive roots $\delta_1,\ldots,\delta_s$ are said to be \emph{strongly orthogonal} if for all $i\neq j$, $\delta_i\pm\delta_j$ is not a root. By \cite{jep93HC56}, or \cite[Ch.~VIII, \S7]{jep93Hel01}, we have
\begin{fact}
All maximal sets of noncompact strongly orthogonal roots have cardinality $r_M$.
\end{fact}
If $(\delta_1,\ldots,\delta_{r_M})$ is such a set then $\mathfrak g_{-\delta_1}\oplus\cdots\oplus\mathfrak g_{-\delta_{r_M}}$ is the tangent space to a maximal polydisc through the origin in $M$.

\subsection{The cominuscule representation of $G$ associated with $M$}
\subsubsection{The representation}
\begin{definition}
A simple root $\beta\in\Pi$, is \emph{cominuscule} w.r.t. the root system $R$ if $n_\beta(\alpha)\in\{-1,0,1\}$ for all $\alpha\in R$, or equivalently if $n_\beta(\Theta)=1$, where $\Theta$ is the highest root of $R$.

An irreducible representation of $G$ whose highest weight is equal to $\varpi_\beta$ for $\beta$ a cominuscule root of $R$ is called a \emph{cominuscule representation}.
\end{definition}
The following is well-known, but we include a proof.
\begin{proposition}
There is a unique simple noncompact root. This root is long and cominuscule.
\end{proposition}
\begin{notation}
This unique simple noncompact root will be denoted by $\zeta$.
\end{notation}
\begin{proof}
Since $R$ is irreducible, the highest root $\Theta$ of $R$ is unique and $\Theta$ has a positive coefficient on every simple root. Now, if there are more than one simple noncompact root, or if $n_\zeta(\Theta)>1$ for some noncompact simple root $\zeta$, then $\langle\Theta,z\rangle>2$. Contradiction. Hence $\zeta$ is unique and cominuscule. Assume that $\zeta$ is short (so that $R$ is not simply laced). Then, since $\Theta$ is long,
\[
\langle\Theta,\zeta^\jepPubCoroot\rangle=\langle\zeta,\Theta^\jepPubCoroot\rangle\frac{(\text{root length of }\Theta)^2}{(\text{root length of }\zeta)^2}\geqslant2.
\]
Since $s_\zeta(\Theta)$ is a positive root, this implies that $n_\zeta(\Theta)\geqslant2$, contradicting the fact that $\zeta$ is cominuscule.
\end{proof}

Note that the maximal parabolic subalgebra $\mathfrak p=\mathfrak k\oplus\mathfrak m^+$ of $\mathfrak g$ is the standard parabolic defined by the root $\zeta$:
\[
\mathfrak p=\mathfrak t\oplus\jepPubOplus_{\alpha:n_\zeta(\alpha)\geqslant0}\mathfrak g_\alpha.
\]
Recall that $P$ is the corresponding parabolic subgroup of $G$ and that the compact dual of $M$ is $\check M=G/P$. By \cite{jep93Mur59}, since $G$ is simply connected, the Picard group $\operatorname{Pic}(G/P)$ is isomorphic to the group of characters $\jepPubScript{X}(P)$ of $P$. Since $\mathfrak p=\mathfrak z\oplus[\mathfrak p,\mathfrak p]$, $\jepPubScript{X}(P)$ is isomorphic to $\mathbb Z$ and thus it is generated by the smallest positive character of $P$, namely $\varpi:=\varpi_\zeta$. Here we embed $\jepPubScript{X}(P)$ in $\mathfrak t^*$ via $\lambda\mapsto d\lambda$, the differential of $\lambda$ at the unit element of $T$. Moreover, the isomorphism $\iota:\jepPubScript{X}(P)\simeq\operatorname{Pic}(G/P)$ is given by $\lambda\mapsto(G\times\mathbb C_\lambda)/P$, where $\mathbb C_\lambda$ denotes the $1$-dimensional $P$-module defined by $\lambda$. In particular, $\jepPubScript{L}:=\iota(\varpi)$ is a generator of $\operatorname{Pic}(\check M)$.

Let now $\jepPubBbb{E}$ be the cominuscule irreducible representation of $G$ whose highest weight is $\varpi=\varpi_\zeta$. Let $\jepPubBbb{E}_\varpi$ be the $\varpi$-eigenspace of $\jepPubBbb{E}$. Then $\jepPubBbb{E}_\varpi$ is 1-dimensional and its stabilizer in $G$ is $P$. This gives a $G$-equivariant holomorphic and isometric embedding of $\check M=G/P$ in the projective space $\mathbb P\jepPubBbb{E}$. It is called the \emph{first canonical embedding} of $\check M$, and since $\jepPubBbb{E}_\varpi\simeq\mathbb C_\varpi$ as $P$-modules, we have $\jepPubScript{L}\simeq\jepPubScript{O}_{\mathbb P\jepPubBbb{E}}(-1)$. See e.g. \cite{jep93NT76} for more details.

\begin{remark}
The complex group $G$ may have several cominuscule simple roots or representations. In fact, cominuscule simple roots are in one to one correspondence with Hermitian real forms $G_{\mathbb R}$ (or $\check G_{\mathbb R}$) of $G$. When we will need a case by case argument, we will rather use the classification of the cominuscule roots than the classification of the real groups $G_{\mathbb R}$. This correspondence is shown in the first two columns of Table~1.
\end{remark}
We now begin our study of the cominuscule representation $\jepPubBbb{E}$ of $G$.
\begin{notation}
We denote by $\mu_0$ the lowest weight of $\jepPubBbb{E}$ and by $X(\jepPubBbb{E})$ the set of weights of $\jepPubBbb{E}$. For $\chi\in X(\jepPubBbb{E})$, we write $\jepPubBbb{E}_\chi$ for the corresponding weight space. Recall that $\varpi$ is the highest weight of $\jepPubBbb{E}$.
\end{notation}
The fact that $\jepPubBbb{E}$ is cominuscule has the following consequence on the weights of $\jepPubBbb{E}$.
\begin{lemma}
For any weight $\chi$ of $\jepPubBbb{E}$, and any root $\alpha\in R$, $\langle\chi,\alpha^\jepPubCoroot\rangle\in\{-2,-1,0,1,2\}$, and the equality $\langle\chi,\alpha^\jepPubCoroot\rangle=\pm2$ implies that $\alpha$ is short (hence that $R$ is not simply laced).
\end{lemma}
\begin{proof}
For the highest weight $\varpi$ of $\jepPubBbb{E}$, the results follows from the facts that $\langle\varpi,\alpha^\jepPubCoroot\rangle=n_\zeta(\alpha)\langle\zeta,\alpha^\jepPubCoroot\rangle/\langle\alpha,\zeta^\jepPubCoroot\rangle$ and $\zeta$ is a long root. The result still holds if $\varpi$ is replaced by $w\cdot\varpi$, where $w\in W$ is arbitrary, and since any weight of $\jepPubBbb{E}$ is in the convex hull of $W\cdot\varpi$, it holds for any weight.
\end{proof}
We deduce that the structure of $\jepPubBbb{E}$ with respect to the action of $\mathfrak g_\alpha$ for $\alpha$ a long root is particularly simple.
\begin{lemma}
Let $\alpha$ be a long root and let $\chi$ be a weight of $\jepPubBbb{E}$. We have
\[
\mathfrak g_{-\alpha}\cdot \jepPubBbb{E}_\chi=\begin{cases}\jepPubBbb{E}_{\chi-\alpha}&\text{if }\langle\chi,\alpha^\jepPubCoroot\rangle=1,\\\{0\}&\text{otherwise}\end{cases}.
\]
\end{lemma}
\begin{proof}
Let $\alpha$ be long and let $\mathfrak{sl}(\alpha)$ be the Lie subalgebra of $\mathfrak g$ isomorphic to $\mathfrak{sl}_2$ corresponding to $\alpha$. Let $S\subset \jepPubBbb{E}$ be the $\mathfrak{sl}(\alpha)$-submodule generated by $\jepPubBbb{E}_\chi$. By Lemma~2.8 and the representation theory of $\mathfrak{sl}_2$, any irreducible component $V$ of $S$ is an $\mathfrak{sl}(\alpha)$-module of dimension $1$ or $2$.

We therefore have only three possibilities. The first case is when $V=V_\chi$, $\mathfrak g_\alpha\cdot V_\chi=\{0\}$ and $\mathfrak g_{-\alpha}\cdot V_\chi=\{0\}$. In this case, $\langle\chi,\alpha^\jepPubCoroot\rangle=0$. The second case is when $V=V_\chi\oplus V_{\chi-\alpha}$, $\mathfrak g_\alpha\cdot V_\chi=\{0\}$ and $\mathfrak g_{-\alpha}\cdot V_\chi=V_{\chi-\alpha}$. In this case, $\langle\chi,\alpha^\jepPubCoroot\rangle=1$ (and $\langle\chi-\alpha,\alpha^\jepPubCoroot\rangle=-1$). The third (symmetric) case is when $V=V_\chi\oplus V_{\chi+\alpha}$, $\mathfrak g_\alpha\cdot V_\chi=V_{\chi+\alpha}$ and $\mathfrak g_{-\alpha}\cdot V_\chi=\{0\}$. In this case, $\langle\chi,\alpha^\jepPubCoroot\rangle=-1$ (and $\langle\chi+\alpha,\alpha^\jepPubCoroot\rangle=1$).

If $\langle\chi,\alpha^\jepPubCoroot\rangle=1$, we deduce that $S=S_\chi\oplus S_{\chi-\alpha}$ and that $\mathfrak g_{-\alpha}\cdot S_\chi=S_{\chi-\alpha}$. We have $s_\alpha(\chi)=\chi-\alpha$ so $\dim(\jepPubBbb{E}_\chi)=\dim(\jepPubBbb{E}_{\chi-\alpha})$. The lemma is proved in this case. If $\langle\chi,\alpha^\jepPubCoroot\rangle\leqslant0$, we see that $\mathfrak g_{-\alpha}\cdot \jepPubBbb{E}_\chi=\{0\}$.
\end{proof}

\subsubsection{The grading}
\begin{notation}
Let $z_{\mathtt{max}}$ denote the number $\langle\varpi,z\rangle$.
\end{notation}
\begin{remark}
The precise value of $z_{\mathtt{max}}$ will not play any role in the following, however it has already been computed in \cite[p.~214-216]{jep93KM10}): with our choice of $z\in\mathfrak z$ we have $z_{\mathtt{max}}=2\dim\check M/c_1(\check M)$, where $c_1(\check M)$ is the first Chern number of $\check M$. For example if $G_{\mathbb R}=\mathrm{SU}(p,q)$, $z_{\mathtt{max}}=2pq/(p+q)$.
\end{remark}
\begin{proposition}
The set $\{\langle\chi,z\rangle\mid\chi\in X(\jepPubBbb{E})\}$ is the set
\[
\{z_{\mathtt{max}},z_{\mathtt{max}}-2,\ldots,z_{\mathtt{max}}-2r_M\}.
\]
\end{proposition}
\begin{proof}
It follows from \cite[Th.~2.1]{jep93RRS92} that the $W$-orbit of the weight $\varpi$ is exactly the set of weights of the form $\varpi-\sum_{i=1}^k\delta_i$, where $(\delta_i)_{1\leqslant i\leqslant k}$ is a family of noncompact long strongly orthogonal roots. In fact, the orbit $W\cdot\varpi$ is in bijection with $W/W_P$, so the equivalence between items~(c) and~(e) in the cited theorem proves the claim. For any $i$, we have $\langle\delta_i,z\rangle=2$, thus we have the equality of sets
\[
\{\langle\mu,z\rangle\mid\mu\in W\cdot\varpi\}=\{z_{\mathtt{max}},z_{\mathtt{max}}-2,\ldots,z_{\mathtt{max}}-2r_M\}.
\]
In particular, $\langle\mu_0,z\rangle=z_{\mathtt{max}}-2r_M$ and for $\chi\in X(\jepPubBbb{E})$, we have $z_{\mathtt{max}}-2r_M\leqslant\langle\chi,z\rangle\leqslant z_{\mathtt{max}}$. The result of the proposition now follows from the fact that $2$ is a divisor of $\langle\alpha,z\rangle$ for any root $\alpha$.
\end{proof}
Now we can introduce the grading of $\jepPubBbb{E}$.
\begin{definition}
For a relative integer $i$, let
\[
\jepPubBbb{E}_i:=\jepPubOplus_{\chi:\langle\chi,z\rangle=z_{\mathtt{max}}-2i}\jepPubBbb{E}_\chi.
\]
\end{definition}
This grading corresponds to the decomposition of $\jepPubBbb{E}$ into irreducible $K$-modules.
\begin{proposition}
The $K$-modules $\jepPubBbb{E}_i$ are irreducible.
\end{proposition}
\begin{proof}
This might be well-known to experts, but we include a proof for completeness. We give a case by case argument and use Table~1.

In type $A_{n-1}$, we have $\jepPubBbb{E}=\wedge^{r_M}(\mathbb C^{r_M}\oplus\mathbb C^{n-r_M})$ and thus $\jepPubBbb{E}_i=\wedge^{r_M-i}\mathbb C^{r_M}\otimes\wedge^i\mathbb C^{n-r_M}$: this is an irreducible $\mathrm S(\mathrm{GL}_{r_M}\times\mathrm{GL}_{n-r_M})$-module.

In type $B_n$, we have $\jepPubBbb{E}=\mathbb C^{2n+1}=\mathbb C\oplus\mathbb C^{2n-1}\oplus\mathbb C$, and each summand is an irreducible $\mathrm{Spin}_{2n-1}$-module, hence an irreducible $K$-module. In type $(D_n,\varpi_1)$ the situation is similar.

In type $(D_n,\varpi_n)$, $\jepPubBbb{E}$ is the spinor representation of the spin group and, according to \cite{jep93Che97}, we have $\jepPubBbb{E}=\wedge^0\mathbb C^n\oplus\wedge^2\mathbb C^n\oplus\cdots\oplus\wedge^{2p}\mathbb C^n$ (where $p=\lfloor n/2\rfloor$). Thus $\jepPubBbb{E}_i=\wedge^{2i}\mathbb C^n$, and this is an irreducible $\mathrm{GL}_n$-module.

For the types $E_6$ and $E_7$, we use models of these exceptional Lie algebras and their minuscule representations, as given for example in \cite{jep93Man06}. In type $E_6$, we have $\jepPubBbb{E}=\mathbb C\oplus V_{16}\oplus V_{10}$, where $V_{16}$ is a spinor representation and $V_{10}$ the vector representation of the spin group $\mathrm{Spin}_{10}$. In type $E_7$, we have $\jepPubBbb{E}=\mathbb C\oplus V_{27}\oplus V'_{27}\oplus\mathbb C$, where $V_{27}$ and $V'_{27}$ are the two minuscule representations of a group of type $E_6$. In both cases, the $K$-modules $\jepPubBbb{E}_i$ are irreducible.

We now deal with the case of type $C_n$. The following representation-theoretic argument has been suggested to us by an anonymous referee that we would like to thank. We denote by $\mathbb C^{2n}=\mathbb C_a^n\oplus\mathbb C_b^n$ a symplectic $2n$-dimensional space, with $\mathbb C_a^n$ and $\mathbb C_b^n$ supplementary isotropic subspaces. We then have $\jepPubBbb{E}=\bigl(\wedge^n\mathbb C^{2n}\bigr)_\omega$, where the symbol $\omega$ means that we take in $\wedge^n\mathbb C^{2n}$ the irreducible $\mathrm{Sp}_{2n}$-submodule containing the highest weight line $\wedge^n\mathbb C_a^n$. More precisely, by \cite[Th.~17.5]{jep93FH91}, we have a decomposition of $\wedge^n\mathbb C^{2n}$ as an $\mathrm{Sp}_{2n}$-module as follows: $\wedge^n\mathbb C^{2n}=\jepPubBbb{E}\oplus\omega\wedge(\wedge^{n-2}\mathbb C^{2n})$.

Thus, $\wedge^{n-i}\mathbb C_a^n\otimes\wedge^i\mathbb C_b^n=\jepPubBbb{E}_i\oplus\omega\wedge(\wedge^{n-i-1}\mathbb C_a^n)\otimes\wedge^{i-1}\mathbb C_b^n.$ We now consider this as a representation of $K=\mathrm{GL}_n$. As $K$-modules, we have $\mathbb C_b^n\simeq(\mathbb C_a^n)^\star$. We thus get $\wedge^{n-i}\mathbb C_a^n\otimes\wedge^{n-i}\mathbb C_a^n\simeq\jepPubBbb{E}_i\oplus\wedge^{n-i-1}\mathbb C_a^n\otimes\wedge^{n-i+1}\mathbb C_a^n.$ From this last equation, it follows that $\jepPubBbb{E}_i$ is the Cartan square of $\wedge^{n-i}\mathbb C_a^n$ and therefore it is an irreducible $\mathrm{GL}_n$-module.
\end{proof}
\begin{proposition}
We have the following properties:
\begin{enumerate}[label=(\alph*)]
\item $\jepPubBbb{E}=\jepPubBbb{E}_0\oplus\jepPubBbb{E}_1\oplus\cdots\oplus\jepPubBbb{E}_{r_M}$
\item $\jepPubBbb{E}_0=\jepPubBbb{E}_\varpi$
\item $\jepPubBbb{E}_{i+1}=\mathfrak m^-\cdot\jepPubBbb{E}_i$
\item The map $\jepPubBbb{E}_0\otimes\mathfrak m^-\to\jepPubBbb{E}_1$ is an isomorphism
\end{enumerate}
\end{proposition}
\begin{proof}
Only the last two points need a proof. Let $U(\mathfrak m^-)$ denote the enveloping algebra of $\mathfrak m^-$. The third point follows from the fact that $\jepPubBbb{E}=U(\mathfrak m^-)\cdot\jepPubBbb{E}_\varpi$ and the fact that for $\alpha$ a root of $\mathfrak m^-$, we have $\langle\alpha,z\rangle=-2$. The last point follows by Schur’s lemma since $\mathfrak m^-$ and $\jepPubBbb{E}_1$ are irreducible $\mathfrak k$-modules and $\jepPubBbb{E}_0$ is 1-dimensional.
\end{proof}
\begin{remark}
Those statements are well-known. In the case where $\jepPubBbb{E}$ is of tube type, they are proved in \cite[Prop.~5.2]{jep93Gro94}.
\end{remark}

\subsection{Dominant orthogonal sequences}
In \cite{jep93Kos12}, Kostant introduced his so-called ``chain cascade'' of orthogonal roots. Here we will need a version of his algorithm where we impose that all the roots of the chain cascade have a positive coefficient on $\zeta$. Note that a similar algorithm is used in \cite{jep93BM15}.

We define an integer $q$ and, for any integer $i$ such that $1\leqslant i\leqslant q$, a root $\alpha_i$ together with the subset $\Pi_i\subset\Pi$, by the following inductive process:
\begin{itemize}
\item We let $\Pi_1=\Pi$
\item Assuming that $\alpha_1,\ldots,\alpha_{i-1}$ and $\Pi_1,\ldots,\Pi_i$ have been defined, we let $\alpha_i$ be the highest root of the root system $R(\Pi_i)$ generated by $\Pi_i$
\item We let $\Sigma_i\subset\Pi_i$ be the set of simple roots $\beta$ such that $\langle\alpha_i^\jepPubCoroot,\beta\rangle\neq0$
\item If $\zeta\in\Sigma_i$, then $q=i$ and the algorithm terminates. Otherwise, $\Pi_{i+1}$ is the connected component of $\Pi_i\smallsetminus\Sigma_i$ containing $\zeta$
\end{itemize}
\begin{definition}
If $(\alpha_i)_{1\leqslant i\leqslant q}$ is the sequence defined by this process, we say that it is the \emph{maximal dominant orthogonal sequence} for $\varpi$. More generally, the sequences $(\alpha_i)_{1\leqslant i\leqslant r}$ for $1\leqslant r\leqslant q$ are called the \emph{dominant orthogonal sequences}.
\end{definition}
\begin{example}
In type $A_{n-1}$ with the standard base $(\beta_1,\ldots,\beta_{n-1})$, and for $\zeta=\beta_p$ with $2p\leqslant n$, we get $\Pi_i=\{\beta_i,\ldots,\beta_{n-i}\}$ and the maximal dominant orthogonal sequence is $(\sum_{i=1}^{n-1}\beta_i,\sum_{i=2}^{n-2}\beta_i,\ldots,\sum_{i=p}^{n-p}\beta_i),$ so $q=p$.
\end{example}
For later use, we record in Table~1 below what are the dominant orthogonal sequences $(\alpha_1,\ldots,\alpha_r)$ in all cases. The number $r$ varies from $1$ to $r_M$, see Proposition~2.19~(2).

The following proposition essentially adapts the results of \cite{jep93Kos12} to our context and explains our terminology.
\begin{proposition}
Let $(\alpha_i)_{1\leqslant i\leqslant q}$ be the maximal dominant orthogonal sequence for $\varpi$. Then
\begin{enumerate}[label=(\arabic*)]
\item The roots $\alpha_i$ are long and strongly orthogonal
\item $q=r_M$
\item For any integer $i\leqslant r_M$, $\alpha_1^\jepPubCoroot+\cdots+\alpha_i^\jepPubCoroot$ is a dominant coweight: for all $\beta$ in $\Pi$, we have $\langle\beta,\alpha_1^\jepPubCoroot+\cdots+\alpha_i^\jepPubCoroot\rangle\geqslant0$
\item $\varpi-\alpha_1-\cdots-\alpha_{r_M}$ is the lowest weight $\mu_0$ of the irreducible $G$-module $\jepPubBbb{E}$
\end{enumerate}
\end{proposition}

\begin{table}[htbp]\centering\footnotesize
\setlength{\tabcolsep}{2pt}\renewcommand{\arraystretch}{1.3}
\begin{tabular}{|c@{\quad}c|c|c|}\hline
$(G,\varpi)$ & $\mathfrak g_{\mathbb R}$, $r_M$ & $\alpha_1^{\jepPubCoroot}+\cdots+\alpha_r^{\jepPubCoroot}$ & $\alpha_r$\\\hline
\multirow{2}{*}{\begin{tabular}{c}$(A_{n-1},\varpi_p)$, $2p\leqslant n$\\[2pt]\jepGraph{0/0/0/{}/{1}/0,1/1.55/0/{}/{}/0,2/2.55/0/{}/{p}/1,3/3.55/0/{}/{}/0,4/5.1/0/{}/{n-1}/0}{1/2,2/3}{0/1,3/4}{}{}\end{tabular}} & \multirow{2}{*}{\begin{tabular}{c}$\mathfrak{su}(p,n-p)$\\[4pt]$r_M=p$\end{tabular}} & \begin{tabular}{r@{\;}c}$\text{if }2r<n:$ & \jepGraph{0/0/0/{0}/{}/0,1/1.55/0/{0}/{}/0,2/2.55/0/{1}/{r}/0,3/3.55/0/{0}/{}/0,4/5.1/0/{0}/{}/0,5/6.1/0/{1}/{n-r}/0,6/7.1/0/{0}/{}/0,7/8.65/0/{0}/{}/0}{1/2,2/3,4/5,5/6}{0/1,3/4,6/7}{}{}\end{tabular} & \jepGraph{0/0/0/{0}/{}/0,1/1.55/0/{0}/{}/0,2/2.55/0/{1}/{r}/0,3/4.1/0/{1}/{n-r}/0,4/5.1/0/{0}/{}/0,5/6.65/0/{0}/{}/0}{1/2,3/4}{0/1,2/3,4/5}{}{}\rule[-12pt]{0pt}{36pt}\\
\cline{3-4}
 & & \begin{tabular}{r@{\;}c}$\text{if }2r=n:$ & \jepGraph{0/0/0/{0}/{}/0,1/1.55/0/{0}/{}/0,2/2.55/0/{2}/{p}/0,3/3.55/0/{0}/{}/0,4/5.1/0/{0}/{}/0}{1/2,2/3}{0/1,3/4}{}{}\end{tabular} & \jepGraph{0/0/0/{0}/{}/0,1/1.55/0/{0}/{}/0,2/2.55/0/{1}/{p}/0,3/3.55/0/{0}/{}/0,4/5.1/0/{0}/{}/0}{1/2,2/3}{0/1,3/4}{}{}\rule[-12pt]{0pt}{36pt}\\
\hline
\multirow{2}{*}{\begin{tabular}{c}$(B_n,\varpi_1)$\\[2pt]\jepGraph{0/0/0/{}/{1}/1,1/1/0/{}/{}/0,2/2.55/0/{}/{}/0,3/3.55/0/{}/{}/0,4/4.55/0/{}/{n}/0}{0/1,2/3}{1/2}{3/4/1}{}\end{tabular}} & \multirow{2}{*}{\begin{tabular}{c}$\mathfrak{so}(2,2n-1)$\\[4pt]$r_M=2$\end{tabular}} & \begin{tabular}{r@{\;}c}$\text{if }r=1:$ & \jepGraph{0/0/0/{0}/{}/0,1/1/0/{1}/{}/0,2/2/0/{0}/{}/0,3/3.55/0/{0}/{}/0,4/4.55/0/{0}/{}/0,5/5.55/0/{0}/{}/0}{0/1,1/2,3/4}{2/3}{4/5/1}{}\end{tabular} & \jepGraph{0/0/0/{1}/{}/0,1/1/0/{2}/{}/0,2/2.55/0/{2}/{}/0,3/3.55/0/{2}/{}/0,4/4.55/0/{2}/{}/0}{0/1,2/3}{1/2}{3/4/1}{}\rule[-12pt]{0pt}{36pt}\\
\cline{3-4}
 & & \begin{tabular}{r@{\;}c}$\text{if }r=2:$ & \jepGraph{0/0/0/{2}/{}/0,1/1/0/{0}/{}/0,2/2.55/0/{0}/{}/0,3/3.55/0/{0}/{}/0,4/4.55/0/{0}/{}/0}{0/1,2/3}{1/2}{3/4/1}{}\end{tabular} & \jepGraph{0/0/0/{1}/{}/0,1/1/0/{0}/{}/0,2/2.55/0/{0}/{}/0,3/3.55/0/{0}/{}/0,4/4.55/0/{0}/{}/0}{0/1,2/3}{1/2}{3/4/1}{}\rule[-12pt]{0pt}{36pt}\\
\hline
\multirow{2}{*}{\begin{tabular}{c}$(C_n,\varpi_n)$\\[2pt]\jepGraph{0/0/0/{}/{1}/0,1/1.55/0/{}/{}/0,2/2.55/0/{}/{}/0,3/3.55/0/{}/{n}/1}{1/2}{0/1}{2/3/-1}{}\end{tabular}} & \multirow{2}{*}{\begin{tabular}{c}$\mathfrak{sp}(2n,\mathbb R)$\\[4pt]$r_M=n$\end{tabular}} & \begin{tabular}{r@{\;}c}$\text{if }r<n:$ & \jepGraph{0/0/0/{0}/{}/0,1/1.55/0/{0}/{}/0,2/2.55/0/{2}/{r}/0,3/3.55/0/{0}/{}/0,4/5.1/0/{0}/{}/0,5/6.1/0/{0}/{}/0,6/7.1/0/{0}/{}/0}{1/2,2/3,4/5}{0/1,3/4}{5/6/-1}{}\end{tabular} & \jepGraph{0/0/0/{0}/{}/0,1/1.55/0/{0}/{}/0,2/2.55/0/{2}/{r}/0,3/4.1/0/{2}/{}/0,4/5.1/0/{2}/{}/0,5/6.1/0/{1}/{}/0}{1/2,3/4}{0/1,2/3}{4/5/-1}{}\rule[-12pt]{0pt}{36pt}\\
\cline{3-4}
 & & \begin{tabular}{r@{\;}c}$\text{if }r=n:$ & \jepGraph{0/0/0/{0}/{}/0,1/1.55/0/{0}/{}/0,2/2.55/0/{0}/{}/0,3/3.55/0/{2}/{}/0}{1/2}{0/1}{2/3/-1}{}\end{tabular} & \jepGraph{0/0/0/{0}/{}/0,1/1.55/0/{0}/{}/0,2/2.55/0/{0}/{}/0,3/3.55/0/{1}/{}/0}{1/2}{0/1}{2/3/-1}{}\rule[-12pt]{0pt}{36pt}\\
\hline
\multirow{2}{*}{\begin{tabular}{c}$(D_n,\varpi_1)$\\[2pt]\jepGraph{0/0/0/{}/{1}/1,1/1/0/{}/{}/0,2/2.55/0/{}/{}/0,3/3.05/0.8/{}/{n-1}/0,4/3.05/-0.8/{}/{n}/0}{0/1,2/3,2/4}{1/2}{}{}\end{tabular}} & \multirow{2}{*}{\begin{tabular}{c}$\mathfrak{so}(2,2n-2)$\\[4pt]$r_M=2$\end{tabular}} & \begin{tabular}{r@{\;}c}$\text{if }r=1:$ & \jepGraph{0/0/0/{0}/{}/0,1/1/0/{1}/{}/0,2/2/0/{0}/{}/0,3/3.55/0/{0}/{}/0,4/4.55/0/{0}/{}/0,5/5.05/0.8/{0}/{}/0,6/5.05/-0.8/{0}/{}/0}{0/1,1/2,3/4,4/5,4/6}{2/3}{}{}\end{tabular} & \jepGraph{0/0/0/{1}/{}/0,1/1/0/{2}/{}/0,2/2.55/0/{2}/{}/0,3/3.55/0/{2}/{}/0,4/4.05/0.8/{1}/{}/0,5/4.05/-0.8/{1}/{}/0}{0/1,2/3,3/4,3/5}{1/2}{}{}\rule[-12pt]{0pt}{36pt}\\
\cline{3-4}
 & & \begin{tabular}{r@{\;}c}$\text{if }r=2:$ & \jepGraph{0/0/0/{2}/{}/0,1/1/0/{0}/{}/0,2/2.55/0/{0}/{}/0,3/3.55/0/{0}/{}/0,4/4.05/0.8/{0}/{}/0,5/4.05/-0.8/{0}/{}/0}{0/1,2/3,3/4,3/5}{1/2}{}{}\end{tabular} & \jepGraph{0/0/0/{1}/{}/0,1/1/0/{0}/{}/0,2/2.55/0/{0}/{}/0,3/3.55/0/{0}/{}/0,4/4.05/0.8/{0}/{}/0,5/4.05/-0.8/{0}/{}/0}{0/1,2/3,3/4,3/5}{1/2}{}{}\rule[-12pt]{0pt}{36pt}\\
\hline
\multirow{3}{*}{\begin{tabular}{c}$(D_n,\varpi_n)$\\[2pt]\jepGraph{0/0/0/{}/{1}/0,1/1/0/{}/{}/0,2/2.55/0/{}/{}/0,3/3.05/0.8/{}/{n-1}/0,4/3.05/-0.8/{}/{n}/1}{0/1,2/3,2/4}{1/2}{}{}\end{tabular}} & \multirow{3}{*}{\begin{tabular}{c}$\mathfrak{so}^{\star}(2n)$\\[4pt]$r_M=\lfloor n/2\rfloor$\end{tabular}} & \begin{tabular}{r@{\;}c}$\text{if }2r\leqslant n-2:$ & \jepGraph{0/0/0/{0}/{}/0,1/1.55/0/{0}/{}/0,2/2.55/0/{1}/{2r}/0,3/3.55/0/{0}/{}/0,4/5.1/0/{0}/{}/0,5/6.1/0/{0}/{}/0,6/6.6/0.8/{0}/{}/0,7/6.6/-0.8/{0}/{}/0}{1/2,2/3,4/5,5/6,5/7}{0/1,3/4}{}{}\end{tabular} & \jepGraph{0/0/0/{0}/{}/0,1/1.55/0/{0}/{}/0,2/2.55/0/{1}/{}/0,3/3.55/0/{2}/{2r}/0,4/5.1/0/{2}/{}/0,5/6.1/0/{2}/{}/0,6/6.6/0.8/{1}/{}/0,7/6.6/-0.8/{1}/{}/0}{1/2,2/3,4/5,5/6,5/7}{0/1,3/4}{}{}\rule[-12pt]{0pt}{36pt}\\
\cline{3-4}
 & & \begin{tabular}{r@{\;}c}$\text{if }2r=n-1:$ & \jepGraph{0/0/0/{0}/{}/0,1/1.55/0/{0}/{}/0,2/2.55/0/{0}/{}/0,3/3.05/0.8/{1}/{}/0,4/3.05/-0.8/{1}/{}/0}{1/2,2/3,2/4}{0/1}{}{}\end{tabular} & \jepGraph{0/0/0/{0}/{}/0,1/1.55/0/{0}/{}/0,2/2.55/0/{1}/{}/0,3/3.05/0.8/{1}/{}/0,4/3.05/-0.8/{1}/{}/0}{1/2,2/3,2/4}{0/1}{}{}\rule[-12pt]{0pt}{36pt}\\
\cline{3-4}
 & & \begin{tabular}{r@{\;}c}$\text{if }2r=n:$ & \jepGraph{0/0/0/{0}/{}/0,1/1.55/0/{0}/{}/0,2/2.55/0/{0}/{}/0,3/3.05/0.8/{0}/{}/0,4/3.05/-0.8/{2}/{}/0}{1/2,2/3,2/4}{0/1}{}{}\end{tabular} & \jepGraph{0/0/0/{0}/{}/0,1/1.55/0/{0}/{}/0,2/2.55/0/{0}/{}/0,3/3.05/0.8/{0}/{}/0,4/3.05/-0.8/{1}/{}/0}{1/2,2/3,2/4}{0/1}{}{}\rule[-12pt]{0pt}{36pt}\\
\hline
\multirow{2}{*}{\begin{tabular}{c}$(E_6,\varpi_1)$\\[2pt]\jepGraph{0/0/0/{}/{1}/1,1/1/0/{}/{}/0,2/2/0/{}/{}/0,3/3/0/{}/{}/0,4/4/0/{}/{}/0,5/2/1/{}/{}/0}{0/1,1/2,2/3,3/4,2/5}{}{}{}\end{tabular}} & \multirow{2}{*}{\begin{tabular}{c}$\mathfrak e_{6(-14)}$\\[4pt]$r_M=2$\end{tabular}} & \begin{tabular}{r@{\;}c}$\text{if }r=1:$ & \jepGraph{0/0/0/{0}/{}/0,1/1/0/{0}/{}/0,2/2/0/{0}/{}/0,3/3/0/{0}/{}/0,4/4/0/{0}/{}/0,5/2/1/{1}/{}/0}{0/1,1/2,2/3,3/4,2/5}{}{}{}\end{tabular} & \jepGraph{0/0/0/{1}/{}/0,1/1/0/{2}/{}/0,2/2/0/{3}/{}/0,3/3/0/{2}/{}/0,4/4/0/{1}/{}/0,5/2/1/{2}/{}/0}{0/1,1/2,2/3,3/4,2/5}{}{}{}\rule[-12pt]{0pt}{36pt}\\
\cline{3-4}
 & & \begin{tabular}{r@{\;}c}$\text{if }r=2:$ & \jepGraph{0/0/0/{1}/{}/0,1/1/0/{0}/{}/0,2/2/0/{0}/{}/0,3/3/0/{0}/{}/0,4/4/0/{1}/{}/0,5/2/1/{0}/{}/0}{0/1,1/2,2/3,3/4,2/5}{}{}{}\end{tabular} & \jepGraph{0/0/0/{1}/{}/0,1/1/0/{1}/{}/0,2/2/0/{1}/{}/0,3/3/0/{1}/{}/0,4/4/0/{1}/{}/0,5/2/1/{0}/{}/0}{0/1,1/2,2/3,3/4,2/5}{}{}{}\rule[-12pt]{0pt}{36pt}\\
\hline
\multirow{3}{*}{\begin{tabular}{c}$(E_7,\varpi_7)$\\[2pt]\jepGraph{0/0/0/{}/{}/0,1/1/0/{}/{}/0,2/2/0/{}/{}/0,3/3/0/{}/{}/0,4/4/0/{}/{}/0,5/5/0/{}/{7}/1,6/2/1/{}/{}/0}{0/1,1/2,2/3,3/4,4/5,2/6}{}{}{}\end{tabular}} & \multirow{3}{*}{\begin{tabular}{c}$\mathfrak e_{7(-25)}$\\[4pt]$r_M=3$\end{tabular}} & \begin{tabular}{r@{\;}c}$\text{if }r=1:$ & \jepGraph{0/0/0/{1}/{}/0,1/1/0/{0}/{}/0,2/2/0/{0}/{}/0,3/3/0/{0}/{}/0,4/4/0/{0}/{}/0,5/5/0/{0}/{}/0,6/2/1/{0}/{}/0}{0/1,1/2,2/3,3/4,4/5,2/6}{}{}{}\end{tabular} & \jepGraph{0/0/0/{2}/{}/0,1/1/0/{3}/{}/0,2/2/0/{4}/{}/0,3/3/0/{3}/{}/0,4/4/0/{2}/{}/0,5/5/0/{1}/{}/0,6/2/1/{2}/{}/0}{0/1,1/2,2/3,3/4,4/5,2/6}{}{}{}\rule[-12pt]{0pt}{36pt}\\
\cline{3-4}
 & & \begin{tabular}{r@{\;}c}$\text{if }r=2:$ & \jepGraph{0/0/0/{0}/{}/0,1/1/0/{0}/{}/0,2/2/0/{0}/{}/0,3/3/0/{0}/{}/0,4/4/0/{1}/{}/0,5/5/0/{0}/{}/0,6/2/1/{0}/{}/0}{0/1,1/2,2/3,3/4,4/5,2/6}{}{}{}\end{tabular} & \jepGraph{0/0/0/{0}/{}/0,1/1/0/{1}/{}/0,2/2/0/{2}/{}/0,3/3/0/{2}/{}/0,4/4/0/{2}/{}/0,5/5/0/{1}/{}/0,6/2/1/{1}/{}/0}{0/1,1/2,2/3,3/4,4/5,2/6}{}{}{}\rule[-12pt]{0pt}{36pt}\\
\cline{3-4}
 & & \begin{tabular}{r@{\;}c}$\text{if }r=3:$ & \jepGraph{0/0/0/{0}/{}/0,1/1/0/{0}/{}/0,2/2/0/{0}/{}/0,3/3/0/{0}/{}/0,4/4/0/{0}/{}/0,5/5/0/{2}/{}/0,6/2/1/{0}/{}/0}{0/1,1/2,2/3,3/4,4/5,2/6}{}{}{}\end{tabular} & \jepGraph{0/0/0/{0}/{}/0,1/1/0/{0}/{}/0,2/2/0/{0}/{}/0,3/3/0/{0}/{}/0,4/4/0/{0}/{}/0,5/5/0/{1}/{}/0,6/2/1/{0}/{}/0}{0/1,1/2,2/3,3/4,4/5,2/6}{}{}{}\rule[-12pt]{0pt}{36pt}\\
\hline
\end{tabular}
\caption{\small Dominant orthogonal sequences $(\alpha_1,\ldots,\alpha_r)$ for $1\leqslant r\leqslant r_M$}
\label{jep93table1}
\begin{minipage}{.92\textwidth}\footnotesize
(We do not indicate all the roots $\alpha_i$, but rather the sum of the corresponding coroots $\alpha_1^\jepPubCoroot+\cdots+\alpha_r^\jepPubCoroot$, by indicating on a Dynkin diagram the values $\langle\alpha_1^\jepPubCoroot+\cdots+\alpha_r^\jepPubCoroot,\beta\rangle$ for all simple roots $\beta$. In the last column, we express the root $\alpha_r$ which is also the smallest root $\alpha$ such that $\langle\alpha_1^\jepPubCoroot+\cdots+\alpha_r^\jepPubCoroot,\alpha\rangle=2$.)
\end{minipage}
\end{table}
\clearpage

\begin{proof}
Kostant \cite[Lem.~1.6]{jep93Kos12} proved the strong orthogonality. For the convenience of the reader, we recall his arguments in our context. The root $\alpha_i$ is the highest root of $R(\Pi_i)$, and $\Pi_i$ contains the long root $\zeta$, so $\alpha_i$ is long. By construction, $\langle\alpha_i^\jepPubCoroot,\alpha_j\rangle=0$ if $j>i$, so $\alpha_i$ and $\alpha_j$ are orthogonal. Since they are both long, they are strongly orthogonal, otherwise $\alpha_i\pm\alpha_j$ would be longer than $\alpha_i$. This proves~(1).

For~(2), by Fact~2.2, it is enough to prove that $(\alpha_1,\ldots,\alpha_q)$ is a maximal sequence of orthogonal roots (see also \cite[Th.~1.8]{jep93Kos12}). Let $\alpha\in R$ be such that $\langle\alpha_i^\jepPubCoroot,\alpha\rangle=0$ holds for all $i$. We can and will assume that $\alpha>0$. Let $i$ be the greatest integer such that $\alpha\in R(\Pi_i)$. By maximality of $i$ there exists a simple root $\beta$ in $\operatorname{supp}(\alpha)\cap\Sigma_i$. Since $\alpha_i$ is dominant on $R(\Pi_i)$, we have $\langle\alpha_i^\jepPubCoroot,\alpha\rangle\geqslant\langle\alpha_i^\jepPubCoroot,\beta\rangle>0$, a contradiction to the assumption made on $\alpha$.

For the third point, let $i\leqslant r_M$ and let $\beta\in\Pi$.
\begin{itemize}
\item If $\beta\in\Pi_i$, by construction, $\langle\alpha_j^\jepPubCoroot,\beta\rangle=0$ for $j<i$. Since $\alpha_i$ is the highest root of $\Pi_i$, $\langle\alpha_i^\jepPubCoroot,\beta\rangle\geqslant0$. Thus $\langle\alpha_1^\jepPubCoroot+\cdots+\alpha_i^\jepPubCoroot,\beta\rangle=\langle\alpha_i^\jepPubCoroot,\beta\rangle\geqslant0$
\item If $\beta\in\Sigma_{i-1}$, then $\langle\alpha_{i-1}^\jepPubCoroot,\beta\rangle\geqslant1$ and $\langle\alpha_j^\jepPubCoroot,\beta\rangle=0$ for $j<i-1$. Since $\alpha_i$ is long, $\langle\alpha_i^\jepPubCoroot,\beta\rangle\geqslant-1$. Thus, $\langle\alpha_1^\jepPubCoroot+\cdots+\alpha_i^\jepPubCoroot,\beta\rangle\geqslant0$
\item If $\beta\in\Sigma_j$ for $j<i-1$, then, by construction of $\Pi_i$, $\langle\alpha_i^\jepPubCoroot,\beta\rangle=0$. By induction on $i$, $\langle\alpha_1^\jepPubCoroot+\cdots+\alpha_{i-1}^\jepPubCoroot,\beta\rangle\geqslant0$, so $\langle\alpha_1^\jepPubCoroot+\cdots+\alpha_i^\jepPubCoroot,\beta\rangle\geqslant0$
\end{itemize}
For~(4), by Lemma~2.8 and since $\alpha_i$ is long, we have $\langle\varpi,\alpha_i^\jepPubCoroot\rangle=1$. Thus $s_{\alpha_i}(\varpi)=\varpi-\langle\varpi,\alpha_i^\jepPubCoroot\rangle\alpha_i=\varpi-\alpha_i$. Therefore $s_{\alpha_1}\cdots s_{\alpha_{r_M}}(\varpi)=\varpi-\alpha_1-\cdots-\alpha_{r_M}$ is a weight of $\jepPubBbb{E}$. We prove that it is a lowest weight. Let $\beta\in\Pi$. If $\beta\neq\zeta$, then
\[
\langle\varpi-\alpha_1-\cdots-\alpha_{r_M},\beta^\jepPubCoroot\rangle=\langle-\alpha_1-\cdots-\alpha_{r_M},\beta^\jepPubCoroot\rangle,
\]
and this is non-positive by~(3) and because all the roots $\alpha_i$ have the same length. If $\beta=\zeta$, then we compute that $\langle\varpi-\alpha_1-\cdots-\alpha_{r_M},\zeta^\jepPubCoroot\rangle=1-\langle\alpha_{r_M},\zeta^\jepPubCoroot\rangle\leqslant0$ since $\zeta\in\Sigma_{r_M}$.
\end{proof}
We make the following observations.
\begin{proposition}
Let $1\leqslant r\leqslant r_M$, $(\alpha_1,\ldots,\alpha_r)$ be the corresponding dominant orthogonal sequence, and $h_r=\alpha_1^\jepPubCoroot+\cdots+\alpha_r^\jepPubCoroot$. Then:
\begin{enumerate}[label=(\arabic*)]
\item $h_r$ is dominant: for any positive root $\alpha$, we have $\langle\alpha,h_r\rangle\geqslant0$
\item For any root $\alpha$, we have $\langle\alpha,h_r\rangle\leqslant2$
\item $\alpha_r$ is the smallest root $\alpha$ such that $\langle\alpha,h_r\rangle=2$
\item If a root $\alpha$ satisfies $\langle\alpha,h_r\rangle=2$, then $\langle\alpha,z\rangle=2$
\item We have $\langle\varpi,h_r\rangle=r$
\end{enumerate}
\end{proposition}
\begin{proof}
Recall Table~1. The first point has been proved in Proposition~2.19~(3). Let $\Theta$ be the highest root of the root system of $G$, which can be found for example in \cite{jep93Bou68}. Since $h_r$ is dominant, the second item follows from the fact that $\langle\Theta,h_r\rangle=2$ in all cases.

For the third item, we have indicated the root $\alpha_r$ in the last column of the table. It readily follows that it is the smallest root $\alpha$ such that $\langle\alpha,h_r\rangle=2$. Since it has coefficient $1$ on $\zeta$, we have $\langle\alpha_r,z\rangle=2$. Any root $\alpha$ as in the fourth item will be bigger than $\alpha_r$, so the claim follows.

To prove that $\langle\varpi,h_r\rangle=r$, recall that if $(\alpha_1,\ldots,\alpha_{r_M})$ is the maximal dominant orthogonal sequence, the weight $\varpi-\sum_{1}^{r_M}\alpha_i$ is the lowest weight $\mu_0$ by Proposition~2.19~(4). Thus, $\langle\varpi-\mu_0,h_r\rangle=\langle\alpha_1+\cdots+\alpha_{r_M},\alpha_1^\jepPubCoroot+\cdots+\alpha_r^\jepPubCoroot\rangle=2r$, since $\langle\alpha_i,\alpha_j^\jepPubCoroot\rangle=2\delta_{i,j}$ by orthogonality of the roots $\alpha_i$. Moreover, $h_r$ is part of some $\mathfrak{sl}_2$-triple $(x,h_r,y)$, so $\langle\mu_0,h_r\rangle=-\langle\varpi,h_r\rangle$, by the representation theory of $\mathfrak{sl}_2$. This proves that $\langle\varpi,h_r\rangle=r$.
\end{proof}

\subsection{The self-dual Higgs submodule associated with a nilpotent element}
We explain in this section that an element $y\in\mathfrak m^-$ defines a special graded subspace in $\jepPubBbb{E}$. This algebraic construction will be used in Section~4.3 to produce the leafwise Higgs subsheaves needed to prove the Milnor-Wood inequality.
\subsubsection{$K$-orbits in $\mathfrak m^-$: ranks of nilpotent elements}
We begin by describing the $K$-orbits in $\mathfrak m^-$, introducing the notion of rank, and then we choose nice representatives in the orbits to simplify our work. Let us recall the following result, which is proved e.g. in \cite{jep93HC56}, \cite[Ch.~VIII, \S7]{jep93Hel01}, or \cite{jep93Wol72}.
\begin{proposition}
The orbits of the complex group $K$ in $\mathfrak m^-$ are parametrized by integers $r\in\{0,\ldots,r_M\}$. A representative of each orbit is $y_{\alpha_1}+\cdots+y_{\alpha_r}$, where $(\alpha_1,\ldots,\alpha_{r_M})$ is the maximal dominant orthogonal sequence.
\end{proposition}
\begin{definition}
Let $y\in\mathfrak m^-$. The rank $r(y)$ of $y$ is the integer $r$ such that $y$ is in the $K$-orbit of $y_{\alpha_1}+\cdots+y_{\alpha_r}$.
\end{definition}
\begin{example}
In case $\mathfrak g$ has type $A$, $\mathfrak m^-$ identifies with a space of matrices, and the rank as defined above of an element in $\mathfrak m^-$ coincides with its rank as a matrix.
\end{example}
The following proposition is a characterization of the rank $r(y)$ of $y\in\mathfrak m^-$ in terms of its action on $\jepPubBbb{E}$. If $y\in\mathfrak g$, we will denote by $\jepPubBmi{y}\in\operatorname{End}(\jepPubBbb{E})$ the image of $y$ by the representation of $\mathfrak g$ on $\jepPubBbb{E}$.
\begin{proposition}
Let $y=y_{\alpha_1}+\cdots+y_{\alpha_{r(y)}}$ with $(\alpha_1,\ldots,\alpha_{r(y)})$ a dominant orthogonal sequence. We have $\jepPubBmi{y}^{r(y)+1}=0$ and $\jepPubBmi{y}^{r(y)}(\jepPubBbb{E}_0)=\jepPubBbb{E}_{\varpi-\alpha_1-\cdots-\alpha_{r(y)}}$, so in particular $\jepPubBmi{y}^{r(y)}(\jepPubBbb{E}_0)\neq\{0\}$. In other words, the rank of $y$ is the order of nilpotency of $\jepPubBmi{y}$.
\end{proposition}
\begin{proof}
We have $y_{\alpha_i}\in\mathfrak g_{-\alpha_i}$. For any pair $(\chi,\delta)$ with $\chi\in X(\jepPubBbb{E})$ and $\delta$ a long root, $\chi-2\delta$ cannot be a weight of $\jepPubBbb{E}$ by Lemma~2.9. We deduce from Proposition~2.19~(1) that $\jepPubBmi{y}_{\alpha_i}^2=0$. On the other hand, for any $i,j$ with $i\neq j$, the maps $\jepPubBmi{y}_{\alpha_i}$ and $\jepPubBmi{y}_{\alpha_j}$ commute since $\alpha_i+\alpha_j$ is not a root by Proposition~2.19~(1) also. Thus we get
\[
\jepPubBmi{y}^k=k!\sum\jepPubBmi{y}_{\alpha_{j_1}}\circ\cdots\circ\jepPubBmi{y}_{\alpha_{j_k}},
\]
where the sum is over the increasing sequences $1\leqslant j_1<\cdots<j_k\leqslant r(y)$.

In particular $\jepPubBmi{y}^{r(y)+1}=0$, and $\jepPubBmi{y}^{r(y)}=r(y)!\jepPubBmi{y}_{\alpha_1}\circ\cdots\circ\jepPubBmi{y}_{\alpha_{r(y)}}$. By Lemma~2.9, we have $\jepPubBmi{y}_{\alpha_1}\circ\cdots\circ\jepPubBmi{y}_{\alpha_{r(y)}}\cdot\jepPubBbb{E}_\varpi=\jepPubBbb{E}_{\varpi-\alpha_1-\cdots-\alpha_{r(y)}}$, so the proposition is proved.
\end{proof}
For the remaining of this section, we let $1\leqslant r\leqslant r_M$ be an integer and $y=y^r:=y_{\alpha_1}+\cdots+y_{\alpha_r}\in\mathfrak m^-$, where $(\alpha_1,\ldots,\alpha_r)$ is the dominant orthogonal sequence of length $r$ obtained by the algorithm of Definition~2.17. By Proposition~2.21, there is no loss of generality in considering this particular element $y$. We denote by $h_r$ the element $\alpha_1^\jepPubCoroot+\cdots+\alpha_r^\jepPubCoroot$ spelled out in Table~1. Recall that $\jepPubBmi{y}$ is $y$ acting on $\jepPubBbb{E}$ via the cominuscule representation of $\mathfrak g$, and that $r$ is the order of nilpotency of $\jepPubBmi{y}$.
\subsubsection{Definition of $\jepPubBbb{V}^r$ and the Higgs property}
It is a very well-known idea that the nilpotent element $\jepPubBmi{y}$ of $\operatorname{End}(\jepPubBbb{E})$ uniquely defines an increasing filtration $\jepPubBbb{W}^r_\jepPubBullet$ of $\jepPubBbb{E}$, called the \emph{Deligne weight filtration} of $\jepPubBmi{y}$, or Jacobson-Morozov filtration. This idea appears in \cite[\S1.6.13]{jep93Del80} in the context of Hodge theory and in \cite[\S\S3.2 \& 3.4]{jep93McG02} in the context of nilpotent orbits in Lie algebras.

The defining properties of the filtration $(\jepPubBbb{W}^r_k)_{k\in\mathbb Z}$ are that for all $k\in\mathbb Z$,
\[
\jepPubBmi{y}(\jepPubBbb{W}^r_k)\subset\jepPubBbb{W}^r_{k-2}
\]
and for all $k\geqslant1$,
\begin{equation}\tag{2.1}
\jepPubBmi{y}^k:\jepPubBbb{W}^r_k/\jepPubBbb{W}^r_{k-1}\longrightarrow\jepPubBbb{W}^r_{-k}/\jepPubBbb{W}^r_{-k-1}\quad\text{is an isomorphism.}
\end{equation}
The filtration $(\jepPubBbb{W}^r_k)_{k\in\mathbb Z}$ can be defined as follows (see e.g. \cite{jep93SZ85}):
\begin{equation}\tag{2.2}
\jepPubBbb{W}^r_k=\sum_{\substack{\ell\geqslant0\\k+\ell+1\geqslant0}}\operatorname{Ker}\jepPubBmi{y}^{k+\ell+1}\cap\operatorname{Im}\jepPubBmi{y}^{\ell}=\sum_{\substack{i-j=k\\i\geqslant0,\ j\geqslant0}}\operatorname{Ker}\jepPubBmi{y}^{i+1}\cap\operatorname{Im}\jepPubBmi{y}^{j}.
\end{equation}
Since $r$ is the rank of nilpotency of $\jepPubBmi{y}$ acting on $\jepPubBbb{E}$, we have $\jepPubBbb{W}^r_k=\{0\}$ for $k<-r$ and $\jepPubBbb{W}^r_k=\jepPubBbb{E}$ for $k\geqslant r$.

The elements $y$ and $h_r$ fit in an $\mathfrak{sl}_2$-triple $(x,h_r,y)$ for some $x\in\mathfrak m^+$, and the subspaces $\jepPubBbb{W}^r_k$ can be alternatively defined by means of this triple:
\begin{equation}\tag{2.3}
\jepPubBbb{W}^r_k=\jepPubOplus_{\chi\in X(\jepPubBbb{E}):\langle\chi,h_r\rangle\leqslant k}\jepPubBbb{E}_\chi.
\end{equation}
We combine the gradation $\jepPubBbb{E}=\jepPubOplus_{i=0}^{r_M}\jepPubBbb{E}_i$ and the weight filtration $(\jepPubBbb{W}^r_k)_{k\in\mathbb Z}$ to define the subspaces of $\jepPubBbb{E}$ that will be the center of our attention for the remaining part of Section~2.
\begin{notation}
Let $\jepPubBbb{V}^r_i=\jepPubBbb{E}_i\cap\jepPubBbb{W}^r_{r-2i}$ for $0\leqslant i\leqslant r_M$, and let $\jepPubBbb{V}^r=\jepPubOplus_i\jepPubBbb{V}^r_i\subset\jepPubBbb{E}$.
\end{notation}
Observe that $\jepPubBbb{V}^r_0=\jepPubBbb{E}_0$ and that $\jepPubBbb{V}^r_i\neq\{0\}$ only if $i\leqslant r$. The subspaces $\jepPubBbb{V}^r_i$ can also be described as sums of weight spaces $\jepPubBbb{E}_\chi$.
\begin{proposition}
For $\chi\in X(\jepPubBbb{E}_i)$, we have $\langle\chi,h_r\rangle\geqslant r-2i$. Hence, for all $i$ such that $0\leqslant i\leqslant r$,
\[
\jepPubBbb{V}^r_i=\jepPubOplus_{\substack{\chi\in X(\jepPubBbb{E}_i):\\\langle\chi,h_r\rangle=r-2i}}\jepPubBbb{E}_\chi=\jepPubOplus_{\substack{\chi\in X(\jepPubBbb{E}):\\\langle\chi,z\rangle=z_{\mathtt{max}}-2i,\\\langle\chi,h_r\rangle=r-2i}}\jepPubBbb{E}_\chi.
\]
\end{proposition}
\begin{proof}
Proposition~2.15 and the fact that, for any root $\alpha$, we have $\langle\alpha,h_r\rangle\leqslant2$ (see Proposition~2.20) imply the first assertion. It is plain from the description of the subspaces $\jepPubBbb{W}^r_k$ given in~(2.3) that the second assertion follows from the first.
\end{proof}
The following lemma is reminiscent of Equation~(2.1). We will prove a stronger statement in Proposition~2.40. More Hodge-theoretic considerations will be made in Section~2.6.
\begin{lemma}
The linear map $\jepPubBmi{y}^{r-2i}$ is an isomorphism between $\jepPubBbb{V}^r_i$ and $\jepPubBbb{V}^r_{r-i}$. In particular, $\dim\jepPubBbb{V}^r_{r-i}=\dim\jepPubBbb{V}^r_i$.
\end{lemma}
\begin{proof}
We know that
\[
\jepPubBbb{V}^r_i=\jepPubBbb{E}_i\cap\jepPubBbb{W}^r_{r-2i},\qquad\jepPubBbb{V}^r_{r-i}=\jepPubBbb{E}_{r-i}\cap\jepPubBbb{W}^r_{2i-r},
\]
and that for $i$ satisfying $0\leqslant i\leqslant r/2$, $\jepPubBmi{y}^{r-2i}$ is an isomorphism between $\jepPubBbb{W}^r_{r-2i}/\jepPubBbb{W}^r_{r-2i-1}$ and $\jepPubBbb{W}^r_{2i-r}/\jepPubBbb{W}^r_{2i-r-1}$. By Proposition~2.15, $\jepPubBmi{y}^{r-2i}$ maps $\jepPubBbb{E}_i$ to $\jepPubBbb{E}_{r-i}$. Since we have $\jepPubBbb{E}_k\cap\jepPubBbb{W}^r_{r-2k-1}=\{0\}$ for all $k$ by Proposition~2.26 and Equation~(2.3), we get that $\jepPubBmi{y}^{r-2i}$ is an isomorphism between $\jepPubBbb{V}^r_i$ and $\jepPubBbb{V}^r_{r-i}$.
\end{proof}
\begin{definition}
A subspace $\jepPubBbb{S}$ of $\jepPubBbb{E}$ is a \emph{$y$-Higgs subspace} if $y\cdot\jepPubBbb{S}\subset\jepPubBbb{S}$ and $\mathfrak m^+\cdot\jepPubBbb{S}\subset\jepPubBbb{S}$.
\end{definition}
The following algebraic fact is at the heart of the construction of leafwise Higgs subsheaves in Section~4.
\begin{proposition}
The subspace $\jepPubBbb{V}^r$ is a $y$-Higgs subspace of $\jepPubBbb{E}$. More precisely, the following inclusions hold:
\begin{itemize}
\item $y\cdot\jepPubBbb{V}^r_i\subset\jepPubBbb{V}^r_{i+1}$
\item $\mathfrak m^+\cdot\jepPubBbb{V}^r_i\subset\jepPubBbb{V}^r_{i-1}$
\end{itemize}
\end{proposition}
\begin{proof}
Using the description of $\jepPubBbb{V}^r_i$ given in Proposition~2.26, the first inclusion holds because $y\in\jepPubOplus_i\mathfrak g_{-\alpha_i}$ and each $\alpha_i$ satisfies $\langle-\alpha_i,z\rangle=\langle-\alpha_i,h_r\rangle=-2$. The second one also holds since for $\alpha$ a root of $\mathfrak m^+$, we have $\langle\alpha,z\rangle=2$ (and $\langle\alpha,h_r\rangle\leqslant2$ by Proposition~2.20).
\end{proof}
\subsubsection{The module structure of $\jepPubBbb{V}^r$}
We now show that the subspace $\jepPubBbb{V}^r$ of $\jepPubBbb{E}$ has the structure of an irreducible $L_r$-module, for $L_r$ a complex reductive group determined by $y$, whereas the $\jepPubBbb{V}^r_i$ are irreducible $H_r$-modules, where $H_r=L_r\cap K$.

We begin by defining the relevant subalgebras of $\mathfrak g$.
\begin{notation}
Let $\mathfrak q_r\subset\mathfrak g$ be the parabolic subalgebra defined by the coweight $z-h_r$:
\[
\mathfrak q_r=\mathfrak t\oplus\jepPubOplus_{\alpha:\langle\alpha,h_r\rangle\leqslant\langle\alpha,z\rangle}\mathfrak g_\alpha.
\]
A Levi factor of $\mathfrak q_r$ is
\[
\mathfrak l_r=\mathfrak t\oplus\jepPubOplus_{\alpha:\langle\alpha,h_r\rangle=\langle\alpha,z\rangle}\mathfrak g_\alpha.
\]
Let $\mathfrak l_r^\pm\subset\mathfrak l_r$ be the nilpotent subalgebras of $\mathfrak l_r$ defined by
\[
\mathfrak l_r^\pm=\jepPubOplus_{\alpha:\langle\alpha,h_r\rangle=\langle\alpha,z\rangle=\pm2}\mathfrak g_\alpha.
\]
The intersection $\mathfrak k\cap\mathfrak q_r$ is a parabolic subalgebra of $\mathfrak k$ and a Levi factor of it is
\[
\mathfrak h_r=\mathfrak t\oplus\jepPubOplus_{\alpha:\langle\alpha,h_r\rangle=\langle\alpha,z\rangle=0}\mathfrak g_\alpha.
\]
We denote by $Q_r$, $L_r$ and $H_r$, the connected subgroups of $G$ with Lie algebra $\mathfrak q_r$, $\mathfrak l_r$ and $\mathfrak h_r$ respectively. Such groups exist because $\mathfrak q_r$ is a parabolic subalgebra, and $\mathfrak h_r$ and $\mathfrak l_r$ are Levi subalgebras. We observe that $H_r$ is a Levi subgroup of $K\cap Q_r$: in fact, we have $K\cap Q_r=H_r\cdot R(K\cap Q_r)$, where $R(K\cap Q_r)$ denotes the unipotent radical of $K\cap Q_r$. Since Levi subgroups correspond to Levi subalgebras, $L_r$ is a Levi factor of $Q_r$.
\end{notation}
For the convenience of the reader, the conditions that define the different Lie subalgebras of $\mathfrak g$ we are considering are displayed in Table~2.
\begin{table}[H]\centering
\begin{tabular}{lcccccc}\hline
Lie algebra & $\mathfrak k$ & $\mathfrak m^\pm$ & $\mathfrak q_r$ & $\mathfrak l_r$ & $\mathfrak l_r^\pm$ & $\mathfrak h_r$\\
Condition & $z=0$ & $z=\pm2$ & $h_r\leqslant z$ & $h_r=z$ & $h_r=z=\pm2$ & $h_r=z=0$\\\hline
\end{tabular}
\caption{\small Lie subalgebras under consideration (We abbreviate the condition $\langle\alpha,z\rangle=0$, resp. $\langle\alpha,h_r\rangle=0$, on a root $\alpha$ as $z=0$, resp. $h_r=0$.)}
\end{table}
We have the following lemmas concerning the $\jepPubBbb{V}^r_i$’s.
\begin{lemma}
We have $\jepPubBbb{V}^r_{i+1}=\mathfrak l_r^-\cdot\jepPubBbb{V}^r_i$ and $\jepPubBbb{V}^r_{i-1}=\mathfrak l_r^+\cdot\jepPubBbb{V}^r_i$.
\end{lemma}
\begin{proof}
Let us denote by $X(\jepPubBbb{E}_i)$ the set of weights of $\jepPubBbb{E}_i$. We know by Proposition~2.15 that $\jepPubBbb{E}_{i+1}=\mathfrak m^-\cdot\jepPubBbb{E}_i$. Thus,
\[
\jepPubBbb{E}_{i+1}=\biggl(\jepPubOplus_{\alpha:\langle\alpha,z\rangle=-2}\mathfrak g_\alpha\biggr)\cdot\biggl(\jepPubOplus_{\chi\in X(\jepPubBbb{E}_i)}\jepPubBbb{E}_\chi\biggr)=\jepPubOplus_{\substack{\chi\in X(\jepPubBbb{E}_i)\\\alpha\in R(\mathfrak m^-)}}\jepPubBbb{E}_{\chi+\alpha}.
\]
Given $\chi\in X(\jepPubBbb{E}_i)$ and $\alpha\in R(\mathfrak m^-)$, we can make two observations:
\begin{itemize}
\item If $\langle\chi,h_r\rangle>r-2i$ then we have $\langle\chi+\alpha,h_r\rangle>r-2(i+1)$ and thus $\jepPubBbb{E}_{\chi+\alpha}\not\subset\jepPubBbb{V}^r_{i+1}$
\item If $\langle\alpha,h_r\rangle>-2$ then $\langle\chi+\alpha,h_r\rangle>r-2(i+1)$ and thus $\jepPubBbb{E}_{\chi+\alpha}\not\subset\jepPubBbb{V}^r_{i+1}$
\end{itemize}
The first equality of the lemma follows from these observations. The proof of the second equality is similar.
\end{proof}
Note that $\mathfrak l_r^-$ is an $\mathfrak h_r$-module. More precisely, we have
\begin{lemma}
The modules $\jepPubBbb{V}^r_i$ are irreducible $\mathfrak h_r$-modules. The Lie algebra $\mathfrak l_r^-$ is an irreducible $\mathfrak h_r$-module.
\end{lemma}
\begin{proof}
Denote by $\mathfrak k^+$, resp. $\mathfrak h_r^+$, the sum of the root spaces for positive roots in $\mathfrak k$, resp. $\mathfrak h_r$. Let $i$ be fixed and such that $\jepPubBbb{V}^r_i\neq\{0\}$. By Proposition~2.15, $\jepPubBbb{E}_i$ is an irreducible $\mathfrak k$-module. Let $\mu_i$ be the lowest weight of $\jepPubBbb{E}_i$. We have $\jepPubBbb{E}_{\mu_i}\subset\jepPubBbb{V}^r_i$. Since $\jepPubBbb{E}_i$ is irreducible, we have $\jepPubBbb{E}_i=U(\mathfrak k^+)\cdot\jepPubBbb{E}_{\mu_i}$ (here $U$ denotes the universal enveloping algebra). Now, arguing as in the proof of the Lemma~2.31, we see that this implies that $\jepPubBbb{V}^r_i=U(\mathfrak h_r^+)\cdot\jepPubBbb{E}_{\mu_i}$. Now, $\jepPubBbb{E}_{\mu_i}$ has dimension 1, hence $\jepPubBbb{V}^r_i$ is indecomposable. Since $\mathfrak h_r$ is reductive, this proves that $\jepPubBbb{V}^r_i$ is irreducible. Now, by Lemma~2.31 again, we have $\jepPubBbb{V}^r_1=\mathfrak l_r^-\cdot\jepPubBbb{V}^r_0\simeq\mathfrak l_r^-\cdot\jepPubBbb{E}_0$: thus $\mathfrak l_r^-$ is also irreducible.
\end{proof}
Concerning the submodule $\jepPubBbb{V}^r$, we have
\begin{proposition}
\begin{enumerate}[label=(\arabic*)]
\item The subspace $\jepPubBbb{V}^r\subset\jepPubBbb{E}$ is an irreducible $Q_r$-module
\item As a consequence, the nilpotent radical of $Q_r$ acts trivially on $\jepPubBbb{V}^r$
\item $\jepPubBbb{V}^r_i$ is a $K\cap Q_r$-module
\item The subgroup of elements in $G$ which preserve $\jepPubBbb{V}^r$ is exactly $Q_r$
\end{enumerate}
\end{proposition}
\begin{proof}
In characteristic $0$, given a representation of a connected group, a subspace of this representation is stable under this group if and only if it is stable for the induced action of the Lie algebra. Thus, to prove the proposition, it is enough to deal with the corresponding Lie algebras.

Let us prove that $\jepPubBbb{V}^r$ is stable under $\mathfrak q_r$. It is clearly stable under $\mathfrak t$. Let $\alpha$ be such that $\mathfrak g_\alpha\subset\mathfrak q_r$, i.e., $\langle\alpha,h_r\rangle\leqslant\langle\alpha,z\rangle$. Let $v\in\jepPubBbb{E}_\chi\subset\jepPubBbb{V}^r_i$, so that we have $\langle\chi,h_r\rangle=r-2i$. For $x\in\mathfrak g_\alpha$, we have $x\cdot v\in\jepPubBbb{E}_{\chi+\alpha}\subset\jepPubBbb{E}_{i-\langle\alpha,z\rangle/2}$. Since $\langle\alpha,h_r\rangle\leqslant\langle\alpha,z\rangle$, we have $\langle\chi+\alpha,h_r\rangle\leqslant r-2i+\langle\alpha,z\rangle$, thus either $x\cdot v=0$ or $\langle\chi+\alpha,h_r\rangle=r-2i+\langle\alpha,z\rangle$, by Proposition~2.26. In the second case, we get $x\cdot v\in\jepPubBbb{V}^r_{i-\langle\alpha,z\rangle/2}$. Combining Lemmas~2.31 and~2.32, we see that $\jepPubBbb{V}^r$ is an irreducible $\mathfrak l_r$-module, hence also an irreducible $\mathfrak q_r$-module.

It follows from Engel’s theorem that there exists a non-zero vector in $\jepPubBbb{V}^r$ which is annihilated by the nilpotent radical of $\mathfrak q_r$. Since this radical is an ideal and $\jepPubBbb{V}^r$ is an irreducible $\mathfrak q_r$-module, this radical acts trivially on $\jepPubBbb{V}^r$.

The fact that $\jepPubBbb{V}^r_i$ is a $K\cap Q_r$-module follows because $\jepPubBbb{V}^r_i=\jepPubBbb{E}_i\cap\jepPubBbb{V}^r$, $\jepPubBbb{E}_i$ is $K$-stable, and $\jepPubBbb{V}^r$ is $Q_r$-stable.

Finally, to prove that the stabilizer of $\jepPubBbb{V}^r$ is exactly $Q_r$, let us denote by $\mathfrak{stab}(\jepPubBbb{V}^r)\subset\mathfrak g$ the Lie subalgebra preserving the subspace $\jepPubBbb{V}^r$ in $\jepPubBbb{E}$. We know by~(1) that $\mathfrak{stab}(\jepPubBbb{V}^r)\supset\mathfrak q_r$. On the other hand, let $\alpha$ be a root such that $\mathfrak g_\alpha\cdot\jepPubBbb{V}^r\subset\jepPubBbb{V}^r$. Let us assume as a first case that $\langle\alpha,z\rangle=-2$, i.e., $\mathfrak g_\alpha\subset\mathfrak m^-$. Since, by Proposition~2.15(d), the action of $\mathfrak g$ on $\jepPubBbb{E}$ induces an isomorphism $\jepPubBbb{E}_1\simeq\jepPubBbb{E}_\varpi\otimes\mathfrak m^-$, we have $\mathfrak g_\alpha\cdot\jepPubBbb{V}^r_0=\mathfrak g_\alpha\cdot\jepPubBbb{E}_\varpi=\jepPubBbb{E}_{\varpi+\alpha}$. Then $\jepPubBbb{E}_{\varpi+\alpha}\subset\jepPubBbb{V}^r\cap\jepPubBbb{E}_1=\jepPubBbb{V}^r_1$, so $\langle\alpha,h_r\rangle=-2$, and $\mathfrak g_a\subset\mathfrak q_r$. Assume now that $\langle\alpha,z\rangle=0$, i.e., $\mathfrak g_\alpha\subset\mathfrak k$, and by contradiction that $\langle\alpha,h_r\rangle>0$. Proposition~2.20 then implies that $\langle\alpha,h_r\rangle=1$. For any integer $1\leqslant i\leqslant r$, we cannot have $\langle\alpha,\alpha_i^\jepPubCoroot\rangle=2$ because this would imply $\alpha=\alpha_i$ and $\langle\alpha,h_r\rangle=2$. Let $i$ be such that $\langle\alpha,\alpha_i^\jepPubCoroot\rangle>0$: then $\langle\alpha,\alpha_i^\jepPubCoroot\rangle=1$. Therefore, $\alpha-\alpha_i$ is a root. Since $\mathfrak g_\alpha\subset\mathfrak k$, $\mathfrak g_\alpha\cdot\jepPubBbb{E}_i\subset\jepPubBbb{E}_i$ and hence $\mathfrak g_\alpha\cdot\jepPubBbb{V}^r_i\subset\jepPubBbb{V}^r_i$. Furthermore $\mathfrak g_\alpha(\jepPubBbb{V}^r_0)=0$ since $\jepPubBbb{V}^r_0=\jepPubBbb{E}_\varpi$. It follows that
\[
\mathfrak g_{\alpha-\alpha_i}\cdot\jepPubBbb{E}_\varpi=[\mathfrak g_\alpha,\mathfrak g_{-\alpha_i}]\cdot\jepPubBbb{E}_\varpi=\mathfrak g_\alpha\cdot(\mathfrak g_{-\alpha_i}\cdot\jepPubBbb{E}_\varpi).
\]
Again by Proposition~2.15(d), $\mathfrak g_{\alpha-\alpha_i}\cdot\jepPubBbb{E}_\varpi=\jepPubBbb{E}_{\varpi+\alpha-\alpha_i}$ and $\mathfrak g_{-\alpha_i}\cdot\jepPubBbb{E}_\varpi=\jepPubBbb{E}_{\varpi-\alpha_i}$. But $\jepPubBbb{E}_{\varpi-\alpha_i}\subset\jepPubBbb{V}^r_1$ whereas $\jepPubBbb{E}_{\varpi+\alpha-\alpha_i}\not\subset\jepPubBbb{V}^r_1$ since we assumed $\langle\alpha,h_r\rangle>0$. This contradicts $\mathfrak g_\alpha\cdot\jepPubBbb{V}^r\subset\jepPubBbb{V}^r$.

Let now $\operatorname{Stab}(\jepPubBbb{V}^r)\subset G$ be the subgroup stabilizing $\jepPubBbb{V}^r$. We have $Q_r\subset\operatorname{Stab}(\jepPubBbb{V}^r)$, thus $\operatorname{Stab}(\jepPubBbb{V}^r)$ is parabolic and therefore connected \cite[Cor.~23.1.B]{jep93Hum75}. It follows that $Q=\operatorname{Stab}(\jepPubBbb{V}^r)$.
\end{proof}

\subsubsection{The self-duality property. Tube type Hermitian symmetric spaces}
By definition, a Hermitian symmetric space $M=G_{\mathbb R}/K_{\mathbb R}$ is of \emph{tube type} if it is biholomorphically equivalent to a domain $F\oplus iC$, where $C\subset F$ is a proper open cone in the real vector space $F$. The simplest example is the upper half-plane, where $F=\mathbb R$ and $C=\mathbb R^*_+$. For irreducible symmetric spaces, this happens when $G_{\mathbb R}$ is isogenous to $\mathrm{SU}(p,p)$, $\mathrm{Spin}(2,m)$, $\mathrm{Sp}(2m,\mathbb R)$, $\mathrm{SO}^{\star}(2m)$ for $m$ even, and $\mathrm E_{7(-25)}$.

We introduce the following definition.
\begin{definition}
A sequence $(\gamma_1,\ldots,\gamma_r)$ of roots will be called \emph{admissible} if all the roots $\gamma_i$ are long and bigger than $\zeta$, and if they are pairwise strongly orthogonal.
\end{definition}
Note that the dominant orthogonal sequences $(\alpha_1,\ldots,\alpha_r)$ are admissible. We can now state:
\begin{proposition}
The following assertions are equivalent:
\begin{enumerate}[label=(\arabic*)]
\item $M$ has tube type
\item For all admissible sequences $\gamma_1,\ldots,\gamma_{r_M}$, we have $\gamma_1+\cdots+\gamma_{r_M}=2\varpi$
\item $\alpha_1+\cdots+\alpha_{r_M}=2\varpi$
\item There exists an admissible sequence $\gamma_1,\ldots,\gamma_{r_M}$ such that $\gamma_1+\cdots+\gamma_{r_M}=2\varpi$
\item $\alpha_1^\jepPubCoroot+\cdots+\alpha_{r_M}^\jepPubCoroot=z$
\item $\jepPubBbb{E}$ is self-dual: $\jepPubBbb{E}\simeq\jepPubBbb{E}^{\star}$ as $G$-modules
\item For all $i$, $\jepPubBbb{E}_i\simeq\jepPubBbb{E}_{r_M-i}^{\star}$ as $K$-modules
\item $\dim\jepPubBbb{E}_{r_M}=1$
\item $\alpha_{r_M}=\zeta$
\end{enumerate}
\end{proposition}
\begin{proof}
Glancing at Table~1, one can readily check that $G_{\mathbb R}/K_{\mathbb R}$ is of tube type if and only if $z=\alpha_1^\jepPubCoroot+\cdots+\alpha_{r_M}^\jepPubCoroot$, and that this occurs also exactly when $\alpha_{r_M}=\zeta$. Hence~(1),~(5) and~(9) are equivalent.

The equality $z=a_1^\jepPubCoroot+\cdots+\alpha_{r_M}^\jepPubCoroot$ is equivalent to $\langle\alpha_1^\jepPubCoroot+\cdots+\alpha_{r_M}^\jepPubCoroot,\beta\rangle=2\delta_{\beta,\zeta}$ for all $\beta\in\Pi$ because $\zeta$ is the only noncompact simple root and it has the same length as each $\alpha_i$ (they are all long). Since all the roots $\alpha_i$ have the same length, for any root $\beta$, the equation $\langle\alpha_1^\jepPubCoroot+\cdots+\alpha_{r_M}^\jepPubCoroot,\beta\rangle=0$ is equivalent to $\langle\alpha_1+\cdots+\alpha_{r_M},\beta^\jepPubCoroot\rangle=0$. It follows that~(5) and~(3) are equivalent.

We now prove that items (2), (3), and (4) are equivalent. In fact, we have the trivial implications $\text{(2)}\Rightarrow\text{(3)}\Rightarrow\text{(4)}$. Moreover, if~(4) holds, then let $\gamma_1,\ldots,\gamma_{r_M}$ be an admissible sequence such that $\gamma_1+\cdots+\gamma_{r_M}=2\varpi$. By \cite[Th.~2.1]{jep93RRS92}, any admissible sequence can be obtained applying to this sequence an element $w$ of the subgroup of the Weyl group that fixes $\varpi$. Condition~(2) follows from the equality $w(\gamma_1)+\cdots+w(\gamma_{r_M})=w(2\varpi)=2\varpi$.

Since the weights of $\jepPubBbb{E}^{\star}$ are the opposite of the weights of $\jepPubBbb{E}$, the highest weight of $\jepPubBbb{E}^{\star}$ is $-\mu_0$. Therefore,~(6) holds if and only if $\mu_0=-\varpi$, which is equivalent to~(3) by Proposition~2.19~(4).

Assuming~(6) and using the decomposition into $K$-modules $\jepPubBbb{E}=\jepPubOplus_{i=0}^{r_M}\jepPubBbb{E}_i$, we deduce that $\jepPubBbb{E}_{r_M-i}\simeq\jepPubBbb{E}_i^{\star}$ as $K$-modules. Conversely, given~(7), we obviously have~(8). Assuming now~(8), we deduce that $\jepPubBbb{E}_{r_M}=\jepPubBbb{E}_{\mu_0}$. Since $\jepPubBbb{E}_{r_M}$ is a $K$-module of dimension $1$ in this case, it follows that $\mu_0$ vanishes on all the roots of $\mathfrak k$, so $\mu_0$ is an integral multiple of $\varpi$. It follows that $\mu_0=-\varpi$ and we get~(3) as above.
\end{proof}
As we just said, the $G$-module $\jepPubBbb{E}$ is not always self-dual. However, the subspace $\jepPubBbb{V}^r$ is always a self-dual module under a subgroup $\overline L_r$ of $G$ that we now define.

Recall that the Lie algebra $\mathfrak l_r$ is by definition generated by the root spaces of the roots $\alpha$ such that $\langle h_r,\alpha\rangle=\langle z,\alpha\rangle$, and that $\mathfrak l_r^\pm$ is the subalgebra generated by the root spaces of the roots $\alpha$ such that $\langle h_r,\alpha\rangle=\langle z,\alpha\rangle=\pm2$. We define $\overline{\mathfrak l}_r$ as the Lie subalgebra of $\mathfrak g$ generated by the spaces $\mathfrak l_r^+$ and $\mathfrak l_r^-$. Two particular cases are easy:
\begin{fact}
If $r=1$, then $\mathfrak{sl}_2\simeq\overline{\mathfrak l}_r\subset\mathfrak g$ is the Lie subalgebra corresponding to the highest root $\Theta$. If $M$ has tube type and $r=r_M$, then $\overline{\mathfrak l}_r=\mathfrak g$.
\end{fact}
\begin{proof}
In case $r=1$, then $h_r=\Theta^\jepPubCoroot$, so $\langle\alpha,h_r\rangle=2$ implies $\alpha=\Theta$. Thus $\mathfrak l_r^+=\mathfrak g_\Theta$ in this case. When $M$ has tube type and $r=r_M$, we have $h_r=z$ (Proposition~2.35~(5)), so $\mathfrak l_r^+=\mathfrak m^+$ in this case.
\end{proof}
We now describe the Lie algebra $\overline{\mathfrak l}_r$ in general. To this end, we introduce a definition.
\begin{definition}
Let $R$ be a (possibly reducible) root system, let $\Pi$ be a basis of $R$, and $\varpi$ be a dominant weight. Given a simple root $\beta$, we say that $\beta$ is \emph{connected to $\varpi$ in $R$} if there is a simple root $\gamma$ in the same connected component as $\beta$ in the Dynkin diagram of $R$ and such that $\langle\varpi,\gamma^\jepPubCoroot\rangle>0$.
\end{definition}
We denote by $\Pi_{h_r=0}$ the set of simple roots $\beta\in\Pi$ such that $\langle\beta,h_r\rangle=0$.
\begin{lemma}
The Lie algebra $\overline{\mathfrak l}_r$ is simple. The lowest root $-\Theta$ together with the simple roots which are connected to $\Theta$ in $\Pi_{h_r=0}$ form a basis of the root system of $\overline{\mathfrak l}_r$.
\end{lemma}
\begin{proof}
Recall that $\zeta$ denotes the only noncompact simple root (see Notation~2.5). Looking at Table~1, we see that the root $\zeta$ cannot be connected to $\Theta$ in $\Pi_{h_r=0}$. Thus, if $\gamma$ is connected to $\Theta$ in $\Pi_{h_r=0}$, it satisfies $\langle\gamma,z\rangle=0$ and, of course, $\langle\gamma,h_r\rangle=0$. Let us assume that $\langle\Theta,\gamma^\jepPubCoroot\rangle>0$. Then $\Theta-\gamma$ is a root and it satisfies $\langle h_r,\Theta-\gamma\rangle=\langle z,\Theta-\gamma\rangle=2$, so $\mathfrak g_{\Theta-\gamma}\subset\overline{\mathfrak l}_r$, and so $\mathfrak g_\gamma\subset\overline{\mathfrak l}_r$. By induction, the root spaces of all simple roots connected to $\Theta$ and their opposite are in $\overline{\mathfrak l}_r$, and, together with $\gamma_{-\Theta}$ and $\gamma_\Theta$, they generate $\overline{\mathfrak l}_r$.

In general, the Dynkin diagram of a Lie algebra has vertices corresponding to the simple roots of the Lie algebra. In our case, given the above description of a basis of $\overline{\mathfrak l}_r$, the set of vertices of its Dynkin diagram is the union of some connected components of $\Pi_{h_r=0}$ and the set $\{-\Theta\}$. Moreover, all the connected components of $\Pi_{h_r=0}$ which occur are connected to $-\Theta$. Thus the Dynkin diagram of $\overline{\mathfrak l}_r$ is connected and $\overline{\mathfrak l}_r$ is simple.
\end{proof}
In Table~3, we indicate in each case which root system for $\overline{\mathfrak l}_r$ is obtained (the cases where $r=1$ or $M$ has tube type and $r=r_M$ are not represented, see Fact~2.36).
\clearpage

\begin{table}[htbp]\centering\normalsize
\setlength{\tabcolsep}{5pt}\renewcommand{\arraystretch}{1.4}
\begin{tabular}{|c|c|c|}\hline
$(G,\varpi)$ & $\overline{\mathfrak l}_r$ & $R(\overline{\mathfrak l}_r)$\\\hline
$(A_{n-1},\varpi_\ell)$ & \jepGraph{0/0/0/{\beta_{r-1}}/{}/0,1/1.55/0/{\beta_2}/{}/0,2/2.55/0/{\beta_1}/{}/0,3/3.55/0/{-\Theta}/{}/1,4/4.55/0/{\beta_n}/{}/0,5/5.55/0/{\beta_{n-1}}/{}/0,6/7.1/0/{\beta_{n-r+1}}/{}/0}{1/2,2/3,3/4,4/5}{0/1,5/6}{}{} & $A_{2r-1}$\rule[-15pt]{0pt}{46pt}\\\hline
$(C_n,\varpi_n)$ & \jepGraph{0/0/0/{\beta_{r-1}}/{}/0,1/1.55/0/{\beta_1}/{}/0,2/2.55/0/{-\Theta}/{}/1}{}{0/1}{1/2/-1}{} & $C_r$\rule[-15pt]{0pt}{46pt}\\\hline
$(D_n,\varpi_n)$ & \jepGraph{0/0/0/{\beta_{2(r-1)}}/{}/0,1/1.55/0/{\beta_3}/{}/0,2/2.55/0/{\beta_2}/{}/0,3/3.05/0.8/{\beta_1}/{}/0,4/3.05/-0.8/{-\Theta}/{}/1}{1/2,2/3,2/4}{0/1}{}{} & $D_{2r-1}$\rule[-15pt]{0pt}{46pt}\\\hline
$(E_6,\varpi_1)$, $r=2$ & \jepGraph{0/0/0/{-\Theta}/{}/1,1/1/0/{\beta_2}/{}/0,2/2/0/{\beta_3}/{}/0,3/2.5/0.8/{\beta_1}/{}/0,4/2.5/-0.8/{\beta_4}/{}/0}{0/1,1/2,2/3,2/4}{}{}{} & $D_5$\rule[-15pt]{0pt}{46pt}\\\hline
$(E_7,\varpi_7)$, $r=2$ & \jepGraph{0/0/0/{-\Theta}/{}/1,1/1/0/{\beta_1}/{}/0,2/2/0/{\beta_3}/{}/0,3/3/0/{\beta_4}/{}/0,4/3.5/0.8/{\beta_5}/{}/0,5/3.5/-0.8/{\beta_2}/{}/0}{0/1,1/2,2/3,3/4,3/5}{}{}{} & $D_6$\rule[-15pt]{0pt}{46pt}\\\hline
\end{tabular}
\caption{\small The Lie algebra $\overline{\mathfrak l}_r$ and its root system (The cases where $r=1$ or $M$ has tube type and $r=r_M$ are not represented.)}
\label{jep93table3}
\end{table}
\clearpage

We consider the root system $R_{\overline{\mathfrak l}_r}$ with its basis $(-\Theta,\beta_1,\ldots,\beta_k)$ given by Lemma~2.38. We have
\begin{lemma}
The simple root $-\Theta$ in $R(\overline{\mathfrak l}_r)$ is cominuscule.
\end{lemma}
In other words, the lemma says that any root in $R(\overline{\mathfrak l}_r)$ has coefficient at most one on the simple root $-\Theta$. Observe that, as a consequence, $-\Theta$ defines a Hermitian noncompact real form $\overline L_{r,\mathbb R}$ of $\overline L_r$. The element $y=y^r$ can be considered as an element of $\mathfrak l_r^-$.
\begin{proof}
Since any root in $R(\overline{\mathfrak l}_r)$ can be obtained from $\Theta$ applying a suitable element of the Weyl group of $\overline{\mathfrak l}_r$, it is enough to show the implication
\[
|n_{-\Theta}(\alpha)|\leqslant1\Longrightarrow|n_{-\Theta}(s_\gamma(\alpha))|\leqslant1,
\]
for any simple root $\gamma$ in $R(\overline{\mathfrak l}_r)$. To prove this, we assume without loss of generality that $n_{-\Theta}(\alpha)\geqslant0$. If $\gamma$ is one of the simple roots $\beta_j$, then $n_{-\Theta}(\alpha)=n_{-\Theta}(s_\gamma(\alpha))$, so the implication clearly holds. If $\gamma=-\Theta$, we consider two cases. If $n_{-\Theta}(\alpha)=1$, then $\alpha$ is a negative root. We have $0\leqslant\langle\alpha,-\Theta^\jepPubCoroot\rangle\leqslant2$, thus $-1\leqslant n_{-\Theta}(s_\gamma(\alpha))\leqslant1$. If $n_{-\Theta}(\alpha)=0$, then $\alpha\neq-\Theta$. Since $\Theta$ is long, it follows that $|\langle\alpha,-\Theta^\jepPubCoroot\rangle|\leqslant1$. Thus $|n_{-\Theta}(s_\gamma(\alpha))|\leqslant1$, as expected.
\end{proof}
As announced, we have the following, which strengthens Lemma~2.27.
\begin{proposition}
The $\overline L_r$-module generated by $\jepPubBbb{E}_\varpi$ is $\jepPubBbb{V}^r$. It is self-dual.
\end{proposition}
\begin{proof}
First, we have $\overline{\mathfrak l}_r\subset\mathfrak l_r$. This follows from the definition of $\overline{\mathfrak l}_r$ and the inclusions $\mathfrak l_r^+\subset\mathfrak l_r,\ \mathfrak l_r^-\subset\mathfrak l_r$.

From the inclusion $\overline{\mathfrak l}_r\subset\mathfrak l_r$ and Proposition~2.33(1), it follows that the $\overline L_r$-submodule generated by $\jepPubBbb{E}_\varpi$ is included in $\jepPubBbb{V}^r$. By Lemma~2.31, $\jepPubBbb{V}^r$ is included in the $\mathfrak l_r^-$-submodule generated by $\jepPubBbb{E}_\varpi$. Since $\mathfrak l_r^-\subset\overline{\mathfrak l}_r$ by definition, we deduce that the $\overline L_r$-module generated by $\jepPubBbb{E}_\varpi$ is exactly $\jepPubBbb{V}^r$.

Now we prove that $\jepPubBbb{V}^r$, as an $\overline L_r$-module, is self-dual. The highest weight of $\jepPubBbb{V}^r$ as an $\overline L_r$-module is the restriction of $\varpi$, which is the fundamental weight $\varpi_{-\Theta}$ corresponding to the simple root $-\Theta$ in the root system of $\overline{\mathfrak l}_r$. The sequence $\alpha_1,\ldots,\alpha_r$ is a sequence of long strongly orthogonal roots. They are all bigger than $-\Theta$. In fact, if $\alpha_i$ was not bigger than $-\Theta$, then it would be a linear combination of the roots $\beta_j$, so it would satisfy $n_\zeta(\alpha_i)=0$, contradicting the definition of the sequence $\alpha_1,\ldots,\alpha_r$.

In view of~(4) in Proposition~2.35, it is enough to prove that $\sum\alpha_i$ and $2\varpi_{-\Theta}$ have the same pairing with the simple roots of $R_{\overline{\mathfrak l}_r}$. For $-\Theta$, we have $\langle\sum\alpha_i^\jepPubCoroot,-\Theta^\jepPubCoroot\rangle=-2$ and $\langle2\varpi_{-\Theta},-\Theta^\jepPubCoroot\rangle=-2$. For $\beta$ a simple root connected to $\Theta$ in $\Pi_{h_r=0}$, we have $\langle\sum\alpha_i,\beta^\jepPubCoroot\rangle=0$ because $\langle\beta,\sum\alpha_i^\jepPubCoroot\rangle=0$. We have $\langle\varpi_{-\Theta},\beta^\jepPubCoroot\rangle=0$. This proves that $\jepPubBbb{V}^r$ is self-dual.
\end{proof}
From the equivalences in Proposition~2.35 and Proposition~2.40, and recalling that $\overline L_{r,\mathbb R}$ is a Hermitian noncompact real form of $\overline L_r$, it follows that the symmetric space $M^r:=\overline L_{r,\mathbb R}/(\overline L_{r,\mathbb R}\cap K_{\mathbb R})$ is a tube type Hermitian symmetric space of rank $r$.
\begin{remark}
This construction is related, but not identical, to the theory of boundary components of bounded symmetric domains (see e.g. \cite{jep93Wol72,jep93PS69,jep93Sat80,jep93AMRT10}). The topological boundary of the symmetric space $M$ inside its compact dual $\check M$ is a union of $r_M$ $G_{\mathbb R}$-orbits, and each $G_{\mathbb R}$-orbit is a union of so-called boundary components, cf. \cite[Def.~3.2 p.~127]{jep93AMRT10}. The components are isomorphic to Hermitian symmetric subspaces $M_\Lambda$ of $M$, where $\Lambda$ is a subset of the set of noncompact strongly orthogonal roots, see \cite[Th.~3.3 p.~127]{jep93AMRT10} for example. If we take our maximal dominant orthogonal sequence $(\alpha_1,\ldots,\alpha_{r_M})$ as the set of strongly orthogonal roots, and $\Lambda=\{\alpha_1,\ldots,\alpha_r\}$ as a subset, boundary components corresponding to this subset are symmetric spaces of rank $r$, as is the subspace $M^r$, but they are not of tube type unless $M$ already is. The Lie algebra of the stabilizer of $M_\Lambda$ in $G_{\mathbb R}$ is $\mathfrak q_\Lambda\cap\mathfrak g_{\mathbb R}$, where $\mathfrak q_\Lambda$ is the parabolic subalgebra of $\mathfrak g$ corresponding to the coweight $h_{r_M}-h_r$, whereas we recall that the parabolic subalgebra $\mathfrak q_r$ we are considering here corresponds to the coweight $z-h_r$. This shows that $\mathfrak q_r=\mathfrak q_\Lambda$ if and only if $M$ is of tube type.
\end{remark}
\clearpage

\begin{example}
We give an example of this construction. We assume that we are in the first row of Table~1, namely that $G$ has type $A_{p+q-1}$ for some positive integers $p\leqslant q$ and that $\varpi=\varpi_p$. The group $G_{\mathbb R}$ is therefore $\operatorname{SU}(p,q)$ acting on $\mathbb C^{p+q}$ preserving a Hermitian form $h$ of signature $(p,q)$ given by the diagonal matrix $\operatorname{diag}(-1,\ldots,-1,1,\ldots,1)$ in the canonical basis $(e_1,\ldots,e_p,e_{p+1},\ldots,e_{p+q})$.

Let $y\in\mathfrak m^-$ be an element of rank $r$. We have a natural decomposition $\mathbb C^{p+q}=A\oplus N\oplus B\oplus I$, with $A$, resp. $N,B,I$ generated by $(e_1,\ldots,e_r)$, resp. $(e_{r+1},\ldots,e_p)$, $(e_{p+1},\ldots,e_{p+q-r})$, $(e_{p+q-r+1},\ldots,e_{p+q})$, so that $y$ is given by the block matrix
\[
\begin{pmatrix}0&0&0&0\\0&0&0&0\\0&0&0&0\\ *&0&0&0\end{pmatrix}.
\]
The element $h_r\in\mathfrak t$ acts on the Lie algebra $\mathfrak{sl}(p+q,\mathbb C)$ of $G$, and we represent the weights of its action as the matrix
\[
\begin{pmatrix}0&1&1&2\\-1&0&0&1\\-1&0&0&1\\-2&-1&-1&0\end{pmatrix},
\]
where the blocks correspond to the decomposition $A\oplus N\oplus B\oplus I$ of $\mathbb C^{p+q}$. Similarly, the weights of the action of the central element $z$ are given as follows
\[
\begin{pmatrix}0&0&2&2\\0&0&2&2\\-2&-2&0&0\\-2&-2&0&0\end{pmatrix}.
\]
Thus the Lie algebra $\mathfrak q_r$, $\mathfrak l_r$ and $\mathfrak h_r$ are respectively
\[
\begin{pmatrix}*&0&*&*\\{*}&*&*&*\\0&0&*&0\\{*}&0&*&*\end{pmatrix},\quad
\begin{pmatrix}*&0&0&*\\0&*&0&0\\0&0&*&0\\{*}&0&0&*\end{pmatrix}
\quad\text{and}\quad
\begin{pmatrix}*&0&0&0\\0&*&0&0\\0&0&*&0\\0&0&0&*\end{pmatrix}.
\]
We see that $Q_r$ is exactly the stabilizer of the flag $(N\subset N\oplus A\oplus I)$, that the subgroup $L_r$ of $G$ corresponding to $\mathfrak l_r$ is $\operatorname S(\operatorname{GL}(N)\times\operatorname{GL}(A\oplus I)\times\operatorname{GL}(B))$ and that $H_r$ is the block diagonal group $\operatorname S(\operatorname{GL}(N)\times\operatorname{GL}(A)\times\operatorname{GL}(I)\times\operatorname{GL}(B))$. Finally the intersections of $G_{\mathbb R}=\operatorname{SU}(p,q)$ with these two latter groups are isomorphic to $\operatorname S(\operatorname U(p-r)\times\operatorname U(r,r)\times\operatorname U(q-r))$ and $\operatorname S(\operatorname U(p-r)\times\operatorname U(r)\times\operatorname U(r)\times\operatorname U(q-r))$ respectively, so that $\overline L_{r,\mathbb R}\simeq\operatorname{SU}(r,r)$.

On the other hand, we have $\jepPubBbb{E}=\wedge^p(N\oplus A\oplus I\oplus B)$ and it is easy to check that $\jepPubBbb{V}^r=\wedge^{p-r}N\otimes\wedge^r(A\oplus I)$, which illustrates that the stabilizer of $\jepPubBbb{V}^r$ is $Q_r$.
\end{example}
\subsection{The slope of the $H_r$-modules $\jepPubBbb{V}^r_i$}
Our results so far allow to compute the slope of the $H_r$-modules $\jepPubBbb{V}^r_i$. We first need two definitions and a lemma. Given an algebraic group $H$, we denote by $X(H)$ its character group considered as a $\mathbb Z$-module with additive law defined by: $(\chi+\psi)(h)=\chi(h)\psi(h)$.
\begin{definition}
Let $H$ be a reductive group and $\jepPubBbb{F}$ be an $H$-module. The slope of $\jepPubBbb{F}$ is the element $\mu_H(\jepPubBbb{F})=\det(\jepPubBbb{F})/\dim(\jepPubBbb{F})$ in $X(H)\otimes_{\mathbb Z}\mathbb Q$. We say that $\jepPubBbb{F}$ is a polystable $H$-module if $\jepPubBbb{F}$ is a direct sum of irreducible $H$-modules of the same slope.
\end{definition}
\begin{lemma}
Let $H$ be a reductive group.
\begin{itemize}
\item Let $\jepPubBbb{F}$ be a polystable $H$-module. For any $H$-submodule $\jepPubBbb{S}$ of $\jepPubBbb{F}$, $\mu_H(\jepPubBbb{S})=\mu_H(\jepPubBbb{F})$ and $\mu_H(\jepPubBbb{F}/\jepPubBbb{S})=\mu_H(\jepPubBbb{F})$.
\item Let $\jepPubBbb{F}$, $\jepPubBbb{F}'$ be polystable $H$-modules. Then $\jepPubBbb{F}\otimes \jepPubBbb{F}'$ is a polystable $H$-module, whose slope is $\mu_H(\jepPubBbb{F})+\mu_H(\jepPubBbb{F}')$.
\item Let $\jepPubBbb{F}$ be a polystable $H$-module. Then $\jepPubBbb{F}^\star$ is a polystable $H$-module, whose slope is $-\mu_H(\jepPubBbb{F})$.
\end{itemize}
\end{lemma}
\begin{proof}
The first assertion is clear because $H$ is reductive and so any submodule is a direct sum of irreducible components. We prove the second assertion. Let $Z$ be the center of $H$. It is known that restriction to $Z$ yields an injection $X(H)\hookrightarrow X(Z)$. In fact, by \cite[Prop.~14.2]{jep93Bor69}, we have $H=Z\cdot(H,H)$, and any character of $H$ is trivial on $(H,H)$.

Let us first assume that $\jepPubBbb{F}$ and $\jepPubBbb{F}'$ are irreducible $H$-modules. By Schur's lemma, there are characters $\chi,\psi$ of $Z$ such that $\forall g\in Z,\ \forall w\in \jepPubBbb{F},\ \forall w'\in \jepPubBbb{F}',\ g\cdot w=\chi(g)w$ and $g\cdot w'=\psi(g)w'$. Therefore, in $X(Z)\otimes\mathbb Q$, we have $\chi=\mu_H(\jepPubBbb{F})\jepPubRestrict_Z$ and $\psi=\mu_H(\jepPubBbb{F}')\jepPubRestrict_Z$. Let now $\jepPubBbb{U}\subset \jepPubBbb{F}\otimes \jepPubBbb{F}'$ be an irreducible component. We have
\[
\forall u\in \jepPubBbb{U},\ \forall g\in Z,\quad g\cdot u=\chi(g)\psi(g)u.
\]
Therefore $\mu_H(\jepPubBbb{U})\jepPubRestrict_Z=\chi+\psi$. Thus, $(\mu_H(\jepPubBbb{F})+\mu_H(\jepPubBbb{F}'))\jepPubRestrict_Z=\mu_H(\jepPubBbb{U})\jepPubRestrict_Z$. Since restriction of characters to $Z$ is injective, we have $\mu_H(\jepPubBbb{U})=\mu_H(\jepPubBbb{F})+\mu_H(\jepPubBbb{F}')$.

Let now $\jepPubBbb{F}$ and $\jepPubBbb{F}'$ be arbitrary $H$-modules, and write the decomposition into irreducible submodules: $\jepPubBbb{F}=\jepPubOplus \jepPubBbb{F}_i$ and $\jepPubBbb{F}'=\jepPubOplus \jepPubBbb{F}'_j$. Let $i,j$ be fixed and let $\jepPubBbb{U}\subset \jepPubBbb{F}_i\otimes \jepPubBbb{F}'_j$ be an irreducible factor. We have proved that $\mu_H(\jepPubBbb{U})=\mu_H(\jepPubBbb{F}_i)+\mu_H(\jepPubBbb{F}'_j)$. Since $\jepPubBbb{F}$ and $\jepPubBbb{F}'$ are polystable, we have $\mu_H(\jepPubBbb{F}_i)=\mu_H(\jepPubBbb{F})$ and $\mu_H(\jepPubBbb{F}'_j)=\mu_H(\jepPubBbb{F}')$. Thus, $\jepPubBbb{U}$ has slope $\mu_H(\jepPubBbb{F})+\mu_H(\jepPubBbb{F}')$ and the lemma is proved.

The third item follows from the identity $\det \jepPubBbb{F}^\star=-\det \jepPubBbb{F}$ in $X(H)$.
\end{proof}
We apply this lemma to the $H_r$-modules $\jepPubBbb{V}^r_i$:
\begin{proposition}
We have $\mu_{H_r}(\jepPubBbb{V}^r_i)=\mu_{H_r}(\jepPubBbb{V}^r_0)+i\mu_{H_r}(\jepPubBbb{V}^r_0{}^\star\otimes\jepPubBbb{V}^r_1)$.
\end{proposition}
\begin{proof}
The $H_r$-modules $\jepPubBbb{V}^r_i$ are irreducible by Lemma~2.32, thus polystable. The same holds for the $H_r$-module $\mathfrak l_r^-\simeq\jepPubBbb{V}^r_0{}^\star\otimes\jepPubBbb{V}^r_1$, and we have $\mu_{H_r}(\mathfrak l_r^-)=\mu_{H_r}(\jepPubBbb{V}^r_0{}^\star\otimes\jepPubBbb{V}^r_1)$. Let $i$ be fixed. By Lemma~2.44, $\jepPubBbb{V}^r_i\otimes\mathfrak l_r^-$ is polystable. Since, by Lemma~2.31, $\jepPubBbb{V}^r_{i+1}$ is a quotient of $\jepPubBbb{V}^r_i\otimes\mathfrak l_r^-$, it follows that $\mu_{H_r}(\jepPubBbb{V}^r_{i+1})=\mu_{H_r}(\jepPubBbb{V}^r_i)+\mu_{H_r}(\mathfrak l_r^-)$.
\end{proof}
\subsection{Interpretation in terms of homogeneous variations of Hodge structure}
As we learned from one of the referees, the cominuscule representation $\jepPubBbb{E}$ of $G$ has been interpreted in the language of Hodge structures, first by Gross \cite{jep93Gro94} in the tube type case, then by Sheng and Zuo \cite{jep93SZ10} in general. Indeed, the representation $\jepPubBbb{E}$ defines a canonical homogeneous complex variation of Hodge structure ($\mathbb C$-VHS) $E_{CY}:=\jepPubOplus_{p+q=r_M}E^{p,q}$ on $M$ or on any quotient $\Lambda\backslash M$, for $\Lambda$ a discrete subgroup of $G_{\mathbb R}$. The Hodge bundle $E^{p,q}$ is just the bundle over $M=G_{\mathbb R}/K_{\mathbb R}$ associated with the $K_{\mathbb R}$-module $\jepPubBbb{E}_{r_M-p}$. Geometrically, the corresponding filtration is given by the successive osculating tangent spaces to the cone over the subvariety $M\subset\mathbb P\jepPubBbb{E}$. This $\mathbb C$-VHS is of particular interest because it is of Calabi-Yau type, meaning that $\dim E^{r_M,0}=1$ and $E^{r_M-1,1}\simeq E^{r_M,0}\otimes\Omega_M^1$. It is a real variation of Hodge structure ($\mathbb R$-VHS) of weight $r_M$ precisely when $M$ is of tube type. This is exactly the same thing as saying, as we did, that $\jepPubBbb{E}$ is self-dual in this case.

From this point of view, what we do here can be interpreted in the following way. Let us work on the projectivized tangent bundle $\mathbb PT_M$ of $M$. The notion of rank of an element of $\mathfrak m^-$ naturally extends to elements in the tangent bundle of $M$, and the bundle $\mathbb PT_M$ is the union of $r_M$ $G_{\mathbb R}$-orbits $\operatorname{Orb}_1,\ldots,\operatorname{Orb}_{r_M}$, where $\operatorname{Orb}_r:=\{[\eta]\in\mathbb PT_M\mid\operatorname{rk}(\eta)=r\}$. It is also the quotient $G_{\mathbb R}/(K_{\mathbb R}\cap Z_r)$, where $Z_r$ is the centralizer of a nilpotent element $y$ in $\mathfrak m^-$ of rank $r$. The sets $\cup_{s\leqslant r}\operatorname{Orb}_s$ are the characteristic varieties of Mok \cite{jep93Mok89}. Let us fix $1\leqslant r\leqslant r_M$ and call $\pi_M$ the projection $\mathbb PT_M\to M$. The definition of the weight filtration $\jepPubBbb{W}^r_\jepPubBullet$ of $\jepPubBbb{E}$ associated with a nilpotent element $y\in\mathfrak m^-$ or rank $r$ can also be extended to give an increasing holomorphic weight filtration on the restriction of the $\mathbb C$-VHS $\pi_M^\star E_{CY}$ to $\operatorname{Orb}_r$, compatible with the decreasing Hodge filtration $(\jepPubOplus_{p\geqslant k}\pi_M^\star E^{p,q})_k$, and this defines a homogeneous family of ``complex'' mixed Hodge structures on $\operatorname{Orb}_r$. When $M$ has tube type, $\pi_M^\star E_{CY}$ is an $\mathbb R$-VHS and we have a real mixed Hodge structure on each fiber of $\pi_M^\star E_{CY}$ above $\operatorname{Orb}_r$.
\begin{example}
We keep the notations of Example~2.42, assuming now that we are in the tube type case: the group $G_{\mathbb R}$ is $\operatorname{SU}(p,p)$ with maximal compact subgroup $K_{\mathbb R}=\operatorname S(\operatorname U(p)\times\operatorname U(p))$. Let $C$, resp. $D$, be the subspace of $\mathbb C^{2p}$ generated by $(e_1,\ldots,e_p)$, resp. $(e_{p+1},\ldots,e_{2p})$. In this case, the cominuscule representation is
\[
\jepPubBbb{E}=\wedge^p(C\oplus D)=\jepPubOplus_{m+n=p}\wedge^m C\otimes\wedge^n D.
\]
Whenever $V$ is a complex vector space, we denote by $\overline V$ the complex vector space $V$ with the $\mathbb C$-action defined by $\lambda\cdot v=\overline\lambda v$. Note that, with this convention, a real structure on $V$ is the same as an involutory $\mathbb C$-linear isomorphism $\overline V\to V$. As representations of $K_{\mathbb R}$, $\overline C\simeq C^\star$, $\overline D\simeq D^\star$, and $\wedge^{2p}(C\oplus D)$ is the trivial representation. We get a $K_{\mathbb R}$-equivariant isomorphism
\[
\overline{\wedge^m C\otimes\wedge^n D}\simeq\wedge^m\overline C\otimes\wedge^n\overline D\simeq\wedge^m C^\star\otimes\wedge^n D^\star\simeq(\wedge^m C\otimes\wedge^n D)^\star\simeq\wedge^{p-m}C\otimes\wedge^{p-n}D,
\]
where the last isomorphism is given by the perfect pairing
\[
(\wedge^m C\otimes\wedge^n D)\otimes(\wedge^{p-m}C\otimes\wedge^{p-n}D)\longrightarrow\wedge^p C\otimes\wedge^p D\simeq\wedge^{2p}(C\oplus D)\simeq\mathbb C.
\]
Therefore we get a real structure on $\jepPubBbb{E}$, invariant by $K_{\mathbb R}$, using the isomorphism
\[
\overline{\jepPubBbb{E}}=\jepPubOplus_{m+n=p}\overline{\wedge^m C\otimes\wedge^n D}\simeq\jepPubOplus_{m+n=p}\wedge^n C\otimes\wedge^m D=\jepPubBbb{E}.
\]
Since we have tautologically, for this real structure, $\overline{\wedge^m C\otimes\wedge^n D}=\wedge^n C\otimes\wedge^m D$, this defines an $\mathbb R$-VHS of weight $p=r_M$ on the symmetric space $M=G_{\mathbb R}/K_{\mathbb R}$. This is the homogeneous VHS defined in \cite{jep93Gro94}.

If now $y$ is an element in $\mathfrak m^-$ of rank $r$, we moreover have the decompositions $C=A\oplus N$ and $D=B\oplus I$, as in Example~2.42, where the restriction of the action of $y$ to $A$ gives an isomorphism $A\simeq I$. The ($p$-shifted) increasing weight filtration $\jepPubBbb{W}^r_\jepPubBullet$ of $\jepPubBbb{E}$ associated to the action of $y$ is given by
\[
\jepPubBbb{W}^r_s=\jepPubOplus_{\substack{i-\ell\leqslant s-p\\i+j+k+\ell=p}}\wedge^i A\otimes\wedge^j N\otimes\wedge^k B\otimes\wedge^\ell I,
\]
whereas the decreasing Hodge filtration $\jepPubBbb{F}^\jepPubBullet$ is given by
\[
\jepPubBbb{F}^t=\jepPubOplus_{\substack{i+j\geqslant t\\i+j+k+\ell=p}}\wedge^i A\otimes\wedge^j N\otimes\wedge^k B\otimes\wedge^\ell I.
\]
Therefore the graded parts are:
\[
\operatorname{Gr}^{\jepPubBbb{F}}_t\operatorname{Gr}^{\jepPubBbb{W}^r}_s\jepPubBbb{E}=\jepPubOplus_{\substack{i-\ell=s-p\\i+j=t\\i+j+k+\ell=p}}\wedge^i A\otimes\wedge^j N\otimes\wedge^k B\otimes\wedge^\ell I.
\]
The intersection $Z_r\cap K_{\mathbb R}$ of $K_{\mathbb R}$ with the centralizer $Z_r$ of $y$ is the group
\[
\{(a,b,c)\in\operatorname U(r)\times\operatorname U(p-r)\times\operatorname U(p-r)\mid\det(a^2bc)=1\}.
\]
As $(Z_r\cap K_{\mathbb R})$-modules, we have:
\begin{align*}
\overline{\wedge^i A\otimes\wedge^j N\otimes\wedge^k B\otimes\wedge^\ell I}&\simeq\wedge^i A^\star\otimes\wedge^j N^\star\otimes\wedge^k B^\star\otimes\wedge^\ell I^\star\\
&\simeq(\wedge^i A^\star\otimes\wedge^j N^\star\otimes\wedge^k B^\star\otimes\wedge^\ell I^\star)\otimes\wedge^{2p}(A\oplus N\oplus B\oplus I)\\
&\simeq\wedge^{r-i}A\otimes\wedge^{p-r-j}N\otimes\wedge^{p-r-k}B\otimes\wedge^{r-\ell}I\\
&\simeq\wedge^{r-\ell}A\otimes\wedge^{p-r-j}N\otimes\wedge^{p-r-k}B\otimes\wedge^{r-i}I\\
&\simeq\wedge^{i'}A\otimes\wedge^{j'}N\otimes\wedge^{k'}B\otimes\wedge^{\ell'}I
\end{align*}
with $i'=r-\ell$, $j'=p-r-j$, $k'=p-r-k$ and $\ell'=r-i$ (on the fourth line we used the isomorphism $I\simeq A$). Therefore, if $i-\ell=s-p$ and $i+j=t$, then $i'-\ell'=i-\ell=s-p$ and $i'+j'=p-j-\ell=s-(i+j)=s-t$. Hence the real structure on $\jepPubBbb{E}$ defined as above with these isomorphisms is such that
\[
\overline{\operatorname{Gr}^{\jepPubBbb{F}}_t\operatorname{Gr}^{\jepPubBbb{W}^r}_s\jepPubBbb{E}}=\operatorname{Gr}^{\jepPubBbb{F}}_{s-t}\operatorname{Gr}^{\jepPubBbb{W}^r}_s\jepPubBbb{E}
\]
and the Hodge filtration induces a real Hodge structure of weight $s$ on $\operatorname{Gr}^{\jepPubBbb{W}^r}_s\jepPubBbb{E}$, so that the two filtrations $\jepPubBbb{W}^r_\jepPubBullet$ and $\jepPubBbb{F}^\jepPubBullet$ indeed define a real mixed Hodge structure on $\jepPubBbb{E}$.
\end{example}
However, again from the Hodge structures point of view, our purpose here is a bit different. We would like to construct a canonical ``weak sub-$\mathbb R$-VHS'' of weight $r$ of the restriction of the $\mathbb C$-VHS $\pi_M^\star E_{CY}$ to $\operatorname{Orb}_r$. This is the role played by the submodule $\jepPubBbb{V}^r$. The meaning of ``weak $\mathbb R$-VHS'' is that if $V^r$ denotes the equivariant holomorphic bundle $\jepPubBbb{V}^r\otimes\jepPubScript{O}_{\operatorname{Orb}_r}$ on $\operatorname{Orb}_r$ with its Gauss-Manin connection $\nabla$, and if $F^\jepPubBullet$ denotes the Hodge filtration on $V^r$ defined by the filtration $\jepPubBbb{F}^\jepPubBullet$ of $\jepPubBbb{E}$, then $(V^r,F^\jepPubBullet)$ does have the symmetries of an $\mathbb R$-VHS by Proposition~2.40, but does not satisfy Griffiths transversality $\nabla(F^t)\subset F^{t-1}\otimes\Omega^1_{\operatorname{Orb}_r}$. Rather, $(V^r,F^\jepPubBullet)$ satisfies a weak transversality condition by Proposition~2.29. Indeed we have $\nabla(F^t\otimes\jepPubScript{O}_{\mathbb PT_M}(-1)\jepPubRestrict_{\operatorname{Orb}_r})\subset F^{t-1}$, where $\jepPubScript{O}_{\mathbb PT_M}(-1)$ is the tautological line bundle on $\mathbb PT_M$. This will allow us to construct \emph{leafwise} Higgs subsheaves (but not plain Higgs subsheaves) in Section~4.3.

\section{Higgs bundles on foliated Kähler manifolds}
\subsection{Harmonic Higgs bundles}
We keep the notation of the previous sections. In particular, $G_{\mathbb R}$ is a simple real algebraic Hermitian group, $K_{\mathbb R}$ its maximal compact subgroup, $M=G_{\mathbb R}/K_{\mathbb R}$ the associated irreducible Hermitian symmetric space of the noncompact type, $G$ and $K$ are the complexifications of $G_{\mathbb R}$ and $K_{\mathbb R}$, and $\mathfrak g=\mathfrak k\oplus\mathfrak m$ and $\mathfrak g_{\mathbb R}=\mathfrak k_{\mathbb R}\oplus\mathfrak m_{\mathbb R}$ are the associated Cartan decompositions of the Lie algebras of $G$ and $G_{\mathbb R}$.

Let now $(Y,\omega_Y)$ be a compact Kähler manifold, $\Gamma=\pi_1(Y)$ its fundamental group, and $\rho:\Gamma\to G_{\mathbb R}$ be a \emph{reductive representation} (group homomorphism) of $\Gamma$ into $G_{\mathbb R}$. The assumption that $\rho$ is reductive means that the real Zariski closure of $\rho(\Gamma)$ in $G_{\mathbb R}$ is a reductive group.

In this case, by \cite{jep93Cor88}, there exists a $\rho$-equivariant harmonic map $f$ from the universal cover $\widetilde Y$ of $Y$ to the symmetric space $M=G_{\mathbb R}/K_{\mathbb R}$ associated with $G_{\mathbb R}$. The fact that $Y$ is Kähler and the nonpositivity of the complexified sectional curvature of $M$ imply by a Bochner formula due to \cite{jep93Sam78,jep93Siu80} that the map $f$ is pluriharmonic (i.e., its restriction to 1-dimensional complex submanifolds of $Y$ is still harmonic), and that the image of its (complexified) differential at every point $y\in Y$ is an abelian subalgebra of $T^{\mathbb C}_{f(y)}M$ identified with $\mathfrak m$.

By the work of Hitchin and Simpson, this gives a \emph{harmonic $G_{\mathbb R}$-Higgs principal bundle} $(P_K,\theta)$ on $Y$. We will now briefly describe the construction and the properties of such a Higgs bundle. Details and proofs can be found in the original papers \cite{jep93Hit87,jep93Hit92,jep93Sim88,jep93Sim92}.

There is a flat principal bundle $P_{G_{\mathbb R}}\to Y$ of group $G_{\mathbb R}$ associated with the representation $\rho$. The $\rho$-equivariant map $f:\widetilde Y\to G_{\mathbb R}/K_{\mathbb R}$ defines a reduction of its structure group to $K_{\mathbb R}$, i.e., a principal $K_{\mathbb R}$ bundle $P_{K_{\mathbb R}}\subset P_{G_{\mathbb R}}$. The differential of $f$ can be seen as a 1-form with values in $P_{K_{\mathbb R}}(\mathfrak m_{\mathbb R}):=(P_{K_{\mathbb R}}\times\mathfrak m_{\mathbb R})/K_{\mathbb R}$, the vector bundle associated with the adjoint action of $K_{\mathbb R}$ on $\mathfrak m_{\mathbb R}$.

If we enlarge the structure group of $P_{K_{\mathbb R}}$ to $K$, the pluriharmonicity of $f$ implies that the $K$-principal bundle $P_K\to Y$ is a holomorphic bundle and that the (1,0)-part $d^{1,0}f:T^{1,0}Y\to T^{\mathbb C}M$ of the complexified differential of $f$ defines a holomorphic section $\theta$ of $P_K(\mathfrak m)\otimes\Omega_Y^1$, where $P_K(\mathfrak m)$ is the holomorphic vector bundle associated with the principal bundle $P_K$ and the adjoint representation of $K$ on $\mathfrak m$. The section $\theta$ is called the \emph{Higgs field} and satisfies the integrability condition $[\theta,\theta]=0$ as a section of $P_K(\mathfrak m)\otimes\Omega_Y^2$. The pair $(P_K,\theta)$ is called a $G_{\mathbb R}$-Higgs principal bundle on $Y$, see e.g. \cite[\S2.2]{jep93BGPG06}.

If now $\jepPubBbb{E}$ is a (complex) representation of $G$ we can construct the associated holomorphic vector bundle $E:=P_K(\jepPubBbb{E})$ over $Y$. Since $\jepPubBbb{E}$ is a representation of $\mathfrak g$ and not only of $\mathfrak k$, we have a morphism $P_K(\mathfrak m)\to P_K(\operatorname{End}(\jepPubBbb{E}))=\operatorname{End}(E)$. Thus, the Higgs field $\theta$ can be seen as a holomorphic 1-form with values in the endomorphisms of $E$, i.e., a section of $\operatorname{End}(E)\otimes\Omega_Y^1$. The pair $(E,\theta)$ is called a \emph{$G_{\mathbb R}$-Higgs vector bundle} on $Y$. The harmonic map $f$, seen as a reduction of the structure group of $P_K$ to the compact subgroup $K_{\mathbb R}$, together with a $K_{\mathbb R}$-invariant metric on $\jepPubBbb{E}$, gives a Hermitian metric on $(E,\theta)$ called the \emph{harmonic metric}.

The existence of this harmonic metric and the fact that $P_K$ comes from a flat principal $G_{\mathbb R}$ bundle imply that for any representation $\jepPubBbb{E}$ of $G$, the associated Higgs vector bundle $(E,\theta)\to Y$ is Higgs polystable of degree 0, see \cite{jep93Sim88}. To explain what Higgs polystability means, we first define Higgs subsheaves of the Higgs bundle $(E,\theta)$. A coherent subsheaf $\jepPubScript{F}$ of $\jepPubScript{O}_Y(E)$ is a \emph{Higgs subsheaf} if it is invariant by the Higgs field, i.e., if it satisfies $\theta(\jepPubScript{F})\subset\jepPubScript{F}\otimes\Omega_Y^1$. The Higgs vector bundle $(E,\theta)$ is said to be \emph{Higgs stable} if for any Higgs subsheaf $\jepPubScript{F}$ of $(E,\theta)$ such that $0<\operatorname{rk}\jepPubScript{F}<\operatorname{rk}E$, we have $\mu_{\omega_Y}(\jepPubScript{F})<\mu_{\omega_Y}(E)$, where $\mu_{\omega_Y}(\jepPubScript{F})$ is the slope of $\jepPubScript{F}$, i.e., its degree (computed w.r.t. the Kähler form $\omega_Y$ of $Y$) divided by its rank. The Higgs bundle $(E,\theta)$ is \emph{Higgs polystable} if it is a direct sum of Higgs stable Higgs vector bundles of the same slope. Note that here $E$ is flat as a $C^{\infty}$ bundle, so that its degree is zero.
\begin{remark}
Since moreover we assumed that $G_{\mathbb R}$ is a Hermitian group, then as a $K$-representation we have $\mathfrak m=\mathfrak m^+\oplus\mathfrak m^-$ and the Higgs field $\theta$ on the principal bundle $P_K$ (or on any associated vector bundle $E$) has two components $\beta\in P_K(\mathfrak m^-)\otimes\Omega_Y^1$ and $\gamma\in P_K(\mathfrak m^+)\otimes\Omega_Y^1$. The vanishing of $\gamma$, resp. $\beta$, means that the harmonic map $f$ is holomorphic, resp. antiholomorphic. The component $\beta$, resp. $\gamma$, will be called the holomorphic, resp. antiholomorphic, component of the Higgs field $\theta$.
\end{remark}
\subsection{Harmonic Higgs bundles on foliated Kähler manifolds}
Assume now that the base Kähler manifold $Y$ of the harmonic $G_{\mathbb R}$-Higgs vector bundle $(E,\theta)\to Y$ of degree $0$ admits a smooth holomorphic foliation by complex curves and that this foliation $\jepPubScript{T}$ admits a transverse volume form $\Omega$. This means that $\Omega$ is a nowhere vanishing semipositive closed $(d-1,d-1)$-form ($d$ is the dimension of $Y$) such that the interior product of $\Omega$ with any vector tangent to the foliation $\jepPubScript{T}$ is zero.

Our goal in this section is to briefly explain the interplay between the Higgs bundle and the foliation (equipped with its transverse volume form). Details can be found in \cite[\S2.2]{jep93KM17}.

We first weaken the notion of Higgs subsheaves of $(E,\theta)$ to \emph{leafwise Higgs subsheaves} by asking only an invariance by the Higgs field along the leaves. More precisely we now consider the Higgs field as a section of $\operatorname{Hom}(E\otimes L,E)$, where $L$ is the holomorphic line subbundle of $T_Y$ tangent of the foliation $\jepPubScript{T}$. A leafwise Higgs subsheaf $\jepPubScript{F}$ of $(E,\theta)$ is a subsheaf of $\jepPubScript{O}_Y(E)$ such that $\theta(\jepPubScript{F}\otimes L)\subset\jepPubScript{F}$.

We define the foliated degree $\deg_{\jepPubScript{T}}\jepPubScript{F}$ and the foliated slope $\mu_{\jepPubScript{T}}(\jepPubScript{F})$ of a coherent sheaf $\jepPubScript{F}$ on $Y$ by wedging the first Chern class of $\jepPubScript{F}$ with the transverse volume form $\Omega$:
\[
\deg_{\jepPubScript{T}}\jepPubScript{F}:=\int_Y c_1(\jepPubScript{F})\wedge\Omega,\quad\text{and}\quad\mu_{\jepPubScript{T}}(\jepPubScript{F})=\frac{\deg_{\jepPubScript{T}}\jepPubScript{F}}{\operatorname{rk}\jepPubScript{F}}.
\]
The link between the notions of slope (and stability) for modules and for sheaves used in this paper will be given by Fact~4.5 below.

In order to state the needed leafwise polystability property of the Higgs bundle $(E,\theta)$ w.r.t. the foliated degree, we first recall a few definitions.

An open subset of $Y$ is called \emph{big} if it is the complement of an analytic set of codimension at least $2$ in $Y$. Two line bundles on $Y$ which have isomorphic restrictions to a big open subset of $Y$ are in fact isomorphic over $Y$.

Let $\jepPubScript{F}$ be a subsheaf of $\jepPubScript{O}_Y(E)$. The complement of the biggest subset of $Y$ where $\jepPubScript{F}$ is the sheaf of sections of a subbundle $F$ of $E$ is called the \emph{singular locus} $\jepPubScript{S}(\jepPubScript{F})$ of $\jepPubScript{F}$. Equivalently, $\jepPubScript{S}(\jepPubScript{F})$ is the analytic subset of $Y$ where the quotient sheaf $\jepPubScript{O}_Y(E)/\jepPubScript{F}$ is not locally free. Recall that a subsheaf $\jepPubScript{F}$ of $\jepPubScript{O}_Y(E)$ is \emph{saturated} if $\jepPubScript{O}_Y(E)/\jepPubScript{F}$ is torsion free. The \emph{saturation} of $\jepPubScript{F}$ is the kernel of the morphism $\jepPubScript{O}_Y(E)\to(\jepPubScript{O}_Y(E)/\jepPubScript{F})_{\mathrm{tf}}$, where $(\jepPubScript{O}_Y(E)/\jepPubScript{F})_{\mathrm{tf}}$ denotes the quotient of $\jepPubScript{O}_Y(E)/\jepPubScript{F}$ by its torsion. If $\jepPubScript{F}$ is saturated, $\jepPubScript{S}(\jepPubScript{F})$ has codimension at least $2$ in $Y$, so that the \emph{regular locus} $Y\backslash\jepPubScript{S}(\jepPubScript{F})$ of $\jepPubScript{F}$ is a big open subset. Moreover a saturated subsheaf $\jepPubScript{F}$ is \emph{normal}, so that in particular for all big open subsets $\jepPubScript{U}$ of $Y$, the restriction $\jepPubScript{F}(Y)\to\jepPubScript{F}(\jepPubScript{U})$ is injective.

A subset of the foliated manifold $(Y,\jepPubScript{T})$ is said to be \emph{saturated under the foliation $\jepPubScript{T}$} if it is a union of leaves of $\jepPubScript{T}$. Obviously if $Z\subset Y$ is saturated under $\jepPubScript{T}$, then so is its complement $Y\backslash Z$. (Unfortunately, the word ``saturated'' is used here with two different meanings but this shouldn't cause any confusion.)

Here is the result (\cite[Prop.~2.2]{jep93KM17}).
\begin{proposition}
The harmonic Higgs bundle $(E,\theta)$ on the foliated Kähler manifold $(Y,\jepPubScript{T},\Omega)$ is weakly polystable along the leaves in the following sense:
\begin{enumerate}[label=(\arabic*)]
\item it is semistable along the leaves of $\jepPubScript{T}$: if $\jepPubScript{F}$ is a leafwise Higgs subsheaf of $(E,\theta)$, then $\deg_{\jepPubScript{T}}\jepPubScript{F}\leqslant0$.
\item if $\jepPubScript{F}$ is a saturated leafwise Higgs subsheaf of $(E,\theta)$ such that $\deg_{\jepPubScript{T}}\jepPubScript{F}=0$, then the singular locus $\jepPubScript{S}(\jepPubScript{F})$ of $\jepPubScript{F}$ is saturated under the foliation $\jepPubScript{T}$. Moreover, on $Y\backslash\jepPubScript{S}(\jepPubScript{F})$, and if $F$ denotes the holomorphic subbundle of $E$ such that $\jepPubScript{F}=\jepPubScript{O}_{Y\backslash\jepPubScript{S}(\jepPubScript{F})}(F)$ and $F^\perp$ its $C^\infty$ orthogonal complement w.r.t. the harmonic metric, then $\theta(F^\perp\otimes L)\subset F^\perp$. Finally, for each leaf $\jepPubLeaf{L}$ of $\jepPubScript{T}$ such that $\jepPubLeaf{L}\cap\jepPubScript{S}(\jepPubScript{F})=\varnothing$, $F^\perp\jepPubRestrict_{\jepPubLeaf{L}}$ is holomorphic on $\jepPubLeaf{L}$ and $(E,\theta)\jepPubRestrict_{\jepPubLeaf{L}}=(F,\theta\jepPubRestrict_F)\jepPubRestrict_{\jepPubLeaf{L}}\oplus(F^\perp,\theta\jepPubRestrict_{F^\perp})\jepPubRestrict_{\jepPubLeaf{L}}$ is a holomorphic orthogonal decomposition of Higgs bundles on $\jepPubLeaf{L}$.
\end{enumerate}
\end{proposition}
\subsection{The tautological foliation on the projectivized tangent bundle of a complex hyperbolic manifold}
An $n$-dimensional complex hyperbolic manifold $X$ is a quotient of the complex hyperbolic $n$-space $\mathbb H^n_{\mathbb C}=\operatorname{SU}(1,n)/\operatorname U(n)$ by a discrete torsion free subgroup $\Gamma$ of $\operatorname{SU}(1,n)$. It is a Hermitian locally symmetric space of rank 1. The complex hyperbolic space $\mathbb H^n_{\mathbb C}$ can be realized as the subset of negative lines in $\mathbb C^{n+1}$ for a Hermitian form $h$ of signature $(n,1)$ (recall our convention in Example~2.1), which is an open subset in the projective space $\mathbb{CP}^n$.

Intersections of lines of $\mathbb{CP}^n$ with $\mathbb H^n_{\mathbb C}$ are totally geodesic complex subspaces of $\mathbb H^n_{\mathbb C}$ isometric to the Poincaré disc. They are called \emph{complex geodesics}. The space $\jepPubScript{G}$ of complex geodesics is the homogeneous space $\operatorname{SU}(1,n)/\operatorname S(\operatorname U(1,1)\times\operatorname U(n-1))$. It is also a complex manifold, since it can be realized as the open subset of the Grassmannian manifold $\operatorname{Gr}_2(\mathbb C^{n+1})$ of $2$-planes in $\mathbb C^{n+1}$ consisting of planes on which the signature of $h$ is $(1,1)$.

The \emph{projectivized tangent bundle} $\mathbb PT_{\mathbb H^n_{\mathbb C}}$ of $\mathbb H^n_{\mathbb C}$ is the space of lines in the holomorphic tangent bundle $T_{\mathbb H^n_{\mathbb C}}$ of $\mathbb H^n_{\mathbb C}$. As a homogeneous space it can be identified with $\operatorname{SU}(1,n)/\operatorname S(\operatorname U(1)\times\operatorname U(1)\times\operatorname U(n-1))$. Again, its also a complex manifold (in fact a Kähler one) since it identifies with an open subset of the manifold $\mathrm F_{1,2}(\mathbb C^{n+1})$ of incomplete flags $(\ell\subset\Pi\subset\mathbb C^{n+1})$ with $\dim\ell=1$, $\dim\Pi=2$. A flag $(\ell\subset\Pi\subset\mathbb C^{n+1})$ belong to $\mathbb PT_{\mathbb H^n_{\mathbb C}}$ if the line $\ell$ is negative and the plane $\Pi$ has signature $(1,1)$ for the Hermitian form $h$.

The natural $\operatorname{SU}(1,n)$-equivariant fibration $\pi_{\jepPubScript{G}}:\mathbb PT_{\mathbb H^n_{\mathbb C}}\to\jepPubScript{G}$ which associates to a tangent line to $\mathbb H^n_{\mathbb C}$ the complex geodesic it defines is a disc bundle over $\jepPubScript{G}$.

By $\operatorname{SU}(1,n)$-equivariance, this fibration endows the projectivized tangent bundle $\mathbb PT_X=\Gamma\backslash\mathbb PT_{\mathbb H^n_{\mathbb C}}$ of $X=\Gamma\backslash\mathbb H^n_{\mathbb C}$ with a smooth holomorphic foliation $\jepPubScript{T}$ by complex curves, whose leaves are given by the tangent spaces of the (immersed) complex geodesics in $X$. This foliation is called the \emph{tautological foliation} of $\mathbb PT_X$ because the tangent line bundle $L$ to the leaves is naturally isomorphic to the tautological line bundle $\jepPubScript{O}_{\mathbb PT_X}(-1)$ on $\mathbb PT_X$.

The space $\jepPubScript{G}$ of complex geodesics of $\mathbb H^n_{\mathbb C}$ is a pseudo-Kähler manifold: it admits an $\operatorname{SU}(1,n)$-invariant Kähler form $\omega_{\jepPubScript{G}}$, which is closed, of type $(1,1)$, non degenerate, but not positive definite. The form $\omega_{\jepPubScript{G}}$ is moreover unique up to scaling. This form defines a transverse volume form $\Omega_{\jepPubScript{G}}$ for the tautological foliation $\jepPubScript{T}$ on $\mathbb PT_X$, and the associated notion of foliated degree $\deg_{\jepPubScript{T}}$ for sheaves on $\mathbb PT_X$ has the following fundamental property (after a suitable normalization of the involved forms) \cite[Prop.~3.1]{jep93KM17}.
\begin{proposition}
Assume that $X=\Gamma\backslash\mathbb H^n_{\mathbb C}$ is compact and let $\pi:\mathbb PT_X\to X$ be the projectivized tangent bundle of $X$. If $\jepPubScript{F}$ is a coherent $\jepPubScript{O}_X$-sheaf, then $\deg_{\jepPubScript{T}}(\pi^*\jepPubScript{F})=\deg_\omega\jepPubScript{F}$, where $\deg_\omega\jepPubScript{F}$ is the usual degree of $\jepPubScript{F}$ computed w.r.t. the Kähler form $\omega$ on $X$.
\end{proposition}

\section{The Milnor-Wood inequality}
Let $\Gamma$ be a uniform lattice in $\operatorname{SU}(1,n)$, $X$ be the compact quotient $\Gamma\backslash\mathbb H^n_{\mathbb C}$, and let $\rho$ be a representation of $\Gamma$ in a simple real algebraic Hermitian Lie group $G_{\mathbb R}$, whose associated symmetric space is denoted by $M$. In this section we use the material developed or recalled in Sections~2 and~3 to prove the Milnor-Wood inequality
\[
|\tau(\rho)|\leqslant r_M\operatorname{vol}(X).
\]
Before going into the proof, we remark that if the lattice $\Gamma$ has torsion, the measure defined by the Kähler form $\omega$ on $\mathbb H^n_{\mathbb C}$ descends to $X$ so that $\operatorname{vol}(X)$ is well-defined. There are several ways to define the Toledo invariant of the representation $\rho$ of $\Gamma$ in this case. Since this is how we will handle lattices with torsion, we observe that by Selberg's Lemma (see \cite[p.~327]{jep93Rat06} or \cite{jep93Sel60}) there exists a finite index torsion free subgroup $\Gamma'$ of $\Gamma$. We then define the Toledo invariant of $\rho:\Gamma\to G_{\mathbb R}$ by $\tau(\rho):=\tau(\rho')/|\Gamma/\Gamma'|$, where $\rho'=\rho\jepPubRestrict_{\Gamma'}$ is the restriction of $\rho$ to $\Gamma'$ and $\tau(\rho')$ is defined as in the introduction. One checks easily that this does not depend on the choice of $\Gamma'$. Moreover, if $X'=\Gamma'\backslash\mathbb H^n_{\mathbb C}$, then $\operatorname{vol}(X)=\operatorname{vol}(X')/|\Gamma/\Gamma'|$ so that proving the Milnor-Wood inequality for $\rho:\Gamma\to G_{\mathbb R}$ is equivalent to proving the Milnor-Wood inequality for $\rho':\Gamma'\to G_{\mathbb R}$.

In addition, the representation $\rho$ can always be deformed to a reductive representation with the same Toledo invariant, see e.g. \cite[Lem.~4.11]{jep93KM17}. Moreover, we may change the complex structure of $M$ to its opposite (i.e., replace $z\in\mathfrak t$ by $-z\in\mathfrak t$, or equivalently $\mathfrak m^{\pm}$ by $\mathfrak m^{\mp}$), thereby changing the sign of the Toledo invariant of $\rho$.

Therefore we may assume without loss of generality that
\begin{itemize}
\item the lattice $\Gamma$ is torsion free, so that $X=\Gamma\backslash\mathbb H^n_{\mathbb C}$ is a compact complex hyperbolic manifold of (complex) dimension $n$;
\item the representation $\rho$ is reductive, so that the results discussed in Section~3 are available;
\item the Toledo invariant $\tau(\rho)$ is positive, so we only have to prove $\tau(\rho)\leqslant r_M\operatorname{vol}(X)$.
\end{itemize}
(Observe that this implies that the harmonic map $f$ is not antiholomorphic.)

We also assume as in Section~2 that
\begin{itemize}\item the complexification $G$ of $G_{\mathbb R}$ is simply connected.\end{itemize}
See Remarks~5.14 and~5.15 where we explain why this is also not a loss of generality.
\subsection{Setup}
Consider the representation $\jepPubBbb{E}$ of $G$ defined in Section~2.2. As explained in Section~3.1, this gives rise to a flat harmonic $G_{\mathbb R}$-Higgs vector bundle $(\overline E,\overline\theta)$ over $X$, associated to the $G_{\mathbb R}$-Higgs principal bundle $(\overline P_K,\overline\theta)$. As a representation of $K$, $\jepPubBbb{E}=\jepPubOplus_{i=0}^{r_M}\jepPubBbb{E}_i$, where $r_M$ is the real rank of $G_{\mathbb R}$. This means that the holomorphic bundle $\overline E$ admits the holomorphic decomposition $\overline E=\jepPubOplus_{i=0}^{r_M}\overline E_i$, which is orthogonal for the harmonic metric. Moreover, the components $\overline\beta\in\overline P_K(\mathfrak m^-)\otimes\Omega_X^1$ and $\overline\gamma\in\overline P_K(\mathfrak m^+)\otimes\Omega_X^1$ of the Higgs field $\overline\theta\in\operatorname{Hom}(\overline E,\overline E)\otimes\Omega_X^1$ (see Remark~3.1) satisfy
\[
\overline\beta\in\jepPubOplus_{i=0}^{r_M-1}\operatorname{Hom}(\overline E_i,\overline E_{i+1})\otimes\Omega_X^1\quad\text{and}\quad\overline\gamma\in\jepPubOplus_{i=0}^{r_M-1}\operatorname{Hom}(\overline E_{i+1},\overline E_i)\otimes\Omega_X^1.
\]
We pull-back the harmonic Higgs bundle $(\overline E,\overline\theta)\to X$ to the projectivized tangent bundle $\mathbb PT_X$ of $X$ to obtain a harmonic Higgs bundle that we denote by $(E,\theta)\to\mathbb PT_X$. We restrict the Higgs field $\theta$ to the tangent space $L$ of the tautological foliation on $\mathbb PT_X$, so that its components $\beta$ and $\gamma$ satisfy
\[
\beta\in\jepPubOplus_{i=0}^{r_M-1}\operatorname{Hom}(E_i\otimes L,E_{i+1})\quad\text{and}\quad\gamma\in\jepPubOplus_{i=0}^{r_M-1}\operatorname{Hom}(E_{i+1}\otimes L,E_i).
\]
We also call $P_K$ the pull-back of the principal bundle $\overline P_K$ to $\mathbb PT_X$.
\begin{notation}
A tangent vector $\xi\neq0$ to $X$ at a point $x\in X$ defines a point $[\xi]$ in the projectivized tangent bundle $\mathbb PT_X$. It also defines an element in the fiber $\jepPubScript{O}_{\mathbb PT_X}(-1)_{[\xi]}=\{v\in T_{X,x}\mid v\in[\xi]\}$ of the tautological line bundle $\jepPubScript{O}_{\mathbb PT_X}(-1)$ at $[\xi]$, hence an element in the fiber $L_{[\xi]}$ of the tangent line bundle $L$ to the tautological foliation $\jepPubScript{T}$ at $[\xi]$, which will also be denoted by $\xi$.
\end{notation}
\begin{definition}
For $\xi\neq0$ a vector tangent to $X$ and $[\xi]$ the corresponding point in $\mathbb PT_X$, the \emph{rank} $\operatorname{rk}\beta_{[\xi]}$ of $\beta_{[\xi]}$ is the rank of the corresponding element $\beta(\xi)$ in $\mathfrak m^-$, as defined in Definition~2.22. By Proposition~2.24, it is also the largest value of $k$ such that $(\beta_{[\xi]})^k:E_0\otimes L^k\to E_k$ is not zero.
\begin{itemize}
\item The \emph{generic rank} $\operatorname{rk}\beta$ of $\beta$ is the maximum of the ranks of $\beta_{[\xi]}$ for $[\xi]\in\mathbb PT_X$.
\item The \emph{singular locus} of $\beta$ is the following subset of $\mathbb PT_X$:
\begin{align*}
\jepPubScript{S}(\beta)&:=\{[\xi]\in\mathbb PT_X\mid\operatorname{rk}\beta_{[\xi]}<\operatorname{rk}\beta\}\\
&=\{[\xi]\in\mathbb PT_X\mid(\beta_{[\xi]})^{\operatorname{rk}\beta}:E_0\otimes L^{\operatorname{rk}\beta}\to E_{\operatorname{rk}\beta}\text{ vanishes}\}.
\end{align*}
\item The \emph{regular locus} of $\beta$ is $\jepPubScript{R}(\beta):=\mathbb PT_X\backslash\jepPubScript{S}(\beta)$.
\item The \emph{singular locus} $\jepPubScript{S}(\overline\beta)$ of $\overline\beta$ is the projection to $X$ of $\jepPubScript{S}(\beta)$.
\end{itemize}
The \emph{regular locus} $\jepPubScript{R}(\overline\beta)$ of $\overline\beta$ is $X\backslash\jepPubScript{S}(\overline\beta)$ (it is not the projection to $X$ of $\jepPubScript{R}(\beta)$).
\end{definition}
Our assumption that $\tau(\rho)>0$ implies that $\overline\beta\neq0$ and $\beta\neq0$, therefore $1\leqslant\operatorname{rk}\beta\leqslant r_M$. Observe that while $\jepPubScript{S}(\beta)$ is an analytic subset of $\mathbb PT_X$ of codimension at least 1, its projection $\jepPubScript{S}(\overline\beta)$, although it is an analytic subset of $X$ (because $\pi:\mathbb PT_X\to X$ is a proper map), might be equal to the whole base $X$. A very important consequence of the weak polystability along the leaves (Proposition~3.2~(2)), is that this does not happen when there is equality in the refined Milnor-Wood inequality of Theorem~4.6, see Proposition~4.7. This will indeed be a crucial point when dealing with maximal representations.
\subsection{Rewording of the inequality}
Since the Hermitian symmetric space $M$ associated with $G_{\mathbb R}$ is a Kähler-Einstein manifold, the first Chern form $c_1(T_M)$ of its tangent bundle is a constant multiple of the $G_{\mathbb R}$-invariant Kähler form $\omega_M$: $c_1(T_M)=-\frac1{4\pi}c\,\omega_M$ for some positive constant $c$ ($c=c_1(\check M)$ is the first Chern number of the compact dual $\check M$ of $M$). On the other hand, the line bundle $\jepPubScript{L}$ associated with the $K$-representation $\jepPubBbb{E}_0$ is a generator of the Picard group of the compact dual $\check M$ of $M$ and it can be checked that the canonical bundle $K_{\check M}$ of $\check M$ is precisely given by $\jepPubScript{L}^{-c}$, see e.g. \cite[\S2]{jep93KM10}. Therefore the pull-back $f^\star\omega_M$ is $4\pi$ times the first Chern form of the line bundle $f^\star\jepPubScript{L}=\overline E_0$, so that the Toledo invariant of $\rho$ is
\begin{equation}\tag{4.1}
\tau(\rho)=4\pi\deg(\overline E_0)=4\pi\deg_{\jepPubScript{T}}(E_0),
\end{equation}
where the last equality follows from Proposition~3.3. Similarly, we get that
\begin{equation}\tag{4.2}
\deg(K_X)=\frac{n+1}{4\pi}\operatorname{vol}(X).
\end{equation}
Recall that $L$ is the tangential line bundle to the tautological foliation $\jepPubScript{T}$ on the projectivized tangent bundle $\mathbb PT_X$, and let $L^\star$ be its dual line bundle. One can compute as explained in \cite[\S4.3.1]{jep93KM17} that
\begin{equation}\tag{4.3}
\deg_{\jepPubScript{T}}(L^\star)=\frac1{2\pi}\operatorname{vol}(X).
\end{equation}
Therefore, the Milnor-Wood inequality can be rephrased as an inequality between the foliated degrees of the line bundles $E_0$ and $L^\star$ on $\mathbb PT_X$ and we need to prove:
\begin{equation}\tag{4.4}
\deg_{\jepPubScript{T}}(E_0)\leqslant\frac{r_M}{2}\deg_{\jepPubScript{T}}(L^\star),
\end{equation}
where $r_M$ is the rank of the symmetric space $M$.
\subsection{Leafwise Higgs subsheaves associated with the holomorphic component of the Higgs field}
We now define a subsheaf $\jepPubScript{V}$ of $\jepPubScript{E}:=\jepPubScript{O}(E)$ associated with $\beta$ (recall that $\beta\neq0$) in the same way we defined the submodule $\jepPubBbb{V}^r$ of $\jepPubBbb{E}$ associated with the nilpotent element $y\in\mathfrak m^-$ of rank $r$ in Definition~2.25 (for all $\xi\in L$, $\beta(\xi)\in P_K(\mathfrak m^-)$ is a nilpotent endomorphism of the bundle $E$). This subsheaf will be shown to be a leafwise Higgs subsheaf of the Higgs bundle $(E,\theta)$ on $\mathbb PT_X$. In Section~4.4 this will be used to prove the Milnor-Wood inequality.

More precisely, mimicking Definition~2.25 and for $-\operatorname{rk}\beta\leqslant k\leqslant\operatorname{rk}\beta$, we define the following subsheaves of $\jepPubScript{E}$:
\[
\jepPubScript{W}_k:=\sum_{\substack{\ell\geqslant0\\k+\ell+1\geqslant0}}\operatorname{Ker}\beta^{k+\ell+1}\cap\operatorname{Im}\beta^\ell,
\]
where in order to define $\operatorname{Ker}\beta^j$ we consider $\beta^j$ as a sheaf morphism from $\jepPubScript{E}$ to $\jepPubScript{E}\otimes(L^\star)^j$ and to define $\operatorname{Im}\beta^j$ we consider $\beta^j$ as a sheaf morphism from $\jepPubScript{E}\otimes L^j$ to $\jepPubScript{E}$.

For $k=0,1,\ldots,\operatorname{rk}\beta$, let $\jepPubScript{V}_k$ be the saturation in $\jepPubScript{E}_k:=\jepPubScript{O}(E_k)$ of the subsheaf $\jepPubScript{E}_k\cap\jepPubScript{W}_{\operatorname{rk}\beta-2k}$ and set
\[
\jepPubScript{V}=\jepPubOplus_{0\leqslant k\leqslant\operatorname{rk}\beta}\jepPubScript{V}_k.
\]
Since the sheaves $\jepPubScript{V}_k$ are saturated subsheaves of $\jepPubScript{O}(E_k)$, there exist a big open subset $\jepPubScript{U}$ of $\mathbb PT_X$ and subbundles $V_k$ of $E_k$ defined on $\jepPubScript{U}$ such that the restriction of the $\jepPubScript{V}_k$'s to $\jepPubScript{U}$ are the sheaves of sections of the $V_k$'s. On $\jepPubScript{U}$ we let $V$ be the subbundle $\oplus_{0\leqslant k\leqslant r}V_k$, so that $\jepPubScript{V}\jepPubRestrict_{\jepPubScript{U}}=\jepPubScript{O}_{\jepPubScript{U}}(V)$.

Observe that on the regular locus $\jepPubScript{R}(\beta)$ of $\beta$, the rank of $\beta^k$, as a vector bundle morphism from $E\otimes L^k$ to $E$, is constant for all $k$. Indeed, by Definition~2.22, the rank of $\beta(\xi)$ viewed as an element of $\mathfrak m^-$ determines its $K$-orbit. Hence on this open subset the formulas used above to define the subsheaves $\jepPubScript{W}_k$ of $\jepPubScript{E}$ in fact define subbundles $W_k$ of $E$ such that $\jepPubScript{W}_k\jepPubRestrict_{\jepPubScript{R}(\beta)}=\jepPubScript{O}_{\jepPubScript{R}(\beta)}(W_k)$. Therefore, on $\jepPubScript{R}(\beta)$, the subbundles $V_k$ such that $\jepPubScript{V}_k\jepPubRestrict_{\jepPubScript{R}(\beta)}=\jepPubScript{O}(V_k)$ are given by $V_k=E_k\cap W_{\operatorname{rk}\beta-2k}$. Thus, we may assume that $\jepPubScript{R}(\beta)$ is contained (and therefore dense) in $\jepPubScript{U}$.

We view an element $p$ of the $K$-principal bundle $P_K\overset{\pi_K}{\longrightarrow}\mathbb PT_X$ above $\xi\in\mathbb PT_X$ as an isomorphism between the fiber $E_\xi$ of $E=P_K(\jepPubBbb{E})$ and the model space $\jepPubBbb{E}$. The component $\beta$ of the Higgs field is a section of $P_K(\mathfrak m^-)\otimes L^\star\subset P_K(\operatorname{End}(\jepPubBbb{E}))\otimes L^\star$.

Let $y=y^{\operatorname{rk}\beta}$ and $Q_{\operatorname{rk}\beta}$ be defined as in Section~2.4.
\begin{lemma}
On the big open set $\jepPubScript{U}$, the subsheaf $\jepPubScript{V}$ defines a reduction $P_{K\cap Q_{\operatorname{rk}\beta}}$ of the structure group of $P_K$ to the subgroup $K\cap Q_{\operatorname{rk}\beta}\subset K$.
\end{lemma}
\begin{proof}
We begin by working on $\jepPubScript{R}(\beta)\subset\jepPubScript{U}$. Since for all $\xi\in\jepPubScript{R}(\beta)$ and all $\eta\in L_\xi$, $\eta\neq0$, we have that $\beta_\xi(\eta)$ has rank $\operatorname{rk}\beta$, there exists $p\in(P_K)_\xi$ such that $p\circ\beta_\xi(\eta)\circ p^{-1}=y\in\mathfrak m^-\subset\operatorname{End}(\jepPubBbb{E})$, so that $p(V_\xi)=\jepPubBbb{V}^{\operatorname{rk}\beta}$. Therefore, on $\jepPubScript{R}(\beta)$, by Proposition~2.33~(4), the subbundle $V$ of $E$ defines a (holomorphic) reduction $P_{K\cap Q_{\operatorname{rk}\beta}}$ of the structure group of $P_K$ to the subgroup $K\cap Q_{\operatorname{rk}\beta}$ of $K$ ($Q_{\operatorname{rk}\beta}$ is the normalizer in $G$ of the parabolic subalgebra $\mathfrak q_{\operatorname{rk}\beta}$, see Definition~2.30). Explicitly $P_{K\cap Q_{\operatorname{rk}\beta}}=\{p\in P_K\mid p(V_{\pi_K(p)})=\jepPubBbb{V}^{\operatorname{rk}\beta}\}$.

We now work on $\jepPubScript{U}$. Enlarge the structure group of $P_K$ to $\operatorname{GL}(\jepPubBbb{E})$. The subbundle $V=\jepPubOplus_{k=0}^{\operatorname{rk}\beta}V_k$ of $E$ defines a reduction $P_S$ of the structure group of $P_{\operatorname{GL}(\jepPubBbb{E})}\overset{\pi_{\operatorname{GL}(\jepPubBbb{E})}}{\longrightarrow}\mathbb PT_X$ to the stabilizer $S$ of $\jepPubBbb{V}^{\operatorname{rk}\beta}$ in $\operatorname{GL}(\jepPubBbb{E})$ by setting $P_S=\{p\in P_{\operatorname{GL}(\jepPubBbb{E})}\mid p(V_{\pi_{\operatorname{GL}(\jepPubBbb{E})}(p)})=\jepPubBbb{V}^{\operatorname{rk}\beta}\}$.

Let $B\subset\jepPubScript{U}$ be an open ball on which $P_K$ is trivial. Then the reductions $P_{K\cap Q_{\operatorname{rk}\beta}}$ of $P_K$ on $B\cap\jepPubScript{R}(\beta)$ and $P_S$ of $P_{\operatorname{GL}(\jepPubBbb{E})}$ on $B$ are respectively given by holomorphic maps $\sigma:B\cap\jepPubScript{R}(\beta)\to K/(K\cap Q_{\operatorname{rk}\beta})$ and $s:B\to\operatorname{GL}(\jepPubBbb{E})/S$. Moreover, if $\iota$ denotes the natural map $K/(K\cap Q_{\operatorname{rk}\beta})\to\operatorname{GL}(\jepPubBbb{E})/S$, which is injective, then we have $s=\iota\circ\sigma$ on $B\cap\jepPubScript{R}(\beta)$. Since $K/(K\cap Q_{\operatorname{rk}\beta})$ is compact, its image by $\iota$ is closed in $\operatorname{GL}(\jepPubBbb{E})/S$. Therefore, since $B\cap\jepPubScript{R}(\beta)$ is dense in $B$, $s$ maps $B$ to $\iota(K/(K\cap Q_{\operatorname{rk}\beta}))$. This means that the reduction $P_{K\cap Q_{\operatorname{rk}\beta}}$ initially defined on $\jepPubScript{R}(\beta)$ extends to $\jepPubScript{U}$.
\end{proof}
We deduce the following result.
\begin{proposition}
The subsheaf $\jepPubScript{V}$ is a leafwise Higgs subsheaf of the Higgs bundle $(E,\theta)$ on $\mathbb PT_X$.
\end{proposition}
\begin{proof}
By Proposition~2.29, we know that $y$ and $\mathfrak m^+$ stabilize $\jepPubBbb{V}^{\operatorname{rk}\beta}$. Therefore, on $\jepPubScript{R}(\beta)$, the two components $\beta$ and $\gamma$ of the Higgs field stabilize the subsheaf $\jepPubScript{V}$ since it is the sheaf of sections of the subbundle $V=P_{K\cap Q_{\operatorname{rk}\beta}}(\jepPubBbb{V}^{\operatorname{rk}\beta})$ of $E=P_{K\cap Q_{\operatorname{rk}\beta}}(\jepPubBbb{E})$. By continuity, this still holds on $\jepPubScript{U}$ since on this big open set $\jepPubScript{V}$ is also the sheaf of section of $V=P_{K\cap Q_{\operatorname{rk}\beta}}(\jepPubBbb{V}^{\operatorname{rk}\beta})$. Now, $\jepPubScript{V}$ is a saturated, hence normal, subsheaf of $\jepPubScript{O}(E)$ by definition. Hence the restriction map $\jepPubScript{V}(\mathbb PT_X)\to\jepPubScript{V}(\jepPubScript{U})$ is an isomorphism since $\jepPubScript{U}$ is big. Therefore $\jepPubScript{V}$ is indeed a leafwise Higgs subsheaf of $(E,\theta)$ on $\mathbb PT_X$.
\end{proof}
\subsection{Proof of the Milnor-Wood inequality}
Recall that the slope of an $H$-module was introduced in Section~2.5 and is an element in $X(H)\otimes\mathbb Q$. To relate it to the slope of the associated vector bundle, we use the following fact which follows immediately from the definitions.
\begin{fact}
Let $H$ be a complex reductive group and let $P_H\to Z$ a holomorphic $H$-principal bundle over a complex manifold $Z$. Let $\jepPubBbb{V}_i,i\in\{1,2\}$, be $H$-modules and let $V_i=P_H(\jepPubBbb{V}_i)$ be the associated vector bundles. Assume that $\mu_H(\jepPubBbb{V}_1)=\mu_H(\jepPubBbb{V}_2)$. Then
\[
\frac{\det V_1}{\operatorname{rk}V_1}=\frac{\det V_2}{\operatorname{rk}V_2},
\]
as elements of $\operatorname{Pic}(Z)\otimes_{\mathbb Z}\mathbb Q$. In other words, the line bundles $(\det V_1)^{\otimes\operatorname{rk}V_2}$ and $(\det V_2)^{\otimes\operatorname{rk}V_1}$ are isomorphic.
\end{fact}
Therefore, if one is given a reasonable way to compute slopes of vector bundles on $Z$, e.g. the usual one if $Z$ is $X$ with its Kähler form $\omega$, or the foliated one if $Z$ is $\mathbb PT_X$ with its tautological foliation $\jepPubScript{T}$ and transverse volume form $\Omega_{\jepPubScript{G}}$, then using the notation of Fact~4.5, $\mu_H(\jepPubBbb{V}_1)=\mu_H(\jepPubBbb{V}_2)$ implies that the slopes of $V_1$ and $V_2$ are equal.

This observation, together with the computation of the slopes of the $H_{\operatorname{rk}\beta}$-submodules $\jepPubBbb{V}^{\operatorname{rk}\beta}_k$ of $\jepPubBbb{E}$ and the construction of the leafwise Higgs subsheaf $\jepPubScript{V}$ of $(E,\theta)$, gives the Milnor Wood inequality.
\begin{theorem}
We have the inequalities
\[
\deg_{\jepPubScript{T}}(E_0)+\frac{\operatorname{rk}\beta}{2}\deg_{\jepPubScript{T}}(L)\leqslant\mu_{\jepPubScript{T}}(\jepPubScript{V})\leqslant0.
\]
Therefore the Milnor-Wood inequality $\deg_{\jepPubScript{T}}(E_0)\leqslant(r_M/2)\deg_{\jepPubScript{T}}(L^\star)$ holds.
\end{theorem}
\begin{proof}
Let $n_k:=\dim\jepPubBbb{V}^{\operatorname{rk}\beta}_k$. First, recall that Proposition~2.33~(2) states that the unipotent radical of $Q_{\operatorname{rk}\beta}$ acts trivially on $\jepPubBbb{V}^{\operatorname{rk}\beta}$. Hence so does the unipotent radical of $K\cap Q_{\operatorname{rk}\beta}$. Thus, in fact, $\jepPubBbb{V}^{\operatorname{rk}\beta}$ is a $(K\cap Q_{\operatorname{rk}\beta})/R_u(K\cap Q_{\operatorname{rk}\beta})$-module, and $(K\cap Q_{\operatorname{rk}\beta})/R_u(K\cap Q_{\operatorname{rk}\beta})\simeq H_{\operatorname{rk}\beta}$ is reductive.

Therefore, by Lemma~4.3, on the big open set $\jepPubScript{U}$, the vector bundles $V_k$ such that $\jepPubScript{V}_k=\jepPubScript{O}(V_k)$ are given by $V_k=P_{K\cap Q_{\operatorname{rk}\beta}}(\jepPubBbb{V}^{\operatorname{rk}\beta}_k)$. By Proposition~2.45 and Fact~4.5 applied to $\jepPubBbb{V}^{\operatorname{rk}\beta}_k$ and $\jepPubBbb{V}'_k:=\jepPubBbb{V}^{\operatorname{rk}\beta}_0\otimes(\jepPubBbb{V}^{\operatorname{rk}\beta}_1\otimes\jepPubBbb{V}^{\operatorname{rk}\beta}_0{}^\star)^{\otimes k}$, we have on $\jepPubScript{U}$
\[
(\det V_k)^{\operatorname{rk}V'_k}\simeq(\det V'_k)^{\operatorname{rk}V_k},
\]
where $V'_k=P_{K\cap Q_{\operatorname{rk}\beta}}(\jepPubBbb{V}'_k)$. Letting $\jepPubScript{V}'_k:=\jepPubScript{V}_0\otimes(\jepPubScript{V}_1\otimes\jepPubScript{V}_0{}^\star)^{\otimes k}$, we deduce that the line bundles $(\det\jepPubScript{V}_k)^{\operatorname{rk}\jepPubScript{V}'_k}$ and $(\det\jepPubScript{V}'_k)^{\operatorname{rk}\jepPubScript{V}_k}$ on $\mathbb PT_X$ are isomorphic on $\jepPubScript{U}$, hence on $\mathbb PT_X$, because $\jepPubScript{U}$ is a big open set of $\mathbb PT_X$. Therefore
\[
\mu_{\jepPubScript{T}}(\jepPubScript{V}_k)=\mu_{\jepPubScript{T}}(\jepPubScript{V}'_k)=\deg_{\jepPubScript{T}}(\jepPubScript{V}_0)+k\,\mu_{\jepPubScript{T}}(\jepPubScript{V}^\star_0\otimes\jepPubScript{V}_1).
\]
Since $\beta^{\operatorname{rk}\beta}:\jepPubScript{V}_0\otimes L^{\operatorname{rk}\beta}\to\jepPubScript{V}_{\operatorname{rk}\beta}$ is not zero and $n_0=n_{\operatorname{rk}\beta}=1$, we also have $\mu_{\jepPubScript{T}}(\jepPubScript{V}_{\operatorname{rk}\beta})\geqslant\mu_{\jepPubScript{T}}(\jepPubScript{V}_0)+\operatorname{rk}\beta\,\mu_{\jepPubScript{T}}(L)$ so that $\mu_{\jepPubScript{T}}(\jepPubScript{V}^\star_0\otimes\jepPubScript{V}_1)\geqslant\deg_{\jepPubScript{T}}(L)$.

Finally, remembering that $n_k=n_{\operatorname{rk}\beta-k}$ by Proposition~2.35 and that $\jepPubScript{V}_0=E_0$, we get
\begin{align*}
2\deg_{\jepPubScript{T}}(\jepPubScript{V})&=\sum_{k=0}^{\operatorname{rk}\beta}\deg_{\jepPubScript{T}}(\jepPubScript{V}_k)+\sum_{k=0}^{\operatorname{rk}\beta}\deg_{\jepPubScript{T}}(\jepPubScript{V}_{\operatorname{rk}\beta-k})\\
&=\sum_{k=0}^{\operatorname{rk}\beta}(n_k\mu_{\jepPubScript{T}}(\jepPubScript{V}_k)+n_{\operatorname{rk}\beta-k}\mu_{\jepPubScript{T}}(\jepPubScript{V}_{\operatorname{rk}\beta-k}))\\
&\geqslant\sum_{k=0}^{\operatorname{rk}\beta}n_k(\deg_{\jepPubScript{T}}(\jepPubScript{V}_0)+k\deg_{\jepPubScript{T}}(L)+\deg_{\jepPubScript{T}}(\jepPubScript{V}_0)+(\operatorname{rk}\beta-k)\deg_{\jepPubScript{T}}(L))\\
&=(\dim\jepPubBbb{V}^{\operatorname{rk}\beta})(2\deg_{\jepPubScript{T}}(E_0)+\operatorname{rk}\beta\deg_{\jepPubScript{T}}(L)).
\end{align*}
The inequality $\mu_{\jepPubScript{T}}(\jepPubScript{V})\leqslant0$ follows from Propositions~4.4 and~3.2, and the conclusion from the fact that $\operatorname{rk}\beta\leqslant r_M$.
\end{proof}
In case of equality in Theorem~4.6, we have
\begin{proposition}
Assume that $\deg_{\jepPubScript{T}}(E_0)+\frac{\operatorname{rk}\beta}{2}\deg_{\jepPubScript{T}}(L)=0$. Then
\begin{enumerate}[label=(\arabic*)]
\item on the regular locus $\jepPubScript{R}(\beta)=\mathbb PT_X\backslash\jepPubScript{S}(\beta)$ of $\beta$, the orthogonal complement $V^\perp=(\jepPubOplus V_i)^\perp$ of the subbundle $V=\jepPubOplus V_i$ of $E$ w.r.t. the harmonic metric is stable under the Higgs field $\theta:E\otimes L\to E$;
\item the regular locus $\jepPubScript{R}(\overline\beta)\subset X$ of $\overline\beta$ is (open and) dense in $X$.
\end{enumerate}
\end{proposition}
\begin{proof}
The first point follows from the discussion after the definition of the subsheaves $\jepPubScript{V}_k$ and the polystability property~(2) in Proposition~3.2, since our hypothesis implies that $\deg_{\jepPubScript{T}}\jepPubScript{V}=0$ by Theorem~4.6.

Proposition~3.2~(2) also implies that the singular locus $\jepPubScript{S}(\beta)$ of $\beta$, which is a closed proper subset of $\mathbb PT_X$, is saturated under the tautological foliation $\jepPubScript{T}$, see the proof of \cite[Lem.~4.5]{jep93KM17}. Since the leaves of the tautological foliation on $\mathbb PT_X$ are projections of $\operatorname{SU}(1,1)$-homogeneous subsets of $\Gamma\backslash\operatorname{SU}(1,n)$, M.~Ratner's results on unipotent flows and the fact that $\mathbb H^n_{\mathbb C}=\operatorname{SU}(1,n)/\operatorname U(n)$ is a rank~$1$ symmetric space, then imply by \cite[Prop.~3.6]{jep93KM17} that the singular locus $\jepPubScript{S}(\overline\beta):=\pi(\jepPubScript{S}(\beta))$ of $\overline\beta$ is a proper analytic subset of $X=\Gamma\backslash\mathbb H^n_{\mathbb C}$, hence the second point of the proposition.
\end{proof}

\section{Maximal representations}
Maximal representations $\rho:\Gamma\to G_{\mathbb R}$, where $\Gamma$ is a uniform lattice in $\operatorname{SU}(1,n)$ with $n\geqslant2$ and $G_{\mathbb R}$ is a classical Hermitian group, were classified in \cite{jep93KM17}. Therefore we focus here on exceptional targets, namely $G_{\mathbb R}$ is either $\operatorname E_{6(-14)}$, which is not of tube type, or $\operatorname E_{7(-25)}$, which is.

In Section~5.1 we exclude the possibility of maximal representations in $\operatorname E_{7(-25)}$. In fact, our uniform approach allows to easily prove that maximal representations in tube type target groups $G_{\mathbb R}$ do not exist (Proposition~C). The case of $\operatorname E_{6(-14)}$ is treated in Section~5.2, Theorem~A and Corollary~B are established in Section~5.3.

\subsection{Tube type targets}
We prove Proposition~C, namely that whenever $G_{\mathbb R}$ has tube type and $n\geqslant2$, the Toledo invariant of a representation from $\Gamma$ to $G_{\mathbb R}$ satisfies an inequality stronger than the Milnor-Wood inequality, preventing a representation in such a group to be maximal.
\begin{proposition}
Let $\Gamma$ be a uniform lattice in $\mathrm{SU}(1,n)$, and let $X=\Gamma\backslash\mathbb H^n_{\mathbb C}$. Assume that the simple real algebraic Hermitian Lie group $G_{\mathbb R}$ has tube type and let $r_M$ be the real rank of $G_{\mathbb R}$. Let $\rho$ be a representation $\Gamma\to G_{\mathbb R}$. Then
\[
|\tau(\rho)|\leqslant\max\Big\{r_M-1,\frac{r_M}{2}\cdot\frac{n+1}{n}\Big\}\operatorname{vol}(X).
\]
\end{proposition}
\begin{remark}
Observe that $\max\big\{r_M-1,\frac{r_M}{2}\cdot\frac{n+1}{n}\big\}=r_M-1$ unless $r_M\leqslant2$ or $r_M=3$ and $n=2$.
\end{remark}
\begin{proof}
We may assume that $\Gamma$ is torsion free, that $\tau(\rho)>0$, that $\rho$ is reductive (see the beginning of Section~4), and that $G$ is simply connected (see Remarks~5.14 and~5.15). Then, the constructions of Sections~2,~3 and~4 are valid and the inequality of the proposition is equivalent to the inequality
\[
\deg_{\jepPubScript{T}}(E_0)\leqslant\max\Big(\frac{r_M-1}{2},\frac{r_M}{2}\cdot\frac{n+1}{2n}\Big)\deg_{\jepPubScript{T}}(L^\star).
\]
We use freely the notation of Section~4. If the generic rank of $\beta$ on the projectivized tangent bundle $\mathbb PT_X$ of $X$ is $\leqslant r_M-1$ then we are done by Theorem~4.6. Therefore we may assume that the generic rank of $\beta$ on $\mathbb PT_X$ is $r_M$.

We come back to the Higgs bundle $(\overline E,\overline\theta)$ on $X$ and we consider $\overline\beta:\overline E\otimes T_X\to\overline E$. The fact that $\operatorname{rk}\beta=r_M$ implies that $\overline\beta^{r_M}$, seen as a morphism from $\overline E_0\otimes\overline E_{r_M}^{\star}$ to the $r_M$-th symmetric power $S^{r_M}\Omega_X^1$ of $\Omega_X^1$, is not zero. Since $\Omega_X^1$ is a semistable bundle over $X$ ($X$ is Kähler-Einstein), so is $S^{r_M}\Omega_X^1$. On the other hand, $\overline E_0\otimes\overline E_{r_M}^{\star}$ is also semistable because it is a line bundle by Section~2.4.4. Therefore $\mu_\omega(\overline E_0\otimes\overline E_{r_M}^{\star})\leqslant\mu_\omega(S^{r_M}\Omega_X^1)$ (recall that $\omega$ is our Kähler form on $X$), so that $\deg_\omega\overline E_0-\deg_\omega\overline E_{r_M}\leqslant r_M\mu_\omega(\Omega_X^1)$. Now, as explained in Section~2.4.4, the $K$-modules $\jepPubBbb{E}_{r_M}$ and $\jepPubBbb{E}_0^{\star}$ are isomorphic, so that $\deg_\omega\overline E_{r_M}=-\deg_\omega\overline E_0$ by Fact~4.5. We get the result, since by Equations~(4.2) and~(4.3) in Section~4.2 and the isomorphism $K_X\simeq\det\Omega_X^1$, we have $\deg_\omega(\Omega_X^1)=\frac{n+1}{2}\deg_{\jepPubScript{T}}(L^\star)$.
\end{proof}

{\footnotesize
\begin{thebibliography}{AMRT10}
\bibitem[AMRT10]{jep93AMRT10} A. Ash, D. Mumford, M. Rapoport \& Y.-S. Tai -- \emph{Smooth compactifications of locally symmetric varieties}, 2nd ed., Cambridge University Press, Cambridge, 2010.
\bibitem[BGPR17]{jep93BGPR17} O. Biquard, O. García-Prada \& R. Rubio -- ``Higgs bundles, the Toledo invariant and the Cayley correspondence'', \emph{J. Topology} \textbf{10} (2017), no.~3, p.~795--826.
\bibitem[Bor69]{jep93Bor69} A. Borel -- \emph{Linear algebraic groups}, W.~A. Benjamin, Inc., New York-Amsterdam, 1969.
\bibitem[Bou68]{jep93Bou68} N. Bourbaki -- \emph{Éléments de mathématique. Fasc.~XXXIV. Groupes et algèbres de Lie. Chapitre IV: Groupes de Coxeter et systèmes de Tits. Chapitre V: Groupes engendrés par des réflexions. Chapitre VI: systèmes de racines}, Actualités Scientifiques et Industrielles, vol.~1337, Hermann, Paris, 1968.
\bibitem[BGPG03]{jep93BGPG03} S.~B. Bradlow, O. Garcia-Prada \& P.~B. Gothen -- ``Surface group representations and $\operatorname U(p,q)$-Higgs bundles'', \emph{J. Differential Geom.} \textbf{64} (2003), p.~111--170.
\bibitem[BGPG06]{jep93BGPG06} \rule{1.2em}{.1pt}, ``Maximal surface group representations in isometry groups of classical Hermitian symmetric spaces'', \emph{Geom. Dedicata} \textbf{122} (2006), p.~185--213.
\bibitem[BH99]{jep93BH99} M.~R. Bridson \& A. Haefliger -- \emph{Metric spaces of non-positive curvature}, Grundlehren Math. Wiss., vol.~319, Springer-Verlag, Berlin, 1999.
\bibitem[BM15]{jep93BM15} A.~S. Buch \& L.~C. Mihalcea -- ``Curve neighborhoods of Schubert varieties'', \emph{J. Differential Geom.} \textbf{99} (2015), no.~2, p.~255--283.
\bibitem[BI07]{jep93BI07} M. Burger \& A. Iozzi -- ``Bounded differential forms, generalized Milnor-Wood inequality and an application to deformation rigidity'', \emph{Geom. Dedicata} \textbf{125} (2007), p.~1--23.
\bibitem[BIW09]{jep93BIW09} M. Burger, A. Iozzi \& A. Wienhard -- ``Tight homomorphisms and Hermitian symmetric spaces'', \emph{Geom. Funct. Anal.} \textbf{19} (2009), no.~3, p.~678--721.
\bibitem[BIW10]{jep93BIW10} \rule{1.2em}{.1pt}, ``Surface group representations with maximal Toledo invariant'', \emph{Ann. of Math.} (2) \textbf{172} (2010), p.~517--566.
\bibitem[Che97]{jep93Che97} C. Chevalley -- \emph{The algebraic theory of spinors and Clifford algebras}, Collected works, vol.~2, Springer-Verlag, Berlin, 1997.
\bibitem[Cor88]{jep93Cor88} K. Corlette -- ``Flat $G$-bundles with canonical metrics'', \emph{J. Differential Geom.} \textbf{28} (1988), p.~361--382.
\bibitem[Del80]{jep93Del80} P. Deligne -- ``La conjecture de Weil. II'', \emph{Publ. Math. Inst. Hautes Études Sci.} \textbf{52} (1980), p.~137--252.
\bibitem[FH91]{jep93FH91} W. Fulton \& J. Harris -- \emph{Representation theory. A first course}, Graduate Texts in Math., vol.~129, Springer-Verlag, New York, 1991.
\bibitem[Gro94]{jep93Gro94} B.~H. Gross -- ``A remark on tube domains'', \emph{Math. Res. Lett.} \textbf{1} (1994), no.~1, p.~1--9.
\bibitem[GW12]{jep93GW12} O. Guichard \& A. Wienhard -- ``Anosov representations: domains of discontinuity and applications'', \emph{Invent. Math.} \textbf{190} (2012), p.~357--438.
\bibitem[Ham13]{jep93Ham13} O. Hamlet -- ``Tight holomorphic maps, a classification'', \emph{J. Lie Theory} \textbf{23} (2013), no.~3, p.~639--654.
\bibitem[HC56]{jep93HC56} Harish-Chandra -- ``Representations of semisimple Lie groups. VI. Integrable and square-integrable representations'', \emph{Amer. J. Math.} \textbf{78} (1956), p.~564--628.
\bibitem[Hel01]{jep93Hel01} S. Helgason -- \emph{Differential geometry, Lie groups, and symmetric spaces}, Graduate Studies in Math., vol.~34, American Mathematical Society, Providence, RI, 2001, Corrected reprint of the 1978 original.
\bibitem[Her91]{jep93Her91} L. Hernández -- ``Maximal representations of surface groups in bounded symmetric domains'', \emph{Trans. Amer. Math. Soc.} \textbf{324} (1991), p.~405--420.
\bibitem[Hit87]{jep93Hit87} N. Hitchin -- ``The self-duality equations on a Riemann surface'', \emph{Proc. London Math. Soc.} \textbf{55} (1987), p.~59--126.
\bibitem[Hit92]{jep93Hit92} \rule{1.2em}{.1pt}, ``Lie groups and Teichmüller space'', \emph{Topology} \textbf{31} (1992), p.~449--473.
\bibitem[Hum75]{jep93Hum75} J.~E. Humphreys -- \emph{Linear algebraic groups}, Graduate Texts in Math., vol.~21, Springer, New York, 1975.
\bibitem[Igu70]{jep93Igu70} J.-i. Igusa -- ``A classification of spinors up to dimension twelve'', \emph{Amer. J. Math.} \textbf{92} (1970), p.~997--1028.
\bibitem[Iha67]{jep93Iha67} S.-i. Ihara -- ``Holomorphic imbeddings of symmetric domains'', \emph{J. Math. Soc. Japan} \textbf{19} (1967), p.~261--302.
\bibitem[Kna02]{jep93Kna02} A.~W. Knapp -- \emph{Lie groups beyond an introduction}, 2nd ed., Progress in Math., vol.~140, Birkhäuser Boston, Inc., Boston, MA, 2002.
\bibitem[Kos12]{jep93Kos12} B. Kostant -- ``The cascade of orthogonal roots and the coadjoint structure of the nilradical of a Borel subgroup of a semisimple Lie group'', \emph{Moscow Math. J.} \textbf{12} (2012), no.~3, p.~605--620.
\bibitem[KM08]{jep93KM08} V. Koziarz \& J. Maubon -- ``Representations of complex hyperbolic lattices into rank 2 classical Lie groups of Hermitian type'', \emph{Geom. Dedicata} \textbf{137} (2008), p.~85--111.
\bibitem[KM10]{jep93KM10} \rule{1.2em}{.1pt}, ``The Toledo invariant on smooth varieties of general type'', \emph{J. reine angew. Math.} \textbf{649} (2010), p.~207--230.
\bibitem[KM17]{jep93KM17} \rule{1.2em}{.1pt}, ``Maximal representations of uniform complex hyperbolic lattices'', \emph{Ann. of Math.} (2) \textbf{185} (2017), p.~493--540.
\bibitem[Man06]{jep93Man06} L. Manivel -- ``Configurations of lines and models of Lie algebras'', \emph{J. Algebra} \textbf{304} (2006), no.~1, p.~457--486.
\bibitem[MX02]{jep93MX02} E. Markman \& E.~Z. Xia -- ``The moduli of flat $\operatorname{PU}(p,p)$-structures with large Toledo invariants'', \emph{Math. Z.} \textbf{240} (2002), p.~95--109.
\bibitem[McG02]{jep93McG02} W.~M. McGovern -- ``The adjoint representation and the adjoint action'', in \emph{Algebraic quotients. Torus actions and cohomology. The adjoint representation and the adjoint action}, Encyclopaedia Math. Sci., vol.~131, Springer, Berlin, 2002, p.~159--238.
\bibitem[Mok89]{jep93Mok89} N. Mok -- \emph{Metric rigidity theorems on Hermitian locally symmetric manifolds}, Series in Pure Mathematics, vol.~6, World Scientific, Teaneck, NJ, 1989.
\bibitem[Mur59]{jep93Mur59} S. Murakami -- ``Sur certains espaces fibrés principaux différentiables et holomorphes'', \emph{Nagoya Math. J.} \textbf{15} (1959), p.~171--199.
\bibitem[NT76]{jep93NT76} H. Nakagawa \& R. Takagi -- ``On locally symmetric Kaehler submanifolds in a complex projective space'', \emph{J. Math. Soc. Japan} \textbf{28} (1976), no.~4, p.~638--667.
\bibitem[PS69]{jep93PS69} I. Piatetski-Shapiro -- \emph{Automorphic functions and the geometry of classical domains}, Mathematics and its applications, vol.~8, Gordon and Breach Science Publishers, New York-London-Paris, 1969.
\bibitem[Rat06]{jep93Rat06} J. Ratcliffe -- \emph{Foundations of hyperbolic manifolds}, Springer, New York, 2006.
\bibitem[RRS92]{jep93RRS92} R. Richardson, G. Röhrle \& R. Steinberg -- ``Parabolic subgroups with abelian unipotent radical'', \emph{Invent. Math.} \textbf{110} (1992), no.~3, p.~649--671.
\bibitem[Roy80]{jep93Roy80} H.~L. Royden -- ``The Ahlfors-Schwarz lemma in several complex variables'', \emph{Comment. Math. Helv.} \textbf{55} (1980), no.~4, p.~547--558.
\bibitem[Sam78]{jep93Sam78} J.~H. Sampson -- ``Some properties and applications of harmonic mappings'', \emph{Ann. Sci. École Norm. Sup.} (4) \textbf{11} (1978), p.~211--228.
\bibitem[Sat80]{jep93Sat80} I. Satake -- \emph{Algebraic structures of symmetric domains}, Kanô Memorial Lectures, vol.~4, Princeton University Press, Princeton, NJ, 1980.
\bibitem[Sel60]{jep93Sel60} A. Selberg -- ``On discontinuous groups in higher-dimensional symmetric spaces'', in \emph{Contributions to function theory (internat. Colloq. Function Theory, Bombay, 1960)}, Tata Institute of Fundamental Research, Bombay, 1960, p.~147--164.
\bibitem[SZ10]{jep93SZ10} M. Sheng \& K. Zuo -- ``Polarized variation of Hodge structures of Calabi-Yau type and characteristic subvarieties over bounded symmetric domains'', \emph{Math. Ann.} \textbf{348} (2010), no.~1, p.~211--236.
\bibitem[Sim88]{jep93Sim88} C. Simpson -- ``Constructing variations of Hodge structure using Yang-Mills theory and applications to uniformization'', \emph{J. Amer. Math. Soc.} \textbf{1} (1988), p.~867--918.
\bibitem[Sim92]{jep93Sim92} \rule{1.2em}{.1pt}, ``Higgs bundles and local systems'', \emph{Publ. Math. Inst. Hautes Études Sci.} \textbf{75} (1992), p.~5--95.
\bibitem[Siu80]{jep93Siu80} Y.-T. Siu -- ``The complex-analyticity of harmonic maps and the strong rigidity of compact Kähler manifolds'', \emph{Ann. of Math.} (2) \textbf{112} (1980), p.~73--111.
\bibitem[SV00]{jep93SV00} T.~A. Springer \& F.~D. Veldkamp -- \emph{Octonions, Jordan algebras and exceptional groups}, Springer Monographs in Mathematics, Springer-Verlag, Berlin, 2000.
\bibitem[SZ85]{jep93SZ85} J. Steenbrink \& S. Zucker -- ``Variation of mixed Hodge structure. I'', \emph{Invent. Math.} \textbf{80} (1985), p.~489--542.
\bibitem[Wol72]{jep93Wol72} J.~A. Wolf -- ``Fine structure of Hermitian symmetric spaces'', in \emph{Symmetric spaces (Short Courses, Washington Univ., St. Louis, Mo., 1969--1970)}, Pure and App. Math., vol.~8, Dekker, New York, 1972, p.~271--357.
\bibitem[Xia00]{jep93Xia00} E.~Z. Xia -- ``The moduli of flat $\operatorname{PU}(2,1)$ structures on Riemann surfaces'', \emph{Pacific J. Math.} \textbf{195} (2000), p.~231--256.
\end{thebibliography}
}

\end{document}
